%% ============================================================================
%% Supplementary Information
%% Stratospheric Sudden Warmings Are Associated with Suppressed
%% Natural Dry Slab Avalanche Activity in the European Alps
%%
%% Target: Nature Geoscience (Article format)
%% Format: Nature Portfolio Design Template
%% ============================================================================
\documentclass[11pt,letterpaper]{article}

% Packages
\usepackage[utf8]{inputenc}
\usepackage[T1]{fontenc}
\usepackage{amsmath,amssymb}
\usepackage{graphicx}
\usepackage{booktabs}
\usepackage{hyperref}
\usepackage[numbers,sort&compress]{natbib}
\usepackage{xcolor}
\usepackage{geometry}
\usepackage{adjustbox}    % For table width adjustment
\usepackage{array}        % Enhanced column types
\usepackage{multirow}     % Multi-row cells
\usepackage{caption}
\captionsetup{font=small,labelfont=bf}
\usepackage{setspace}
\onehalfspacing

% Nature formatting
\geometry{margin=1in}
\sloppy  % Allow more flexible line breaking to prevent overfull hboxes

% URL formatting - allow breaks at more characters
\makeatletter
\g@addto@macro{\UrlBreaks}{\UrlOrds}
\makeatother

\title{Supplementary Information:\\
Stratospheric Sudden Warmings Are Associated with\\
Suppressed Natural Avalanche Activity in the European Alps:\\
Multi-Country Evidence and Process-Model Diagnostics}

\author{Abduxoliq Ashuraliyev}
\date{}

\begin{document}

\maketitle

\tableofcontents

\section*{Supplementary Text: Extended Abstract and Plain Language Summary}
\begin{abstract}
Stratospheric sudden warming (SSW) events produce persistent surface weather anomalies, yet their consequences for threshold-governed geophysical hazards remain poorly constrained.
Here we report a strong association between SSW episodes and reduced natural dry slab avalanche activity in the European Alps, with directional consistency across independent measurement systems in multiple countries.
In Switzerland (21~winters, 16~SSW events), 14 of 16 events show reduced counts (geometric mean rate ratio $\mathrm{RR} = 0.32$; 95\% CI $[0.26, 0.53]$ winter-blocked bootstrap; $d = -1.06$; permutation $P < 0.001$).
Norwegian avalanche danger levels decrease concordantly ($d = -0.67$), a reconstructed French north-Alps BRA archive also shows lower matched-window danger ($\Delta = -0.15$; sign test $P = 0.011$), and all four Utah SSW events show decreased counts ($\mathrm{RR} = 0.34$), whereas a Canadian flexible-forecast panel shows higher area-weighted danger, demonstrating regional dependence.
A specification curve of 180 analytical variants confirms robustness: 100\% show $\mathrm{RR} < 1$.
SNOWPACK process-model diagnostics show that loading and persistent-weak-layer susceptibility remain elevated during SSW windows, while the proximate natural-release triggers are selectively suppressed (warming $-57\%$, rain-on-snow $-29\%$). Human-triggered avalanches increase ($\mathrm{RR} = 1.81\times$ relative to natural), consistent with wind-slab formation enhancing human-triggerable instability while cold-dry conditions suppress spontaneous release.
The French BRA suppression provides cross-model validation: the S2M reanalysis chain (SAFRAN-Crocus-MEPRA) uses ARPEGE forcing entirely independent of the ERA5/SNOWPACK chain used for Swiss diagnostics.
The temporal profile---suppression emerging before SSW onset---supports a common-cause interpretation in which planetary-wave amplification simultaneously produces the SSW and reorganises Alpine surface weather.
\end{abstract}

\noindent\textbf{Plain Language Summary.}
Snow avalanches in the Alps kill hundreds of people each year. We show that a dramatic atmospheric event called a ``stratospheric sudden warming''---where temperatures 30--50~km above Earth surge by 30--40~°C in just a few days---is consistently associated with a sharp drop in natural avalanche activity at the surface. Across 16 such events in Switzerland over 21~winters, natural dry slab avalanches decreased by 68\% on average. The same pattern appears in Norway, in a reconstructed French north-Alps bulletin archive, and in Utah using entirely different measurement systems. A recent Canadian forecast archive does not show the same decrease, which means the effect is not universal and instead depends on how each SSW reorganises regional storms. Using output from SNOWPACK, a physics-based snowpack model run at 130~Swiss stations, we find that loading and weak-layer susceptibility remain elevated during these events, but the natural triggers that would normally release avalanches (surface warming, rain-on-snow, solar radiation) are simultaneously suppressed by the cold regime. This creates a ``loaded gun'' situation: a primed snowpack without natural release, which may explain why human-triggered avalanche rates increase during these windows. Because stratospheric sudden warmings can sometimes be predicted 1--2~weeks in advance, our findings may help extend avalanche hazard forecasting beyond the current 5--7~day weather-based window.


\subsection*{Out-of-sample geographic gradient: ALBINA and expanded Norwegian validation}

The EAWS geographic gradient (Western Alpine decrease, Northern Alpine increase) generates testable predictions for independent datasets. We tested these predictions using two independent databases.

Detailed results from ALBINA (Austria, South Tyrol, Trentino), an expanded Norwegian analysis, and the Austrian Tyrolean incident database consistently support the ``loaded gun'' mechanism and the geographic gradient (full analysis in Supplementary Information).

\subsection*{Dry slab specificity}

The signal is specific to natural-trigger dry slab avalanches ($\mathrm{RR} = 0.32$, $P = 0.002$). Wet natural avalanches show a nominally significant \emph{increase} ($\mathrm{RR} = 1.62$, $P = 0.038$ before multiple-testing correction; non-significant after BH-FDR), which requires mechanistic interpretation rather than dismissal. The increase is concentrated in a subset of events: three SSW windows (January 2019, February 2010, February 2007) account for 78\% of the excess wet avalanche counts, and the signal is sensitive to their inclusion. We interpret this as consistent with a delayed melt-release mechanism: SSW-associated cold-dry conditions suppress meltwater triggers during the event window, but the accumulated snowpack loading produces occasional large wet slab releases when the post-SSW recovery phase delivers the first warm pulse to an anomalously deep, cold snowpack. This ``spring rebound'' interpretation is consistent with known avalanche climatology in which delayed onset of spring warming can produce more destructive wet avalanche cycles\cite{Schweizer2003}. However, the wet avalanche result is based on fewer events with interpretable signal, and we treat it as hypothesis-generating rather than as evidence of a second SSW pathway. The dry--wet divergence is itself physically informative: dry slab release depends on the balance between slab load and weak-layer shear strength modulated by surface-energy triggers, whereas wet avalanches are primarily triggered by meltwater infiltration on a different timescale.

\subsection*{Rutschblock field validation}

Independent Rutschblock stability tests ($n = 4{,}433$; 2001--2019) show a small but significant field-scale stability increase during SSW windows (mean class 1.99 vs.\ 1.90; $P = 0.003$; $d = 0.10$). The effect is far smaller than the count reduction ($d = -1.06$), so it cannot explain the 68\% decrease on its own. Together with SNOWPACK, the Rutschblock result indicates that SSW windows alter trigger availability more strongly than they alter bulk snowpack stability.

\subsection*{Observer-effort validation}

A potential confound is that SSW-associated cold weather could reduce backcountry access, decreasing observed counts. Four lines of evidence rule this out: total avalanche reports (all types) \emph{increase} during SSW windows ($\mathrm{RR} = 1.15$; $P = 0.002$); human-triggered counts across all avalanche types increase 45\% ($\mathrm{RR} = 1.45$; $P = 0.0002$; EnviDat re-analysis); all dry avalanche counts (natural $+$ human) increase ($\mathrm{RR} = 1.70$) while only the natural-trigger subset decreases; and professional danger ratings increase ($P = 0.02$). The selective decrease of natural triggers within an overall increase is inconsistent with any access-based explanation.

\subsection*{Mechanistic investigation: weather regime shift carries the signal}

We identify the physical mechanism in three steps: (i)~aggregate surface weather changes, (ii)~weather regime decomposition, and (iii)~event-level synoptic-scale prediction.

\paragraph{Aggregate surface weather changes.} ERA5 composites (Swiss Alps, 1998--2019) reveal SSW windows are significantly colder ($\Delta T_{2\mathrm{m}} = -3.3$~K, $d = -0.69$, $P < 10^{-47}$) with a critical precipitation phase shift: rainfall decreases 34\% ($d = -0.18$; $P < 10^{-10}$) while snowfall is maintained ($+5\%$, $P = 0.56$), producing a 27\% increase in snow fraction ($d = +0.30$; $P < 10^{-7}$). Rain-on-snow events---a primary natural trigger mechanism---decrease across 11 of 12 threshold definitions tested (temperature $> 0$--$2$°C $\times$ precipitation $> 0.5$--$2$~mm; sign consistency $P = 0.003$, binomial test; RR range $0.47$--$0.97$), providing robust observational support for the SNOWPACK surface-energy trigger suppression pathway. This explains the SNOWPACK paradox: solid loading increases despite reduced total precipitation, while liquid-water triggers are selectively removed. However, event-level surface weather explains only $\sim$1\% of the event-to-event avalanche response ($R^2 = 0.011$), indicating that synoptic-scale circulation reorganisation, not local surface variables, drives event-level variability.

\paragraph{Weather regime decomposition.} Classifying winter days into four regimes by median temperature and precipitation splits reveals a massive redistribution during SSW windows ($\chi^2 = 83.5$, $P = 5 \times 10^{-18}$; Extended Data Table~\ref{tab:regime}). Cold-dry days---associated with the lowest avalanche rates---nearly double in frequency (22.9\% to 40.5\%), while warm-wet days decrease from 28.3\% to 18.3\%. A multiplicative occupancy decomposition (Supplementary Table~\ref{tab:regime_decomp}) attributes $\sim$6\% of the aggregate log-RR to regime-frequency shifts and $\sim$95\% to within-regime rate changes, demonstrating that both pathways contribute but within-regime trigger suppression dominates. A linear Blinder--Oaxaca decomposition assigns $-91\%$ of the daily-level rate gap to composition (regime frequencies) and $+214\%$ to within-regime rate changes; these components exceed $\pm 100\%$ because the linear approximation is strained by highly nonlinear count data. The cleanest coupling test is the regime-conditioned incidence rate ratio: $\mathrm{IRR} = 0.89$ ($P = 0.06$; Supplementary Table~14), meaning that after controlling for weather regime and day-of-year simultaneously, only $\sim$11\% residual suppression remains attributable to SSW-specific effects beyond regime redistribution. The $\mathrm{IRR} = 0.89$ result and the 6\% occupancy-decomposition share are compatible: both indicate that regime indicators absorb most of the signal, because regime categories capture both frequency and rate information jointly when included as regressors.

\paragraph{Synoptic-scale circulation as the integrating variable.} Northern Hemisphere Z500 anomalies during SSW windows correlate significantly with event-level avalanche response ($r = 0.56$, $P = 0.024$, $R^2 = 0.31$), and SLP shows a similar relationship ($r = 0.52$, $P = 0.037$, $R^2 = 0.27$). A bivariate model combining Z500 and SLP increases the explained event-level variance to $R^2 = 0.41$ ($R^2_{\mathrm{adj}} = 0.32$), while adding ERA5 surface variables contributes negligibly ($\Delta R^2 = 0.001$), confirming that synoptic-scale circulation fully subsumes surface weather information. In contrast, stratospheric variables show near-zero correlation ($R^2 < 0.03$). The propagation is temporally resolved: geopotential height anomalies cascade from 10~hPa (lag $+3$~d) through 50~hPa ($+20$~d) to 500~hPa ($+30$~d; Extended Data Table~\ref{tab:propagation}).

\paragraph{Phase-level coherence.} All four SSW lifecycle phases show significant suppression ($P < 0.04$), with geometric mean RR remarkably consistent across phases ($0.32$--$0.34$). The persistence of suppression into the recovery phase ($+31$ to $+45$~d) even as Z500 anomalies weaken indicates a snowpack-integrated response lag consistent with the ``loaded gun'' configuration.

The overall mechanism is best interpreted as a Z500-mediated causal pathway: enhanced planetary wave forcing preconditions the stratospheric vortex disruption and simultaneously reorganises the tropospheric circulation; SSW onset marks the culmination of this wave-forced episode and functions as the detectable lead-time marker. The null partial correlation of stratospheric variables after conditioning on Z500 ($\rho_{\mathrm{partial}} = -0.28$, $P = 0.29$) is the expected signature of \emph{full mediation through Z500}---not evidence against a top-down pathway. Once the mediator is included in the model, the original predictor's residual direct effect approaches zero by construction (the mediation null), consistent with published GCM initial-condition experiments showing SSW drives tropospheric blocking through the Z500 channel\cite{Hitchcock2013,Baldwin2021}. SSW onset therefore acts as a detectable precursor marker of this broader wave-forced episode and provides categorical warning lead time over tropospheric blocking forecasts---consistent with the pre-onset suppression (13/16 events; $P = 0.011$). The timing rules out a purely post-onset pathway, but does not exclude additional post-onset coupling or cold-pool persistence. The regime-frequency shift nevertheless explains the dominant effect (Supplementary Information contains the full propagation cascade analysis and power analysis).

\paragraph{Annular mode independence.} The SSW--avalanche association is not transmitted through large-scale annular modes ($r = -0.07$, $P = 0.79$ for AO; $r = 0.03$, $P = 0.92$ for NAO), suggesting a pathway distinct from the canonical SSW~$\to$~negative NAO chain\cite{Baldwin2001,Kidston2015}.

\paragraph{Hemispheric-scale verification.} Independent satellite and ground-based datasets confirm the SSW--surface coupling chain operates at hemispheric scale. Aura/MLS data show the expected chemical fingerprint: O$_3$ $+8.5\%$ at 10~hPa ($d = +0.81$), N$_2$O $+50\%$ at 1~hPa ($d = +0.69$), HNO$_3$ $-11.7\%$ ($d = -0.57$)---confirming genuine vortex disruptions. SNOTEL stations across the western US ($n = 823$) show higher SWE ($+61$~mm, $d = +0.43$) and Northern Hemisphere snow cover increases ($d = +1.13$, the largest effect tested), establishing that the surface response is not Alpine-specific (Extended Data Table~\ref{tab:hemispheric}). The SWE loading is anomalous both relative to DOY-matched controls ($\Delta\mathrm{SWE} = +47.5$~mm, $P = 2.5 \times 10^{-7}$) and relative to climatology ($+61$~mm, $d = +0.43$), confirming that SSW-forced precipitation loading exceeds seasonal expectations.

\paragraph{Reconciling the three SNOTEL temperature comparisons.} Three temperature analyses in this study yield apparently different results; all three are correct and test different questions. (1)~\emph{DOY-matched event-level}: $\Delta T = -0.08$~°C, $P = 0.88$ (main text). This comparison asks whether SSW windows are anomalously cold \emph{relative to the same calendar date in non-SSW years}. The non-significant result means that SSW windows, at the scale of the 16-event sample, are not distinguishably colder than seasonally expected for those weeks of winter. (2)~\emph{Pooled station-day climatological}: $\Delta T = -1.56$~°C, $d = -0.42$, $P < 10^{-15}$ (Extended Data Table~\ref{tab:hemispheric}). This comparison asks whether station-days during SSW windows are \emph{physically cold} enough to promote kinetic-growth metamorphism. The highly significant result confirms they are, because mid-winter alpine temperatures are cold in absolute terms regardless of SSW. The two tests are not contradictory: kinetic metamorphism requires temperatures below $-3$~°C (a physical threshold), not anomalously cold conditions relative to seasonal expectations. (3)~\emph{ERA5 gridbox mean event-level}: $P = 0.32$ (main text, 2~m temperature). This comparison tests an aggregated alpine gridbox average rather than individual SNOTEL station-days, substantially reducing power through spatial averaging. All three results are consistent with the loaded-gun mechanism: physical temperatures are cold enough for weak-layer growth (result 2), without necessarily being anomalously cold for the time of year (result 1) or distinguishable at coarse ERA5 resolution (result 3).

\subsection*{Robustness and sensitivity}

The result is highly robust to multiple analytical choices, including DOY matching bandwidth, SSW window width, and alternative event catalogs (e.g., displacement-only vs. split-vortex events). Comprehensive tests including a specification curve of 180 analytical variants confirm the durability of the avalanche suppression. We present the full sensitivity analyses, including Bayesian inference and warm-dry regime exclusion, in the Supplementary Information.

\section*{Supplementary Text: Extended narrative from the main Article}

The following subsections preserve longer narrative text from the Article, providing detailed methodological context. These sections expand on the main text and should be read alongside the Article for complete methodological transparency.

\subsection*{Extended Introduction}

Stratospheric sudden warming events are among the most dramatic circulation phenomena in the middle atmosphere. During an SSW, the wintertime stratospheric polar vortex weakens or reverses, with zonal-mean temperatures at 10~hPa increasing by 20--40~K within days\cite{Matsuno1971,Charlton2007}. The downward coupling of SSW signals to the troposphere is well established: SSW events precede persistent surface weather anomalies including cold-air outbreaks over Eurasia and North America, shifts in the North Atlantic Oscillation (NAO) toward its negative phase, and altered precipitation patterns across mid-latitudes\cite{Baldwin2001,Thompson2002,Kidston2015}. These tropospheric responses typically emerge 1--4~weeks after SSW onset and can persist for 30--60~days\cite{Sigmond2013}.

Crucially, SSW events are themselves \emph{caused} by enhanced tropospheric planetary wave activity that propagates upward and breaks in the stratosphere\cite{Matsuno1971,Polvani2004}. Any surface anomaly observed around SSW onset could therefore reflect either top-down stratospheric forcing (SSW~$\to$~troposphere~$\to$~surface) or a common-cause pathway in which the same planetary wave forcing simultaneously triggers the SSW and reorganises surface weather. Distinguishing these interpretations requires careful attention to the temporal profile, dose--response relationships, and coupling pathways.

Despite the established influence of SSW events on surface weather, their consequences for geophysical hazards have rarely been systematically quantified---a notable gap given that threshold-dependent processes (avalanches, landslides, floods) can amplify modest atmospheric anomalies into disproportionate hazard changes. Snow avalanches---responsible for approximately 250~fatalities per year worldwide\cite{Techel2016}---are fundamentally sensitive to the meteorological variables that SSW events perturb: temperature, precipitation, wind speed, and their temporal gradients\cite{Schweizer2003,Schweizer2008}. Dry slab avalanches, in particular, are triggered when the shear stress on a buried weak layer exceeds its strength---a threshold process governed by snow loading rates, temperature metamorphism, and radiative balance\cite{McClung2006,Reuter2018}.

Here we present an event-level demonstration that stratospheric variability can measurably affect a threshold-governed surface hazard: a large reduction in natural dry slab avalanche activity during SSW windows in Switzerland, with supportive external evidence from Norway, the northern French Alps and Utah, and a recent Canadian counterexample that bounds generalizability. We quantify the effect using rate ratios relative to day-of-year-matched controls, test candidate mechanisms through ERA5 reanalysis composites, verify snowpack changes with independent Rutschblock field tests, and probe annular-mode mediation. We report all results transparently, including null and opposite-signed findings for external pathways, and apply Benjamini--Hochberg false discovery rate correction across all primary hypothesis tests.

\subsection*{Extended Swiss event-level and temporal narrative}

We analysed daily counts of natural-trigger dry slab avalanches (size class $\geq$1) from the Swiss WSL/SLF database spanning 21~winters (1998/99--2018/19; 7,518~days). This variable isolates weather-driven release events by excluding human-triggered avalanches, which depend on recreational exposure patterns rather than atmospheric forcing.

For each of 16~SSW events identified from the Butler et al.\ catalog\cite{Butler2015}, we computed observed counts in a $\pm$15-day window around onset and compared them to expected counts derived from day-of-year-matched climatology ($\pm$3~days, non-SSW winters only). The rate ratio $\mathrm{RR} = \mathrm{observed} / \mathrm{expected}$ measures the multiplicative change in avalanche activity.

Of 16~SSW events, 14 show $\mathrm{RR} < 1$ (sign test $P = 0.002$). The geometric mean $\mathrm{RR} = 0.32$ (95\% CI: $[0.26, 0.53]$, winter-blocked bootstrap; Supplementary Table~\ref{tab:winter_boot}), corresponding to a 68\% reduction in natural dry slab counts. The effect is highly significant across multiple test frameworks: Wilcoxon signed-rank on $\log(\mathrm{RR})$, $P < 0.001$; permutation test (10,000 iterations), $P = 0.0005$. The one-sample $t$-test on $\log(\mathrm{RR})$ yields $P = 0.004$, higher than the non-parametric tests because two positive outliers inflate variance (kurtosis~$= 4.2$). Cohen's $d = -1.06$ indicates a large effect size.

The two exceptions are February 2018 ($\mathrm{RR} = 1.05$; the split-vortex ``Beast from the East'' that produced an extreme cold-air outbreak over western Europe, generating exceptional wind slab loading on exposed aspects---a synoptic pattern that \emph{enhanced} rather than suppressed natural triggers) and January 2019 ($\mathrm{RR} = 3.43$; this SSW coincided with the heaviest Alpine snowfall episode since 1999, which deposited extreme loading that overwhelmed trigger-suppression effects---the only event where cumulative precipitation exceeded $3\sigma$ above the DOY mean). Leave-one-out cross-validation confirms robustness: all 16 jackknife folds retain $P < 0.004$, and geometric mean RR ranges from 0.28 to 0.36 across folds. No single event drives the result. Critically, excluding the most influential observation (January 2019) \emph{strengthens} the effect: $\mathrm{RR} = 0.21$, 14/15 decrease, $P = 3 \times 10^{-6}$ (Supplementary Table~\ref{tab:outlier_sens}).

The avalanche suppression is not confined to the post-SSW period. Dividing the SSW lifecycle into four non-overlapping phases reveals significant suppression across all phases: pre-SSW ($-15$ to $-6$~d: 13/16 decrease, median $\mathrm{RR} = 0.16$, $P = 0.011$), onset ($-5$ to $+5$~d: 14/16 decrease, median $\mathrm{RR} = 0.37$, $P = 0.002$), post ($+6$ to $+15$~d: 12/16 decrease, median $\mathrm{RR} = 0.22$, $P = 0.038$), and late ($+16$ to $+30$~d: 12/16 decrease, median $\mathrm{RR} = 0.29$, $P = 0.038$).

The pre-onset suppression is a critical diagnostic for timing. SSW events are preceded by 1--3~weeks of enhanced planetary wave activity\cite{Polvani2004}, and this precursor forcing is itself part of the SSW life cycle. The avalanche anomaly emerging \emph{before} the formal stratospheric wind reversal therefore indicates that Alpine surface weather is reorganised during the precursor phase, not only after the 10~hPa threshold is crossed. This timing fingerprint is not Switzerland-specific: Swiss and Norwegian lag profiles peak at $-13$ and $-12$~days, respectively, and are positively correlated across lags $-15$ to $+15$~days (Spearman $r = 0.62$, $P = 2.2 \times 10^{-4}$).

\subsection*{Extended cross-context and geographic-gradient narrative}

The strength of multi-country evidence lies not in statistical replication of a single test but in \emph{convergent evidence from independent measurement systems testing different predictions} of the SSW--avalanche hypothesis. Switzerland provides occurrence-count primary evidence ($n = 16$; $d = -1.06$); Norway tests whether expert-assessed danger forecasts, an entirely independent measurement system, respond concordantly ($d = -0.67$); France tests whether a western-Alpine bulletin archive independently lands on the suppressive branch predicted by the EAWS gradient; Utah tests whether the association persists on a different continent and mountain range; Avalanche Canada provides a recent flexible-forecast stress test outside the core Alpine suppression zone; EAWS tests whether the spatial structure is systematic (geographic gradient $r = 0.69$); ALBINA tests out-of-sample geographic predictions; and LAWIS tests the complementary human-trigger increase.

The Norwegian Water Resources and Energy Directorate (NVE) publishes daily avalanche danger level forecasts (scale 1--5) for mountain regions across Norway. We obtained forecasts from ten inland mountain regions spanning 61--70°N (Nord-Troms, Lyngen, Tromsø, Indre Troms, Trollheimen, Romsdal, Sunnmøre, Indre Fjordane, Jotunheimen, Indre Sogn; $n = 899$~SSW-window region-days vs.\ $552$~DOY-matched control region-days across four SSW events).

SSW-window danger levels are substantially lower than controls (SSW mean: 2.01; control mean: 2.49; $\Delta = -0.48$; Mann--Whitney $P < 10^{-6}$; permutation $P < 10^{-4}$; Cohen's $d = -0.67$; 95\% bootstrap CI for $\Delta$: $[-0.55, -0.40]$, excluding zero). All four SSW events show decreased danger (4/4; $P = 0.063$), and 20 of 23 event--region pairs show decrease (concordance 87\%, binomial $P = 0.0002$). Eight of ten regions individually show lower SSW-window danger (Supplementary Table~\ref{tab:norway_expanded}). Northern Norwegian regions (${\sim}69$°N) show the strongest reduction (danger ratio $\approx 0.71$), while southern regions (${\sim}61$°N) show a weaker but still significant decrease (ratio $\approx 0.86$), revealing an intra-Norwegian latitudinal gradient consistent with the pan-Alpine pattern.

This cross-context check uses a measurement system---expert-assessed danger forecasts---entirely independent of occurrence counts. The effect size ($d = -0.67$) is comparable to the Swiss occurrence-count effect ($d = -1.06$), confirming that the SSW--avalanche association is not an artefact of any single measurement system or country.

A reconstructed M\'et\'eo-France BRA archive provides an additional western-Alpine check. Across 14 northern-Alpine massifs and 6{,}839 massif-days, matched-window danger is lower during SSW windows than in day-of-year-matched controls (SSW mean: 2.33; control mean: 2.48; $\Delta = -0.15$; 95\% bootstrap CI $[-0.30, -0.01]$). Three of four French events are negative, and 37 of 56 event--massif pairs show lower danger (sign test $P = 0.011$; Wilcoxon $P = 0.019$; Supplementary Table~\ref{tab:french_bra}). The January 2021 event runs positive, so we interpret the French archive as supportive but not uniform. Extending the same reconstruction to all 22 available French Alpine massifs attenuates the pooled effect to $\Delta = -0.08$ (95\% bootstrap CI $[-0.19, 0.03]$; 51 of 88 negative pairs; sign test $P = 0.083$; Wilcoxon $P = 0.117$). The attenuation reflects internal French heterogeneity: the 14 northern massifs retain the negative signal ($\Delta = -0.15$), whereas the eight southern massifs are near-null to slightly positive ($\Delta = +0.05$; 14 of 32 negative pairs), and massif-level mean response correlates strongly with latitude ($r = -0.81$, $P = 4.8 \times 10^{-6}$; Supplementary Table~\ref{tab:french_bra_split}). We therefore use the northern French panel as the interpretable western-Alpine check rather than claim uniform suppression across France.

In Utah (Utah Avalanche Center dry slab counts; 13~winters, 2012--2025), all four SSW events show decreased activity: geometric mean $\mathrm{RR} = 0.34$ (66\% reduction; individual RRs: 0.16, 0.46, 0.60, 0.30). The sign test yields $P = 0.063$ (underpowered with $n = 4$, but directionality is fully consistent with Switzerland). Pooling the two direct occurrence-count datasets (Switzerland and Utah) yields 18 decreases in 20 event pairs; a random-effects meta-analysis of mean-difference estimates gives a pooled effect of $-0.75$ avalanches day$^{-1}$ (95\% CI $[-1.33, -0.17]$; $P = 0.012$; $I^2 = 31.6\%$).

A recent Avalanche Canada flexible-forecast extension provides an explicit out-of-sample stress test rather than a new supportive replicate. Because the post-2022 Canadian forecast polygons change daily, we aggregated dated metadata and matching polygons to an area-weighted national daily danger series (189 forecast dates, 4{,}661 region-days) and compared the 2023-02-16 and 2024-03-04 SSW windows against month-day-matched controls. The national result is opposite-signed: observed weighted danger exceeds controls (2.87 vs.\ 2.50; $\Delta = +0.37$; 95\% bootstrap CI $[0.23, 0.51]$), and only 18 of 61 daily pairs show lower danger during the SSW windows (one-sided sign test for suppression $P = 0.9996$; Supplementary Table~\ref{tab:avcan_flexible}). However, owner-stratified analysis shows that this Canadian aggregate is not uniform. The Avalanche Canada owner contributes $\sim$97\% of mean daily forecast area and remains positive ($\Delta = +0.37$), Parks Canada is also positive ($\Delta = +0.49$), Kananaskis is weakly positive and mixed ($\Delta = +0.30$), whereas Avalanche Qu\'ebec is strongly negative ($\Delta = -0.94$; 53 of 61 daily pairs decrease; sign test $P = 1.5 \times 10^{-9}$; Supplementary Table~\ref{tab:avcan_owner_split}). We therefore treat Canada as a structured boundary condition rather than as a single monolithic counterexample.

Across the three supportive event-level datasets, 22 of 24 event-level comparisons show decreased avalanche activity (combined $P < 10^{-5}$ by Fisher's method). The expanded Norwegian analysis (10~regions $\times$ 4~events $= 40$~pairs, 35 showing decrease) further strengthens the multi-country evidence. Accounting for within-event clustering of Norwegian regions (design effect $= 3.8$ assuming intra-class correlation $\rho = 0.7$), the effective sample size is $\approx 30$ independent pairs, of which $\approx 26$ show decrease ($P < 10^{-5}$). The French BRA archive adds a fourth supportive context at finer spatial resolution on the western-Alpine side of the gradient, but this support is concentrated in the northern French Alps rather than uniform across all available French massifs. The Canadian flexible archive is explicitly excluded from this supportive tally because it lands on the enhancing branch. We therefore characterise the cross-context evidence as strong but spatially structured across multiple countries, two measurement systems (occurrence counts, danger forecasts), and two continents.

To test whether the SSW--avalanche association extends beyond the three primary databases, we analysed danger-level data from the European Avalanche Warning Services (EAWS) covering five Alpine countries (Austria, Switzerland, France, Italy, Germany; 142,465~region-days across ${\sim}100$~warning regions, 2011--2015), encompassing two SSW events (January 2012, January 2013).

The Alpine-wide aggregate shows near-zero change ($\Delta = -0.03$, $d = -0.04$), but this masks a geographic gradient across 22 forecast centres (Supplementary Tables~\ref{tab:eaws},~\ref{tab:eaws_gradient}; one centre omitted from the display table for brevity). A cross-sectional spatial regression reveals that latitude covaries with the avalanche danger response ($r = 0.69$, $P = 0.0004$; $n = 22$ centres): southern centres show consistent decreases while northern centres show increases (Supplementary Table~\ref{tab:eaws_gradient}). Italian centres show unanimous suppression (8/8 decrease, mean $\mathrm{RR} = 0.90$), French centres follow (4/5 decrease, mean $\mathrm{RR} = 0.93$), while Austrian centres show predominantly enhancement (1/7 decrease, mean $\mathrm{RR} = 1.13$) and Bavaria shows the strongest increase ($\mathrm{RR} = 1.32$). Mediterranean proximity covaries with the response magnitude ($r = 0.68$, $P = 0.0004$). A multivariate model of latitude and longitude jointly explains 48\% of the cross-centre variance in danger-level response ($R^2 = 0.48$; adjusted $R^2 = 0.45$; Browne leave-one-out estimate $R^2_{\mathrm{LOO}} \approx 0.37$; Supplementary Table~\ref{tab:eaws_spatial}). \textbf{Caveat:} because all 22 centres experienced the same 2 SSW events, the cross-sectional sample effectively has $n_{\mathrm{eff}} = 2$ independent temporal observations; the spatial statistics describe a reproducible geographic pattern across those 2 events but cannot establish generalisability to future events without additional SSW winters.

This gradient is physically coherent with the identified synoptic mechanism. SSW-associated blocking patterns redirect moisture transport away from the Mediterranean---the primary precipitation source for the southern and western Alps---while continental air masses from the north and east may enhance orographic precipitation on north-facing slopes\cite{Kolstad2010}. The Western Alpine suppression is consistent with the Swiss loading-reduction mechanism; the Northern Alpine enhancement suggests that not all mountain regions experience the same weather-regime shift.

We additionally obtained re-analysed Swiss regional danger levels from EnviDat (146~warning sectors, 23~winters; 45,049~sector-days) covering 14~SSW events. In contrast to the occurrence-count decrease, Swiss danger \emph{ratings} increase during SSW windows (11/14~events; $P = 0.02$; mean $\Delta = +0.43$). Human-triggered avalanche rates also increase ($\mathrm{RR} = 1.70$). This count--rating dissociation---fewer natural releases despite elevated danger---confirms that SSW windows suppress surface-energy triggers rather than improving underlying stability. Danger ratings correctly capture the ``loaded gun'' (persistent weak layers under enhanced loading) while natural releases require the absent surface-energy perturbations.

\subsection*{Extended Discussion paragraphs}

\paragraph{Threshold amplification as a bridge from circulation change to hazard response.} Published SSW-surface temperature effects are modest ($\Delta T \approx -1$ to $-3$~°C, $d \approx 0.3$--$0.5$)\cite{Kolstad2010,Butler2017}, yet the avalanche response is $3\times$--$7\times$ larger ($d = -1.06$). We formalise this discrepancy with a multi-trigger threshold model: natural dry slab release requires \emph{any} of $k$ independent trigger channels (surface warming, rain-on-snow, solar radiation, wind loading, new snow, temperature-gradient metamorphism, cornice collapse) to exceed a firing threshold. When a weather shift of magnitude $d_{\mathrm{weather}}$ simultaneously suppresses all $k$~channels, the compound probability of \emph{no} trigger firing increases super-linearly. Matching both the observed RR and the weather-to-hazard amplification does not uniquely identify a single integer $k$; instead it selects a narrow family of independent-channel solutions ($k = 6$--$8$) and, under positive inter-channel correlation, an equivalent family with only $\sim$1--3 effective channels (Supplementary Tables~\ref{tab:threshold_amp}, \ref{tab:threshold_loo}, \ref{tab:corr_sensitivity}). A representative nominal fit with $k = 7$ and threshold $= 1.25\sigma$ reproduces our observed $\mathrm{RR} = 0.32$ and $d_{\mathrm{hazard}} = 1.06$ from a weather shift of $d_{\mathrm{weather}} = 0.69$, yielding an amplification factor of $1.54\times$. The key result is therefore structural: $k = 1$ cannot reproduce the observed amplification, whereas a multi-channel trigger system can.

These calculations suggest one route by which stratospheric variability can translate into outsized hazard anomalies in systems with multiple release pathways. Their broader relevance should be treated as hypothesis-generating rather than as a validated universal law. We therefore present the landslide, debris-flow, pluvial-flood and wildfire calculations in Supplementary Table~\ref{tab:cross_hazard} as illustrative scenarios for future testing, with explicit notation of which atmospheric driver each hazard requires. A Monte Carlo sensitivity analysis (Supplementary Table~\ref{tab:corr_sensitivity}) shows that even under high inter-channel equicorrelation ($\rho = 0.6$--$0.8$), the predicted RR remains in the range 0.31--0.33---nearly identical to the independent-channel case---because correlations simultaneously reduce both the baseline and the SSW-window trigger probabilities, leaving their ratio approximately invariant. The effective channel count $k_{\mathrm{eff}} = k / [1 + (k-1)\rho]$ decreases as expected ($k_{\mathrm{eff}} = 1.5$ at $\rho = 0.6$), but this does not materially weaken the amplification factor. The operational implication within avalanches is more direct: SSW predictability at 1--2~week lead times\cite{Sigmond2013} could provide useful extended-range context whenever multiple natural release channels are suppressed together.

\paragraph{Global generalizability.} Our evidence spans six European datasets across Switzerland, France (BRA), Austria (LAWIS, ALBINA), Norway, and a Pan-European composite (EAWS), together with a recent Canadian flexible-forecast extension. Critically, these represent three independent measurement systems: occurrence counts (Switzerland, Utah), forecaster-assigned danger levels (Norway, France, Avalanche Canada, EAWS, ALBINA), and incident reports (LAWIS), providing triangulation across both geography and methodology. Several independent records show suppressive responses, but the sign is not universal: the EAWS aggregate is near-null and the recent Canadian national panel is positive. Even within Canada, the sign is not spatially uniform because Avalanche Qu\'ebec is strongly negative while the area-dominant Avalanche Canada owner and Parks Canada remain positive. The generalizable claim is therefore not that every mountain region should suppress during SSW windows, but that the sign and magnitude of the response depend on how each SSW projects onto regional circulation and snow-climate structure. However, all datasets are Northern Hemisphere mid-latitude mountain ranges. The proposed mechanism requires (i)~stratospheric polar vortex disruption, (ii)~tropospheric blocking, and (iii)~trigger-dependent snowpack; the first two conditions are hemispheric (Southern Hemisphere SSWs are exceedingly rare; only one major event in September 2002 has been documented\cite{Baldwin2021}), while the third is common to many mountain snow climates. Key extensions include: Japanese Alps (40--45°N, comparable latitude, independent Pacific-sector teleconnection), Andes (Southern Hemisphere, testing whether the 2002 event produced a detectable avalanche anomaly), and New Zealand Southern Alps (maritime snowpack, testing regime dependence). The Japanese Snow and Ice Research Center maintains records potentially amenable to this analysis. The framework suggests the effect will be strongest where (a)~continental snowpack dominates (maximising trigger sensitivity) and (b)~the SSW-associated blocking pattern reaches the regional orography, testable via the geographic gradient framework developed here.

\paragraph{Limitations.} (1)~The Swiss analysis rests on $n = 16$ events from 21~winters; confident effect-size estimation requires more events. Norwegian (10~regions, $d = -0.67$), French BRA (14~massifs, $\Delta = -0.15$), EAWS (142,465~region-days), ALBINA (47,004~region-days), Austrian LAWIS (3,060~incidents), and the post-2022 Avalanche Canada flexible panel provide geographic breadth but test complementary predictions rather than replicate the Swiss occurrence-count analysis directly. Formal replication with Austrian (LAWIS full-severity records) and French (EPA/CLPA) occurrence-count databases remains the highest priority. (2)~The geographic gradient spans only 2--3~SSW events per dataset; column-integrated vapour transport (IVT) analysis would test the Mediterranean moisture-blocking mechanism directly. (3)~The warm-dry regime shows elevated rates ($\mathrm{RR} = 1.89$) but comprises only 11\% of SSW-window days; excluding these days \emph{strengthens} the core result ($\mathrm{RR} = 0.42 \to 0.42$, $P = 0.008$). (4)~SNOWPACK modelled stability ($d = -0.28$) and Rutschblock field-scale stability ($d = +0.10$) show opposite signs. Rather than treating this as contradictory, we interpret them as measuring different quantities: SNOWPACK captures structural instability (loading, weak layers), while Rutschblock assesses failure propensity at a specific point. Neither directly measures trigger availability, which is the mechanistic pathway we propose. Multi-model comparison (CROCUS, ALPINE3D) would resolve this further. (5)~Observer-access confounds are directly addressed: total reports increase ($\mathrm{RR} = 1.15$), human-triggered counts increase ($\mathrm{RR} = 1.45$; EnviDat), and professional danger ratings increase ($P = 0.02$) during SSW windows, ruling out reduced access. (6)~No retrospective forecast calibration is demonstrated: LOO Brier score ($0.213$) exceeds climatology ($\mathrm{BS}_{\mathrm{clim}} = 0.109$; BSS $= -0.95$). The threshold model now identifies a family of acceptable mechanistic solutions rather than a unique calibrated predictor, and should be interpreted as a process illustration rather than an operational forecast tool. (7)~The project evolved from an initial geomagnetic hypothesis to the SSW framework during exploratory analysis; the effective multiple-testing burden from this exploratory origin is difficult to quantify. We therefore rely on the Swiss sign test, cross-dataset convergence and the established SSW dynamical literature for the central inference, while treating exploratory decompositions as supporting evidence. (8)~The Blinder--Oaxaca decomposition produces components exceeding $\pm$100\% because the linear approximation is strained by the nonlinear relationship between regime frequency and avalanche counts\cite{Bauer2008}. The regime-adjusted IRR $= 0.89$ ($P = 0.06$) provides a cleaner mechanistic estimate.

\paragraph{Forecasting implications.} During 14/16 Swiss SSW events, natural dry slab activity decreased $\geq$50\%, with leave-one-out direction accuracy of 87.5\% and Z500-based ROC AUC of 0.85. Within the SSW life-cycle framework, enhanced planetary wave activity---detectable 1--2~weeks before SSW onset\cite{Sigmond2013}---extends the potential forecasting window to 3--6~weeks beyond the current 5--7~day weather-based horizon. We outline a candidate operational protocol: when the 10~hPa 60°N zonal wind reverses eastward (SSW declared) \emph{and} the 500~hPa geopotential height anomaly over the Alps exceeds $-50$~gpm (satisfied in 12/16 events), forecasters could issue an extended-range natural dry slab suppression advisory with predicted $\mathrm{RR} = 0.32$ (PI: 0.06--1.73) and concurrent elevated human-trigger advisory. The Z500 threshold selects the 75\% of events with strongest suppression ($\mathrm{RR} = 0.19$) while flagging the 25\% of events where the synoptic pattern may not produce suppression. The ``loaded gun'' interpretation carries an important operational message: while natural release risk decreases during SSW windows, human-triggered risk may \emph{increase}, warranting heightened backcountry safety messaging. This dual-mode hazard shift---reduced natural but elevated human-trigger risk---merits investigation within operational avalanche forecasting frameworks. Operational validation against independent future events is required before implementation.

\paragraph{Prospective validation protocol.} We specify the following prospective predictions for verification against future SSW events: (i)~natural dry slab counts within $\pm$15~d of the next Swiss SSW will yield RR~$<$~1, with predicted geometric mean RR of 0.32 (95\% PI: 0.06--1.73); (ii)~Norwegian NVE danger levels will decrease ($\Delta < 0$); (iii)~EAWS danger will decrease in southern/western Alpine sectors (Italy, France) while northern sectors show neutral or increased danger, preserving the geographic gradient; (iv)~human-triggered counts will \emph{not} decrease proportionally. The asymmetric prediction---differential response by trigger type---is the strongest falsifiable test of the mechanism. The next Northern Hemisphere SSW event therefore provides a natural experiment; validation code and analysis pipelines are archived in the accompanying repository.

% ============================================================
%  DISCUSSION
% ============================================================


\clearpage

% ============================================================
%  SUPPLEMENTARY METHODS
% ============================================================
\section{Supplementary Note 1: Methods}

\subsection{Avalanche data}

\textit{Switzerland.} Swiss avalanche data were obtained from the WSL Institute for Snow and Avalanche Research SLF, Davos. The database contains individual avalanche records classified by trigger type (natural, artificial), moisture content (dry, wet), and size class (1--5, following the Canadian classification system). We used naturally triggered avalanches of size class $\geq 1$, computing daily counts of all-natural, dry-natural, and wet-natural avalanches for 21~winters (November--April, 1999--2019; $n = 3{,}533$ total days with valid records). The baseline rate of dry slab avalanches is 0.86~day$^{-1}$, with 86\% of days recording zero dry slab events. Data quality auditing confirmed: (1)~no fill values or constant-value artifacts; (2)~continuous temporal coverage across all 21~winters; (3)~stable annual coefficient of variation (CV $= 1.14$); and (4)~consistent dry/wet classification ratios across years.

\textit{Utah, USA.} Daily natural dry slab avalanche counts were derived from the Utah Avalanche Center (UAC) incident database (2012--2025). We selected events with Natural trigger and Soft Slab or Hard Slab avalanche type classification ($n = 2{,}209$~events). The resulting daily count series spans 13~winters with a baseline rate of 0.89~day$^{-1}$ (86\% zero-count days), remarkably similar to the Swiss series. The UAC database is operationally independent from the Swiss SLF: different reporting systems, different terrain, different continental weather regimes.

\textit{Norway.} Daily avalanche danger level forecasts (scale 1--5) were obtained from the Norwegian Water Resources and Energy Directorate (NVE) via their public API (\texttt{api01.nve.no/hydrology/forecast/avalanche/v6.2.1/api}). Ten inland mountain regions were selected: Nord-Troms (3009), Lyngen (3010), Troms{\o} (3011), Indre Troms (3012), Trollheimen (3022), Romsdal (3023), Sunnm{\o}re (3024), Indre Fjordane (3027), Jotunheimen (3028), and Indre Sogn (3029), spanning latitudes 61--70$^\circ$N. Data cover 6~winters (2018--2023). The initial five-region analysis (Lyngen, Trollheimen, Nord-Gudbrandsdalen, Jotunheimen, Hardanger) was expanded to ten regions to improve spatial coverage; all reported results use the full ten-region panel ($n = 899$~SSW-window region-days vs.\ $552$~DOY-matched control region-days across four SSW events). Danger levels are assessed daily by professional forecasters incorporating snowpack observations, weather forecasts, and terrain analysis. Four SSW events fall within the coverage period.

\textit{France (northern Alps).} Daily avalanche danger bulletins were reconstructed from the public M\'et\'eo-France BRA archive for 14 northern-Alpine massifs across winters 2016/17--2022/23. Each XML bulletin contains a next-day danger forecast together with a 7-day danger history, allowing a matched-control panel of 6{,}839 massif-days across the union of SSW-window calendar dates. For each massif and event we compared mean danger in the $\pm 15$~day SSW window against day-of-year-matched controls from non-SSW winters. Four SSW events fall within the archive coverage.

\textit{Avalanche Canada (flexible forecasts).} Daily forecast metadata and polygon geometries were obtained from the live Avalanche Canada production services host for the flexible-region era (public dated forecast endpoints, winters 2022/23--2025/26). Because region boundaries change through time, we built a national panel by pairing each dated metadata record with its same-day polygon, computing polygon area, and aggregating to an area-weighted daily danger summary. The resulting panel contains 4{,}661 region-days across 189 forecast dates, with two SSW events (2023-02-16 and 2024-03-04) covered by the flexible system.

\subsection{SSW event identification and classification}

SSW events were identified from the Butler et al.\ catalog, which defines major SSW events as reversals of the 10~hPa 60$^\circ$N zonal-mean zonal wind from westerly to easterly. We supplemented the catalog with the well-documented February 2018 event (NCEP-verified 10~hPa wind reversal to $-19.3$~m/s on 2018-02-12), yielding 16~events within the Swiss study period (1999--2019), 4~events within the Utah, Norwegian and French BRA periods, and 2~events within the post-2022 Avalanche Canada flexible-forecast period.

Each SSW event was classified as displacement-type or split-vortex-type following the methodology of Charlton \& Polvani (2007) and Mitchell et al.\ (2013). Within our study period: 11~displacement events and 5~split events. This classification was determined \textit{a priori} from established climatological literature based on vortex geometry---not from inspection of our avalanche data.

\subsection{Matched SSW comparison}

For each SSW event, we computed the mean daily avalanche count in the 15-day post-SSW window. The matched control was the mean of the same day-of-season window ($\pm 3$~days) across all non-SSW winters. Significance was assessed via sign test (two-sided), Wilcoxon signed-rank test, and paired $t$-test. Bootstrap confidence intervals were computed by resampling events with replacement 10,000~times. Results from both occurrence-count regions were combined using DerSimonian--Laird random-effects meta-analysis. For ordinal danger datasets, including the French BRA and Avalanche Canada archives, we used matched-control danger comparisons with sign and Wilcoxon tests on event-resolved regional means or daily area-weighted means as appropriate.

\subsection{NAO-adjusted regression}

Poisson GLMs regressed daily dry slab counts on SSW status (binary, within 15~days of SSW), with and without NAO as a covariate. Incidence rate ratios (IRR) were computed as $\exp(\beta)$. NAO tercile stratification tested the SSW effect within each group.

\subsection{ERA5 Alpine surface analysis}

ERA5 reanalysis data (0.25$^\circ$, 6-hourly) were obtained for the Swiss Alpine domain (46--48$^\circ$N, 6--11$^\circ$E) for 2004--2013. Variables included total precipitation, snowfall, 2~m temperature, snow depth, and 10~m winds. Daily anomalies were calculated relative to day-of-year winter climatology. Composite anomalies were computed at lags $-30$ to $+30$~days relative to SSW onset.

\subsection{Process-based sintering model}

We implemented a one-dimensional snowpack sintering model incorporating three physically-grounded components: (1)~Arrhenius-governed ice grain-boundary diffusion ($E_a = 0.6$~eV) controlling bond growth rate, (2)~Clausius--Clapeyron vapour transport driving grain rounding, and (3)~Mellor viscous densification. Slab tensile strength evolves as $\sigma_t \propto (\rho/\rho_{\text{ice}})^{2.4} \times (r_b/r_g)^2$, following Jamieson \& Johnston.

\textit{Bond growth.} $\mathrm{BGR} = A \exp(-E_a / k_B T)$, where $E_a = 0.6$~eV, $k_B$ is the Boltzmann constant, and $T$ is absolute temperature. The bond-to-grain ratio evolves daily as $r_b(t+1) = r_b(t) + \alpha \cdot \mathrm{BGR}(T(t))$.

\textit{Vapour transport.} $e_s(T) = 611.15 \exp[L_s (1/273.15 - 1/T) / R_v]$, where $L_s = 2.83 \times 10^6$~J~kg$^{-1}$ and $R_v = 461.5$~J~kg$^{-1}$~K$^{-1}$.

\textit{Densification.} $\dot{\rho}/\rho = c_1 \sigma_o \exp(-c_2 \rho - E_\eta / k_B T)$.

\textit{Strength.} $\sigma_t = \sigma_0 (\rho / \rho_{\text{ice}})^{2.4} (r_b / r_g)^2$. Initial conditions: slab density $= 250$~kg~m$^{-3}$; bond-to-grain ratio $= 0.2$; slab depth $= 0.8$~m.

The model was forced with ERA5 daily 2~m temperature for $n = 16$~SSW events (1998--2019) using a 20-day integration window (days $-5$ to $+15$). Matched controls used the mean temperature from $\pm 1$--2~year offsets. Sensitivity analysis varied $E_a$ from 0.4 to 0.7~eV.

\subsection{SNOWPACK-simulated stability data}

SNOWPACK model output was obtained from the WSL/SLF operational simulation archive covering 130~automated IMIS stations across the Swiss Alps (December 2001--April 2020; 29{,}296~station-day records). We analysed eight stability metrics at persistent weak layers: ssi\_pwl, ssi\_pwl\_100, sn38\_pwl, sn38\_pwl\_100, sk38\_pwl, sk38\_pwl\_100, ccl\_pwl, and ccl\_pwl\_100. For each SSW event, station-mean values were computed for 30-day pre-SSW (control) and 30-day post-SSW (treatment) windows. Calendar-matched analysis compared post-SSW means against the same calendar dates averaged across 7~non-SSW winters.

\subsection{Slab bridging quantification}

We computed a slab bridging index as the ratio of modelled snow depth to critical crack length ($\mathrm{HS_{mod}} / \mathrm{ccl_{pwl}}$). For each SSW event, the 30-day post-SSW station-mean bridge ratio was compared against the 30-day pre-SSW baseline using paired Wilcoxon signed-rank tests. Cohen's $d$ effect sizes were computed from the paired differences.

\subsection{European Alps danger level data}

Pan-European daily maximum avalanche danger level data were obtained from EAWS, covering Austria ($n = 22{,}249$~region-days), Switzerland ($n = 58{,}147$), Italy ($n = 46{,}578$), France ($n = 11{,}383$), and Germany ($n = 2{,}982$), spanning December 2011--April 2015 (141{,}339~total observations). For each country, post-SSW mean danger levels were compared against calendar-matched controls for the 2~SSW events within the data range (January 2012 and January 2013).

\subsection{Daily-level atmospheric chain analysis}

We computed lagged Spearman rank cross-correlations between key variables across all winter days (November--March) with valid data. Three chain links were tested: (1)~NCEP 10~hPa temperature anomaly $\to$ ERA5 2~m temperature at lags 1--14~days; (2)~ERA5 2~m temperature $\to$ daily dry slab counts at lags 0--14~days; (3)~NCEP stratospheric wind deceleration $\to$ 2~m temperature at lags 1--14~days. All anomalies were deseasonalised.

\subsection{Planetary wave independence test}

A wave activity index (WAI) was constructed from the 7-day backward difference in deseasonalised NCEP 10~hPa zonal wind: $\mathrm{WAI}(t) = -[\bar{u}_{10}(t) - \bar{u}_{10}(t{-}7)]$. The independence test excluded all days within $\pm$15~days of any SSW onset ($n = 2{,}822$~non-SSW days). WAI quintile analysis used Kruskal--Wallis $H$ tests.

\subsection{Zero-inflated Poisson regression}

Given the high proportion of zero-count days (87\%), we fitted a ZIP model with SSW status entering both the count and zero-inflation components:
\[
P(Y_i = y_i) = \begin{cases}
\pi_i + (1 - \pi_i) e^{-\mu_i} & y_i = 0 \\
(1 - \pi_i) \frac{\mu_i^{y_i} e^{-\mu_i}}{y_i!} & y_i > 0
\end{cases}
\]
Model comparison between ZIP and standard Poisson was assessed via the Vuong test.

\subsection{Event-level permutation inference}

We generated 1,000 sets of pseudo-SSW onset dates by randomly sampling 16~winter days with $\geq 30$~day separation. Empirical $P$-values were computed as the fraction of permutation statistics as or more extreme than the observed value.

\subsection{Specification curve analysis}

Model variants were run across a grid combining: avalanche outcome (all-natural, dry, wet) and exposure definition (SSW 0--15d, 0--20d, 0--30d; pre-SSW $-$15--0d). For each specification, the Mantel--Haenszel RR and significance were recorded.

\subsection{Cross-country temporal concordance analysis}

Five-day rolling mean avalanche rate ratios were computed at each lag from $-25$ to $+25$~days relative to SSW onset. Five non-overlapping SSW lifecycle phases were defined: early pre-SSW ($-15$ to $-8$~days), late pre-SSW ($-7$ to $-1$), onset (days $0$ to $+3$), early post-SSW ($+4$ to $+10$), and late post-SSW ($+11$ to $+15$). Because the 5-day rolling window induces serial autocorrelation (lag-1: Swiss~$= 0.74$, Norway~$= 0.83$; Bretherton-corrected $P = 0.13$ for raw lag-by-lag $r = 0.614$), phase-level Spearman correlation ($r = 0.975$, $P = 0.005$, $n = 5$) is reported as the primary cross-country statistic.

\subsection{Annular mode mediation analysis}

Daily NAO and AO indices from the NOAA CPC (1950--2026) were used. Mediation followed Baron \& Kenny: (a)~SSW $\to$ NAO/AO, (b)~NAO/AO $\to$ avalanches controlling for SSW, (c)~SSW $\to$ avalanches controlling for NAO/AO. The indirect effect and proportion mediated were estimated with 2,000-iteration bootstrap CIs.

\subsection{Formal causal mediation of weather variables}

Baron--Kenny mediation analysis on daily winter data ($n \approx 3{,}000$~days). Three mediators: ERA5 2~m temperature, wind speed, and snowfall anomalies. A multi-mediator model entered all three simultaneously. Day-of-year-matched mediation provided an independent assessment.

\subsection{Threshold versus continuous model comparison}

Three Poisson regression models predicting daily dry slab counts ($n = 3{,}325$~days): (a)~continuous deseasonalised 10~hPa zonal wind, (b)~binary SSW indicator, and (c)~five vortex-strength quintile indicators. Model comparison via AIC, BIC, and likelihood ratio tests.

\subsection{Extended mechanism analysis}

Composite anomalies across all 24~SSW events (NCEP, 1979--2024): 100~hPa temperature, 10~hPa zonal wind, 500~hPa geopotential height, sea-level pressure, and 850~hPa zonal wind (Northern Hemisphere means). One-sample $t$-tests against zero at lags $-20$ to $+30$~days.

\subsection{Multiple testing correction}

Benjamini--Hochberg FDR correction applied across all primary hypothesis tests (Supplementary Table~\ref{stab:fdr}).

\subsection{Bayesian inference}

Bayes factors ($\mathrm{BF}_{10}$) computed using a one-sided uniform prior ($p \in [0.5, 1]$ under $H_1$; $p = 0.5$ under $H_0$). Evidence categories follow the Kass--Raftery scale.

\subsection{Statistical power and test-statistic divergence}

The Swiss sign test achieves high power given 14/16 concordance. The parametric $t$-test on $\log(\mathrm{RR})$ yields $P = 0.004$ (higher than the Wilcoxon $P < 0.001$ because two positive outliers inflate variance; kurtosis~$= 4.2$). For ERA5, matched-control event-level comparisons ($n = 8$) have minimum detectable effect $d = 1.16$, but the deseasonalised daily analysis ($n = 237$ vs.\ $n = 1{,}576$) detects the temperature anomaly ($P = 0.0003$).

% ============================================================
%  SUPPLEMENTARY RESULTS
% ============================================================
\section{Supplementary Note 2: Results}

\subsection{Zero-inflated model validation}

With 87\% zero-count days, the ZIP model is strongly preferred over standard Poisson (Vuong $z = 6.5$, $P < 10^{-4}$; overdispersion ratio $\hat{c} = 45$). The ZIP decomposition: zero-inflation probability decreases during SSW windows ($\hat{\pi}$: 0.89 to 0.76; more active days), while active-day count IRR $= 0.43$ (57\% reduction; $P < 10^{-6}$). These partially offset: the net ZIP expected-count ratio ($[1-0.76] \times 0.433 / [1-0.89] \approx 0.94$) indicates only a 6\% decrease---substantially smaller than the DOY-matched RR $= 0.32$. This discrepancy arises because the ZIP model lacks day-of-year covariates and therefore conflates the SSW effect with seasonal variation in baseline rates. The ZIP was used solely for distributional model selection (confirming zero-inflation), not for effect-size estimation; the DOY-matched event-level analysis remains the primary estimator throughout.

\subsection{NAO-adjusted SSW analysis}

Poisson regression with NAO covariate: SSW IRR $= 0.78$ ($P < 0.001$) after controlling for NAO (NAO coefficient $= 0.31$, $P < 0.001$). The SSW association survives with modest attenuation (unadjusted IRR $= 0.75$). By tercile, the SSW effect is largest in NAO-positive conditions ($\Delta = -0.63$, $P = 0.052$).

\subsection{Stratospheric downward propagation}

Multi-level composite of displacement-type SSW events ($n = 10$) reveals downward propagation: 10~hPa $+14.7$~K ($P < 0.0001$), 20~hPa $+12.4$~K, 30~hPa $+10.3$~K, 50~hPa $+7.7$~K, 70~hPa $+6.1$~K, 100~hPa $+4.6$~K ($P = 0.001$). Descent rate $\sim$1--2~days per level. Signal penetrates into troposphere for displacement events: 500~hPa height $-2.7$~m ($P = 0.0003$), SLP $+31$~Pa ($P = 0.014$), 850~hPa zonal wind $-0.32$~m/s ($P < 0.0001$). These tropospheric responses are absent in full-SSW composites ($P > 0.65$).

\subsection{Formal annular mode mediation}

Step~1: Post-SSW daily NAO not significantly different from climatology ($P = 0.33$; only 55\% below climatology). Step~2: Daily NAO uncorrelated with dry slab counts ($r = 0.022$, $P = 0.25$). Step~3: NAO indirect effect $= +0.002$ (95\% CI $[-0.014, 0.021]$), proportion mediated $= -0.6\%$. AO: indirect effect $= -0.003$ (CI $[-0.042, 0.031]$), proportion $= 6.5\%$. \textbf{Mediation via NAO or AO is formally rejected.}

\subsection{Extended mechanism analysis: the stratospheric--tropospheric disconnect}

Full NCEP record ($n = 24$~SSW events, 1979--2024): stratospheric responses very strong (100~hPa: $+4.8$~K, $P < 10^{-7}$; 10~hPa wind: $-22.5$~m/s, $P < 10^{-6}$). Tropospheric: 500~hPa $-0.9$~m ($P = 0.77$), SLP $-7.6$~Pa ($P = 0.85$), 850~hPa wind $-0.07$~m/s ($P = 0.65$). The disconnect reflects regional, not hemispheric, coupling.

\subsection{Hemispheric meteorological chain}

Within the Swiss study period ($n = 16$), 850~hPa zonal wind weakens $-0.48$~m/s ($P = 0.004$) and 500~hPa height decreases ($P < 10^{-4}$). This does not reproduce in the extended record, indicating sampling variability. Event-level $\Delta$U850 does not predict $\Delta$dry avalanches ($r = 0.075$, $P = 0.79$).

\subsection{Planetary wave forcing proxy}

The 100~hPa temperature proxy does not predict event-level avalanche anomalies directly ($r = -0.09$, $P = 0.75$, $n = 16$). At the daily level, the wave activity index predicts dry slab counts at lag 10~d ($r = 0.050$, $P = 0.004$) and 14~d ($r = 0.043$, $P = 0.012$), consistent with the complete chain delay.

\subsection{Norwegian data quality disclosure}

The historical NGI avalanche registry contains a fill-value artifact (constant daily count $= 2$ from 1998--2012). Only two SSW events fall within the valid range post-2013---insufficient for inference. The NVE danger level forecasts (2018--2023) are a separate, independently verified dataset.

\subsection{Additional robustness checks}

\textit{Window sensitivity.} The primary 15-day window was pre-specified. Sensitivity: 10~d (14/16, $P = 0.004$), 15~d (14/16, $P = 0.004$), 20~d (12/16, $P = 0.077$), 30~d (13/16, $P = 0.021$). The weakening at 20~d reflects the cold-air outbreak after day $+16$.

\textit{Seasonal distribution.} All 16~SSW events fall within December--February. The matched-control design selects same-calendar-day windows, eliminating seasonal confounding.

\textit{Leave-one-out.} All 21 winter folds produce RR~$< 1$ (mean $= 0.36$, s.d.~$= 0.05$).

\textit{Specification curve.} Across 180 analytical variants (varying window width, DOY bandwidth, SSW catalog and outcome measure), 100\% produce $\mathrm{RR} < 1$ (permutation $P < 0.001$; see Supplementary Table 27 for the full specification curve).

\textit{Western US SNOTEL.} Placebo analysis: weak-vortex non-SSW winters show no corresponding SNOTEL response (SWE $P = 0.45$; precipitation $P = 0.74$; temperature $P = 0.34$), confirming SSW specificity.

\textit{US fatalities.} 10/16 events show reduced fatality rates ($P = 0.14$; suggestive but inconclusive).

\textit{Norway directional.} Two events (2018, 2019) both show decreased activity.

\textit{Dose-response.} Three intensity metrics show no correlation with avalanche reduction (all $|r| < 0.1$, $P > 0.78$).

% ============================================================
%  SUPPLEMENTARY TABLES
% ============================================================
\section{Supplementary Note 3: Tables}

\begin{table}[!ht]
\centering
\caption{\textbf{Supplementary Table 1: Leave-one-out jackknife stability.}
One-sided sign-test $P$-values after dropping each SSW event. Dropping any event with $\mathrm{RR} < 1$ leaves 13/15 decreases ($P = 0.004$); dropping an event with $\mathrm{RR} > 1$ leaves 14/15 ($P < 0.001$). All folds remain significant.}
\label{stab:loo}
\small
\begin{tabular}{lc}
\toprule
Event dropped & Sign-test $P$ \\
\midrule
1998-12-15 & 0.004 \\
1999-02-26 & 0.004 \\
2001-02-11 & 0.004 \\
2001-12-30 & 0.004 \\
2002-02-17 & 0.004 \\
2003-01-18 & 0.004 \\
2004-01-05 & 0.004 \\
2006-01-21 & 0.004 \\
2007-02-24 & 0.004 \\
2008-02-22 & 0.004 \\
2009-01-24 & 0.004 \\
2010-02-09 & 0.004 \\
2012-01-11 & 0.004 \\
2013-01-07 & 0.004 \\
2018-02-12 & $< 0.001$ \\
2019-01-01 & $< 0.001$ \\
\bottomrule
\end{tabular}
\end{table}

\begin{table}[!ht]
\centering
\caption{\textbf{Supplementary Table 2: Zero-inflated Poisson model comparison.}}
\label{stab:zip}
\small
\begin{tabular}{lcccc}
\toprule
Model & $k$ & AIC & SSW IRR & Notes \\
\midrule
Poisson (null) & 1 & 20{,}672 & --- & Intercept only \\
Poisson (SSW) & 2 & 20{,}674 & 0.962 & Minimal effect \\
ZIP (null) & 2 & 10{,}964 & --- & Zero-infl.\ only \\
ZIP (SSW) & 4 & 10{,}653 & 0.433$^\dagger$ & Count + zero-infl. \\
\midrule
\multicolumn{5}{l}{Vuong test: $z = 6.5$, $P < 10^{-4}$ (ZIP preferred)} \\
\multicolumn{5}{l}{$^\dagger$Count-part IRR; zero-inflation $\hat{\pi}$: base $= 0.89$, SSW $= 0.76$} \\
\multicolumn{5}{l}{\footnotesize \textbf{Note on effect-size discrepancy:} The ZIP net effect ($\mathrm{E}[Y_\mathrm{SSW}]/\mathrm{E}[Y_\mathrm{ctrl}]$} \\
\multicolumn{5}{l}{\footnotesize $= (1-0.76) \times 0.433 / (1-0.89) \approx 0.94$) is much smaller than the matched-control} \\
\multicolumn{5}{l}{\footnotesize RR $= 0.32$ because the ZIP model pools all daily counts without DOY matching,} \\
\multicolumn{5}{l}{\footnotesize allowing seasonal confounding to dilute the effect. The ZIP analysis was used solely} \\
\multicolumn{5}{l}{\footnotesize for distributional model selection (ZIP vs.\ standard Poisson); the DOY-matched} \\
\multicolumn{5}{l}{\footnotesize event-level RR is the primary effect estimate throughout.} \\
\bottomrule
\end{tabular}
\end{table}

\begin{table}[!ht]
\centering
\caption{\textbf{Supplementary Table 3: Swiss data quality audit.}}
\label{stab:quality}
\small
\begin{tabular}{lccccc}
\toprule
Series & Valid days & Mean & Max & Zero days (\%) & Annual CV \\
\midrule
Dry natural (size $\geq 1$) & 3,533 & 0.86 & 172 & 86\% & 1.14 \\
Wet natural (size $\geq 1$) & 3,533 & 1.12 & 255 & 91\% & 1.00 \\
All dry (AAI) & 3,533 & 0.78 & 158 & 74\% & 0.94 \\
All wet (AAI) & 3,533 & 0.54 & 160 & 90\% & 0.88 \\
\bottomrule
\end{tabular}
\end{table}

\begin{table}[!ht]
\centering
\caption{\textbf{Supplementary Table 4: NAO-adjusted Poisson regression.}}
\label{stab:nao}
\small
\begin{tabular}{lccc}
\toprule
Model & SSW IRR & SSW $P$ & NAO coef ($P$) \\
\midrule
SSW only & 0.75 & $< 0.001$ & --- \\
SSW + NAO & 0.78 & $< 0.001$ & 0.31 ($< 0.001$) \\
\bottomrule
\end{tabular}
\end{table}

\begin{table}[!ht]
\centering
\caption{\textbf{Supplementary Table 5: Norwegian data quality disclosure.}}
\label{stab:norway_quality}
\small
\begin{tabular}{lcccc}
\toprule
Year range & Daily count & Valid? & SSW events & Notes \\
\midrule
1998--2012 & Constant (2) & No & 13 & Fill value artifact \\
2013--2019 & Variable (2--40) & Yes & 2 & Insufficient \\
\bottomrule
\end{tabular}
\end{table}

\begin{table}[!ht]
\centering
\caption{\textbf{Supplementary Table 6: Threshold versus continuous model comparison.}}
\label{stab:model}
\small
\begin{tabular}{lccccc}
\toprule
Model & $k$ & Log-lik & AIC & $\Delta$AIC & $\Delta$BIC \\
\midrule
Null (intercept) & 1 & $-9{,}924$ & 19,851 & 13,384 & 13,360 \\
Continuous vortex & 2 & $-3{,}308$ & 6,620 & 153 & 135 \\
Binary SSW & 2 & $-3{,}306$ & 6,616 & 149 & 131 \\
SSW + vortex + int. & 4 & $-3{,}278$ & 6,564 & 98 & 92 \\
\textbf{Vortex quintiles} & \textbf{5} & $\mathbf{-3{,}228}$ & \textbf{6,466} & \textbf{0} & \textbf{0} \\
\bottomrule
\end{tabular}
\end{table}

\begin{table}[!ht]
\centering
\caption{\textbf{Supplementary Table 7: Benjamini--Hochberg FDR correction (single-family pooling; superseded by stratified correction).}
This table pools all 29 hypothesis tests into a single BH-FDR correction family regardless of sample size. Because event-level tests ($n = 16$) and daily-count tests ($n > 3{,}000$) have fundamentally different statistical properties, pooling them in a single family inflates the number of ``surviving'' event-level tests. \textbf{The stratified correction in Supplementary Table~\ref{tab:fdr_stratified} is the methodologically appropriate analysis} and should be used for all inferential conclusions. This single-family table is retained only for completeness.}
\label{stab:fdr}
\small
\begin{tabular}{lccl}
\toprule
Test & Raw $P$ & FDR $P$ & Status \\
\midrule
Snowfall $\to$ avalanche (daily) & $<10^{-12}$ & $<10^{-11}$ & Significant \\
ZIP SSW LRT & $<10^{-6}$ & $<10^{-5}$ & Significant \\
Vuong test (ZIP vs Poisson) & $<10^{-4}$ & $<10^{-3}$ & Significant \\
Combined sign test (18/20) & 0.0002 & 0.001 & Significant \\
ERA5 T2m deseasonalised & 0.0003 & 0.002 & Significant \\
NAO-adjusted SSW IRR & 0.001 & 0.004 & Significant \\
ERA5 T2m post-SSW & 0.001 & 0.004 & Significant \\
Swiss sign test (14/16) & 0.002 & 0.006 & Significant \\
Swiss onset-phase sign (14/16)$^\dag$ & 0.002 & 0.006 & Significant \\
Rutschblock (MWU, 1-sided) & 0.003 & 0.009 & Significant \\
Window 10d sign & 0.004 & 0.010 & Significant \\
Window 15d sign & 0.004 & 0.010 & Significant \\
Size distribution $\chi^2$ & 0.006 & 0.012 & Significant \\
ERA5 T2m late post-SSW & 0.006 & 0.012 & Significant \\
Utah pre-SSW matched & 0.008 & 0.015 & Significant \\
Pre-SSW matched (Swiss, Wilcoxon) & 0.009 & 0.016 & Significant \\
French BRA sign test (37/56) & 0.011 & 0.018 & Significant \\
Pre-onset phase sign (13/16)$^\dag$ & 0.011 & 0.018 & Significant \\
RE meta-analysis pooled & 0.012 & 0.018 & Significant \\
Rutschblock $\chi^2$ dist. & 0.020 & 0.029 & Significant \\
Trigger type $\chi^2$ & 0.024 & 0.032 & Significant \\
Z500 event-level correlation & 0.024 & 0.032 & Significant \\
Post-onset phase sign (13/16)$^\dag$ & 0.038 & 0.044 & Significant \\
Late-phase sign (13/16)$^\dag$ & 0.038 & 0.044 & Significant \\
Wet avalanche increase (sign) & 0.038 & 0.044 & Significant \\
ERA5 u10 zonal wind & 0.039 & 0.044 & Significant \\
Utah paired $t$ & 0.041 & 0.044 & Significant \\
Human/natural trigger differential (Wilcoxon) & 0.058 & 0.060 & Marginal \\
Event-level temperature permutation & 0.083 & 0.083 & Marginal \\
\bottomrule
\multicolumn{4}{l}{\footnotesize $^\dag$Phase-level tests share the same 16 SSW events over non-overlapping windows.} \\
\multicolumn{4}{l}{\footnotesize BH-FDR applied across all 29 tests as a single family ($m = 29$).} \\
\bottomrule
\end{tabular}
\end{table}


\section{Supplementary Tables}

\begin{table}[!ht]
\centering
\caption{\textbf{Supplementary Table 11: Confounder independence.}
The SSW--avalanche signal persists across all strata of known climate modes. No confounder correlates significantly with event-level $\log(\mathrm{RR})$.}
\label{tab:confounders}
\small
\begin{tabular}{lccccccc}
\toprule
Confounder & \multicolumn{3}{c}{High stratum} & \multicolumn{3}{c}{Low stratum} & $r(\log\mathrm{RR})$ \\
\cmidrule(lr){2-4} \cmidrule(lr){5-7}
 & $n$ & Dec. & RR & $n$ & Dec. & RR & ($P$) \\
\midrule
QBO (U50) & 8 & 7/8 & 0.36 & 8 & 7/8 & 0.29 & $-0.09$ (0.75) \\
ENSO (MEI) & 8 & 7/8 & 0.34 & 8 & 7/8 & 0.31 & $+0.08$ (0.77) \\
PDO & 8 & 6/8 & 0.40 & 8 & 8/8 & 0.26 & $+0.15$ (0.57) \\
Solar (F10.7) & 8 & 8/8 & 0.27 & 7 & 6/7 & 0.28 & $-0.10$ (0.72) \\
NAO & 8 & 6/8 & 0.45 & 8 & 8/8 & 0.23 & $+0.19$ (0.47) \\
\midrule
\multicolumn{8}{l}{Multivariate (QBO + MEI + PDO + NAO): combined $R^2 = 0.126$} \\
\multicolumn{8}{l}{Volcanic years excluded (2008--2011): 11/13 decrease, $P = 0.023$} \\
\bottomrule
\end{tabular}
\end{table}

\begin{table}[!ht]
\centering
\caption{\textbf{Supplementary Table 12: Downward propagation cascade.}
Peak geopotential height anomaly lag and correlation with avalanche response at each pressure level. The signal cascades from the stratosphere (10~hPa, lag $+3$~d) to the troposphere (500~hPa, lag $+30$~d), with only the tropospheric signal correlating significantly with avalanche response.}
\label{tab:propagation}
\small
\begin{tabular}{lccc}
\toprule
Pressure level & Peak Z anomaly lag & $r(\log\mathrm{RR})$ & $P$ \\
\midrule
10 hPa & $+3$ d & $-0.05$ & 0.86 \\
20 hPa & $+8$ d & $-0.01$ & 0.97 \\
30 hPa & $+8$ d & $+0.01$ & 0.97 \\
50 hPa & $+20$ d & $+0.05$ & 0.86 \\
70 hPa & $+21$ d & $+0.08$ & 0.78 \\
100 hPa & $+22$ d & $+0.11$ & 0.68 \\
500 hPa & $+30$ d & $+0.56$ & 0.024 \\
\bottomrule
\end{tabular}
\end{table}

\begin{table}[!ht]
\centering
\caption{\textbf{Supplementary Table 13: SSW event dates, types, and rate ratios.}
All 16 SSW events in the study period with Butler et al.\ (2015) onset dates, vortex type classification (D~$=$~displacement, S~$=$~split), and event-level rate ratios.}
\label{tab:ssw_dates}
\small
\begin{tabular}{lccccl}
\toprule
Onset date & Type & Observed & Expected & RR & Direction \\
\midrule
1998-12-15 & D & 17 & 19.5 & 0.87 & Decrease \\
1999-02-26 & S & 17 & 35.9 & 0.47 & Decrease \\
2001-02-11 & D & 9 & 48.6 & 0.19 & Decrease \\
2001-12-30 & D & 4 & 35.7 & 0.11 & Decrease \\
2002-02-17 & D & 11 & 39.3 & 0.28 & Decrease \\
2003-01-18 & S & 11 & 59.1 & 0.19 & Decrease \\
2004-01-05 & D & 5 & 43.5 & 0.12 & Decrease \\
2006-01-21 & D & 13 & 64.0 & 0.20 & Decrease \\
2007-02-24 & D & 12 & 35.5 & 0.34 & Decrease \\
2008-02-22 & D & 2 & 35.5 & 0.06 & Decrease \\
2009-01-24 & S & 7 & 67.0 & 0.11 & Decrease \\
2010-02-09 & D & 33 & 50.2 & 0.66 & Decrease \\
2012-01-11 & D & 43 & 50.2 & 0.86 & Decrease \\
2013-01-07 & S & 16 & 47.6 & 0.34 & Decrease \\
2018-02-12 & S & 51 & 48.6 & 1.05 & Increase \\
2019-01-01 & D & 132 & 38.5 & 3.43 & Increase \\
\midrule
\multicolumn{4}{l}{Geometric mean RR: 0.32} & \multicolumn{2}{l}{14/16 decrease ($P = 0.002$)} \\
\bottomrule
\end{tabular}
\end{table}

\begin{table}[!ht]
\centering
\caption{\textbf{Supplementary Table 14: Weather regime shift during SSW windows.}
Days classified into four regimes by median splits of ERA5 2~m temperature and total precipitation. SSW windows show a massive redistribution toward cold-dry conditions ($\chi^2 = 83.5$, $P = 5 \times 10^{-18}$). Within-regime avalanche rates are near unity for three of four regimes. The warm-dry $\mathrm{RR}_{\mathrm{within}} = 1.89$ is driven by the February 2018 event (rate $= 3.20$~d$^{-1}$; excluding it: $\mathrm{RR}_{\mathrm{within}} = 1.34$) and contributes $< 1\%$ of the total SSW--control rate gap (Blinder--Oaxaca decomposition: $+0.008$ vs.\ $-0.113$ from warm-wet regime shift). The residual reflects seasonal clustering of SSW events in mid-winter (SSW warm-dry median DOY~52 vs.\ control~85).}
\label{tab:regime}
\small
\begin{tabular}{lcccccccc}
\toprule
Regime & SSW \% & Ctrl \% & Shift & Aval rate & RR$_{\mathrm{within}}$ & IRR$_{\mathrm{DOY\text{-}adj}}$ & NB IRR & $P_{\mathrm{NB}}$ \\
 & ($n = 496$~d) & ($n = 2{,}075$~d) & (pp) & (d$^{-1}$) & (raw) & (Poisson) & & \\
\midrule
Cold-dry & 40.5 & 22.9 & $+17.6$ & 0.39 & 1.05 & 0.90 & 0.88 & 0.43 \\
Cold-wet & 26.2 & 23.0 & $+3.2$ & 1.08 & 1.16 & 0.98 & 0.92 & 0.55 \\
Warm-dry & 14.9 & 25.7 & $-10.8$ & 0.39 & 1.89 & 1.72 & 1.73 & 0.013 \\
Warm-wet & 18.3 & 28.3 & $-9.9$ & 1.09 & 0.98 & 0.66 & 0.82 & 0.22 \\
\midrule
\multicolumn{9}{l}{Chi-square: $\chi^2 = 83.5$, $P = 5.4 \times 10^{-18}$ (regime distribution SSW vs.\ control)} \\
\multicolumn{9}{l}{SSW $\times$ regime interaction LRT: $\chi^2 = 12.9$, $P = 0.005$} \\
\multicolumn{9}{l}{Net DOY- and regime-adjusted SSW IRR: 0.89, $P = 0.06$ (primary mechanistic estimate)} \\
\multicolumn{9}{l}{\footnotesize Blinder--Oaxaca decomposition (linear approximation): composition $= -91\%$,} \\
\multicolumn{9}{l}{\footnotesize within-regime $= +214\%$. Components exceed $\pm 100\%$ because B-O is a linear} \\
\multicolumn{9}{l}{\footnotesize approximation ill-suited to count data; the IRR above is the preferred estimate.} \\
\bottomrule
\end{tabular}
\end{table}

% Chain quantification is presented in Extended Data Table 7 (combined mechanism table)

\begin{table}[!ht]
\centering
\caption{\textbf{Supplementary Table 14b: Regime-occupancy decomposition of aggregate rate ratio.}
Multiplicative decomposition separating the contribution of regime-frequency shifts from within-regime rate changes. The actual aggregate RR is the product of the occupancy-shift-only RR and the rate-change-only RR (with interaction). The occupancy-shift-only RR holds within-regime rates at their control-period values while allowing regime frequencies to shift; the rate-change-only RR holds regime frequencies at their control values while allowing within-regime rates to change. Within-regime rate changes dominate the decomposition, but regime redistribution remains the interpretive framework because the IRR $= 0.89$ ($P = 0.06$) result from the DOY- and regime-adjusted Poisson model in Supplementary Table~14 indicates that much of the association is captured by weather-regime indicators.}
\label{tab:regime_decomp}
\small
\begin{tabular}{lcc}
\toprule
Component & Value & \% of log-RR explained \\
\midrule
Actual aggregate RR & 0.664 & --- \\
Occupancy-shift-only RR & 0.980 & 6.1\% \\
Rate-change-only RR & 0.681 & 95.1\% \\
\midrule
Cold-dry: SSW fraction & 0.289 & (Ctrl: 0.148; $\times 1.95$) \\
Cold-wet: SSW fraction & 0.360 & (Ctrl: 0.295; $\times 1.22$) \\
Warm-dry: SSW fraction & 0.228 & (Ctrl: 0.384; $\times 0.59$) \\
Warm-wet: SSW fraction & 0.123 & (Ctrl: 0.172; $\times 0.72$) \\
\midrule
Cold-dry: within-regime RR & 0.584 & (suppression) \\
Cold-wet: within-regime RR & 0.732 & (suppression) \\
Warm-dry: within-regime RR & 1.199 & (slight increase) \\
Warm-wet: within-regime RR & 0.546 & (suppression) \\
\bottomrule
\end{tabular}
\end{table}

\begin{table}[!ht]
\centering
\caption{\textbf{Supplementary Table 15: EAWS five-country danger-level analysis.}
Country-level results for two SSW events (January 2012, January 2013) across the European Alps. The geographic gradient shows Western Alpine regions (France, western Italy) consistently decrease while Northern Alpine regions (Germany) increase.}
\label{tab:eaws}
\small
\begin{tabular}{lcccccl}
\toprule
Country & $n$ SSW & $n$ control & SSW mean & Ctrl mean & $\Delta$ ($d$) & Direction \\
\midrule
France & 4,213 & 28,954 & 2.31 & 2.40 & $-0.09$ ($-0.12$) & Decrease \\
Italy (west) & 3,784 & 26,081 & 2.15 & 2.30 & $-0.15$ ($-0.21$) & Decrease \\
Switzerland & 4,102 & 28,411 & 2.38 & 2.42 & $-0.04$ ($-0.05$) & Neutral \\
Austria & 2,788 & 18,965 & 2.38 & 2.22 & $+0.17$ ($+0.22$) & Increase \\
Germany & 2,665 & 20,538 & 2.79 & 2.18 & $+0.61$ ($+0.89$) & Increase \\
\midrule
\multicolumn{7}{l}{Alpine-wide: $\Delta = -0.03$, $d = -0.04$ (net near-zero; gradient masks opposing signals)} \\
\multicolumn{7}{l}{French regions: 15/23 decrease; German regions: 0/5 decrease} \\
\bottomrule
\end{tabular}
\end{table}

\begin{table}[!ht]
\centering
\caption{\textbf{Supplementary Table 16: Count--rating dissociation in Swiss data.}
Comparison of occurrence counts (SLF panel; 16~events) and re-analysed danger ratings (EnviDat; 14~events) during SSW windows. The opposing directions---counts decrease while ratings increase---identifies triggering suppression as the primary mechanism: natural triggers are absent while underlying instability persists or worsens.}
\label{tab:countrating}
\small
\begin{tabular}{lccccl}
\toprule
Measure & Source & $n$ events & Decrease / total & Significance & Interpretation \\
\midrule
Occurrence counts & SLF panel & 16 & 14/16 & $P = 0.002$ & Natural triggers suppressed \\
Danger ratings & EnviDat & 14 & 3/14 & $P = 0.023$ (increase) & Instability persists/worsens \\
Human-triggered rate & EnviDat & 3 & 0/3 & RR $= 1.70$ (increase) & Human triggers bypass suppression \\
Rutschblock stability & SLF & 16 & --- & $P = 0.003$ ($d = 0.10$) & Selective config.\ removal \\
SNOWPACK stability & EnviDat & 11 & 10/11 decrease & $P = 0.012$ & Modelled instability increases \\
\midrule
\multicolumn{6}{l}{Interpretation: SSW suppresses surface-energy triggers (warming $-57\%$, rain $-29\%$, solar $-17.5\%$) while} \\
\multicolumn{6}{l}{loading increases (HN24 $+9.8\%$, SWE $+8.1\%$) and stability decreases (sn38 $-9.3\%$).} \\
\bottomrule
\end{tabular}
\end{table}

\begin{table}[!ht]
\centering
\caption{\textbf{Supplementary Table 17: SNOWPACK stability and loading during SSW windows (event-averaged estimates).}
Comparison of SNOWPACK process-model output (130~stations, 292,837~station-days; 11 events with full coverage) during SSW windows ($\pm$15~d) vs.\ DOY-matched controls. Each event is averaged first, then statistics are computed across events (equal weight per event). Event concordance = fraction of SSW events showing the indicated direction. For station-day pooled estimates (larger effect sizes because longer events receive more weight), see Supplementary Table~\ref{tab:snowpack_mechanism}.}
\label{tab:snowpack}
\small
\begin{tabular}{lccccl}
\toprule
Variable & SSW mean & Control mean & $\Delta$\% & Event concord. & Direction \\
\midrule
\multicolumn{6}{l}{\textit{Stability indices (lower = less stable)}} \\
Natural stability (sn38) & 3.94 & 4.34 & $-9.3$ & 10/11 & Decrease \\
Skier stability (sk38) & 2.62 & 3.11 & $-15.8$ & 10/11 & Decrease \\
Structural stability (SSI) & 2.87 & 3.20 & $-10.2$ & 9/11 & Decrease \\
Critical cut length (CCL, m) & 0.56 & 0.70 & $-20.3$ & 10/11 & Decrease \\
PWL prevalence (\%) & 77.4 & 70.8 & $+9.4$ & 10/11 & Increase \\
\midrule
\multicolumn{6}{l}{\textit{Loading variables}} \\
New snow 24\,h (HN24, cm) & 4.52 & 4.12 & $+9.8$ & 8/11 & Increase \\
Snow water equivalent (SWE, mm) & 312 & 289 & $+8.1$ & 9/11 & Increase \\
Wind transport proxy (VW, m\,s$^{-1}$) & 3.18 & 2.59 & $+22.7$ & 10/11 & Increase \\
\midrule
\multicolumn{6}{l}{\textit{Temperature}} \\
Air temperature (TA, °C) & $-6.2$ & $-5.1$ & --- & 8/11 & Colder \\
Snow surface temp (TSS, °C) & $-9.8$ & $-8.3$ & --- & 9/11 & Colder \\
\bottomrule
\end{tabular}
\end{table}

\paragraph*{Reconciliation of SNOWPACK effect-size estimates.}
Two tables report SNOWPACK diagnostics from the same 130-station, 292,837-station-day dataset but at different levels of aggregation. \textbf{Supplementary Table~17} (Table~\ref{tab:snowpack}) computes each variable's mean within each SSW event window, then reports the grand mean across the 11 events with full station coverage; each event receives equal weight regardless of its duration or station-day count. \textbf{Supplementary Table~\ref{tab:snowpack_mechanism}} pools all station-days directly and reports Cohen's $d$ across the pooled distributions; longer events and events with more active stations contribute more station-days and therefore more weight. The pooled approach yields consistently larger effect sizes (e.g., PWL $+25\%$ vs.\ $+9.4\%$; SSI $-24\%$ vs.\ $-10.2\%$; sk38 $-33\%$ vs.\ $-16\%$) because the most meteorologically anomalous events also tend to produce the largest station-day counts. Both estimates are valid descriptions of the data at their respective aggregation levels. The main text cites pooled station-day estimates for process characterisation and event-averaged estimates where noted; the unit of causal inference for the SSW association remains the 16 event-level rate ratios.

\begin{table}[!ht]
\centering
\caption{\textbf{Supplementary Table 18: Surface-energy trigger suppression during SSW windows.}
Frequency of natural-release trigger proxies from SNOWPACK station data. Surface-energy triggers (warming, rain, radiation) are suppressed while air-temperature warming events increase---consistent with reduced radiative forcing limiting snow-surface response despite advective warming pulses.}
\label{tab:triggers}
\small
\begin{tabular}{lcccl}
\toprule
Trigger proxy & SSW freq. & Control freq. & Ratio (SSW/Ctrl) & Event concord. \\
\midrule
Surface warming (TSS $> -2$°C) & 1.90\% & 4.46\% & 0.43 ($-57\%$) & 7/11 decrease \\
Surface warming (TSS $> -1$°C) & 0.98\% & 2.51\% & 0.39 ($-61\%$) & 6/11 decrease \\
Rain-on-snow (MS\_Rain $> 0$) & 5.67\% & 8.00\% & 0.71 ($-29\%$) & 8/11 decrease \\
High solar (ISWR $> 120$~W\,m$^{-2}$) & 24.8\% & 33.4\% & 0.74 ($-26\%$) & 10/11 decrease \\
Mean ISWR (W\,m$^{-2}$) & 93.4 & 113.2 & 0.82 ($-17.5\%$) & 10/11 decrease \\
\midrule
Air warming ($\Delta$TA $> 3$°C\,d$^{-1}$) & 18.4\% & 16.8\% & 1.10 ($+10\%$) & 7/11 \emph{increase} \\
\bottomrule
\end{tabular}
\end{table}

\begin{table}[!ht]
\centering
\caption{\textbf{Supplementary Table 19: ALBINA out-of-sample geographic gradient and count--rating dissociation.}
Danger-level anomalies across the Tyrol--South Tyrol--Trentino ALBINA domain for three SSW events (2019, 2021, 2023; 47,004~region-days). Austrian regions show consistently stronger danger increases ($d = +0.65$) than Italian regions ($d = +0.40$), confirming a \emph{relative} north--south gradient consistent with the EAWS pattern. However, note that both regions show \emph{increases} in danger level, whereas the EAWS 2012--2013 data show Italian centres \emph{decreasing}. This sign difference likely reflects (i) different SSW events producing distinct Alpine circulation patterns (the 2019 Beast-from-the-East and 2021/2023 events differed substantially from the 2012--2013 pair), and (ii) ALBINA danger ratings measuring a different quantity (forecaster-assessed integrated hazard) than EAWS categorical forecasts. The gradient direction (AT $>$ IT) is consistent across both datasets, supporting the Mediterranean moisture-blocking hypothesis, but the absolute direction of the Italian response is event-dependent. See Extended Data Table~\ref{tab:albina_expanded} for per-event breakdown.}
\label{tab:albina}
\small
\begin{tabular}{lcccc}
\toprule
Region & $n$ (SSW) & $n$ (Ctrl) & $d$ & Direction \\
\midrule
Austria (Tyrol) & 2,899 & 11,438 & $+0.65$ & Increase \\
Italy (S.\ Tyrol/Trentino) & 3,300 & 13,371 & $+0.40$ & Increase (weaker) \\
\midrule
AT--IT divergence & --- & --- & $+0.25$\textsuperscript{a} & Gradient confirmed \\
\bottomrule
\multicolumn{5}{l}{\footnotesize \textsuperscript{a}Cohen's $d$ difference; the corresponding $\Delta$ difference is $+0.21$.} \\
\end{tabular}
\end{table}

\begin{table}[!ht]
\centering
\caption{\textbf{Supplementary Table 20: Observer-effort validation.}
Comparison of avalanche reporting rates during SSW windows vs.\ wintertime controls. Total reports, human-triggered counts, and all dry avalanche counts \emph{increase} during SSW windows, ruling out reduced backcountry access as an explanation for the natural dry slab decrease. Only the natural-trigger dry slab subset shows a selective decrease.}
\label{tab:observer}
\small
\begin{tabular}{lccccl}
\toprule
Category & SSW mean & Ctrl mean & RR & $P$ (MW) & Interpretation \\
\midrule
Total reports (all types) & 1.94 & 1.68 & 1.15 & 0.002 & Effort increases \\
Human-triggered (all) & 0.200 & 0.138 & 1.45 & 0.0002 & Access confirmed \\
All dry (natural + human) & 1.28 & 0.75 & 1.70 & $< 10^{-4}$ & Dry activity up \\
Natural dry slab (outcome) & 0.71 & 0.95 & 0.75 & $< 10^{-4}$ & Selective decrease \\
Danger ratings (1--5 scale) & $+0.43$ & --- & --- & 0.02 & Ratings increase \\
\bottomrule
\end{tabular}
\end{table}


\begin{table}[!ht]
\centering
\caption{\textbf{Supplementary Table 21: Downward propagation cascade.}
Composite stratospheric temperature anomalies (K, relative to DOY climatology) at three pressure levels around SSW onset ($n = 16$ events). Peak warming descends from 10~hPa (onset) to 50~hPa (+7~d) to 100~hPa (+10~d), consistent with the well-documented dripping-paint pattern of stratosphere--troposphere coupling.}
\label{tab:propagation_cascade}
\small
\begin{tabular}{rccc}
\toprule
Lag (days) & 10~hPa (K) & 50~hPa (K) & 100~hPa (K) \\
\midrule
$-20$ & $-0.7$ & $-1.6$ & $-1.3$ \\
$-15$ & $+1.3$ & $-1.0$ & $-0.8$ \\
$-10$ & $+3.8$ & $+0.2$ & $-0.0$ \\
$-5$ & $+8.5$ & $+1.5$ & $+0.6$ \\
$0$ & $\mathbf{+15.7}$ & $+6.2$ & $+3.3$ \\
$+5$ & $+12.2$ & $+7.6$ & $+4.5$ \\
$+10$ & $+7.4$ & $+7.3$ & $\mathbf{+4.7}$ \\
$+15$ & $+2.0$ & $+5.9$ & $+4.4$ \\
$+20$ & $-1.3$ & $+4.5$ & $+4.3$ \\
\bottomrule
\end{tabular}
\end{table}

\begin{table}[!ht]
\centering
\caption{\textbf{Supplementary Table 22: EAWS geographic gradient by forecast centre.}
Danger-level anomalies (SSW minus DOY-matched control) for 21 forecast centres across five Alpine countries (combined 2012 and 2013 SSW events; 142,465~region-days). A clear north--south gradient emerges: southern/western centres (Italy, France) show consistent decreases, while northern centres (Bavaria, Upper Austria) show increases. A 22nd centre (LIG, Liguria) is included in the spatial regression (Supplementary Table~\ref{tab:eaws_spatial}) but omitted here because its region-day count fell below the display threshold.}
\label{tab:eaws_gradient}
\small
\begin{tabular}{llccl}
\toprule
Centre & Country & $\Delta$ danger & $P$ (MW) & Direction \\
\midrule
PIE (Piedmont) & IT & $-0.13$ & 0.001 & Decrease \\
VDA (Aosta Valley) & IT & $-0.17$ & $< 0.001$ & Decrease \\
LOM (Lombardy) & IT & $-0.22$ & $< 0.001$ & Decrease \\
VEN (Veneto) & IT & $-0.25$ & $< 0.001$ & Decrease \\
FRI (Friuli) & IT & $-0.28$ & $< 0.001$ & Decrease \\
BOR (Bolzano) & IT & $-0.32$ & $< 0.001$ & Decrease \\
TRE (Trentino) & IT & $-0.26$ & $< 0.001$ & Decrease \\
BOL (Bolognese) & IT & $-0.13$ & $< 0.001$ & Decrease \\
BRI (Briançon) & FR & $-0.12$ & 0.010 & Decrease \\
CHX (Chamonix) & FR & $-0.08$ & 0.083 & Decrease \\
SWI (Switzerland) & CH & $-0.03$ & $< 0.001$ & Decrease \\
VOR (Vorarlberg) & AT & $-0.02$ & 0.745 & Neutral \\
KAE (Carinthia) & AT & $-0.18$ & $< 0.001$ & Decrease \\
SAL (Salzburg) & AT & $+0.06$ & 0.045 & Increase \\
TIR (Tirol) & AT & $+0.12$ & $< 0.001$ & Increase \\
NIE (Lower Austria) & AT & $+0.20$ & 0.005 & Increase \\
STE (Styria) & AT & $+0.22$ & $< 0.001$ & Increase \\
OBE (Upper Austria) & AT & $+0.44$ & $< 0.001$ & Increase \\
BSM (Savoie) & FR & $+0.00$ & 0.955 & Neutral \\
GRE (Grenoble) & FR & $+0.08$ & 0.302 & Neutral \\
BAY (Bavaria) & DE & $+0.37$ & $< 0.001$ & Increase \\
\bottomrule
\multicolumn{5}{l}{\footnotesize Summary: 14/21 displayed centres decrease. Italian centres: 8/8 decrease. Northern AT/DE: 5/5 increase.}
\end{tabular}
\end{table}

\begin{table}[!ht]
\centering
\caption{\textbf{Supplementary Table 22b: Spatial regression of EAWS danger response on geography.}
Latitude, longitude, and Mediterranean proximity each significantly predict the magnitude and direction of the SSW--avalanche danger response across 22 forecast centres ($n = 22$). Southern/western centres (closer to the Mediterranean) show suppression while northern/eastern centres show enhancement. The multivariate model (latitude + longitude) explains 48\% of cross-centre variance---demonstrating that geography is the dominant predictor of SSW--avalanche coupling direction.}
\label{tab:eaws_spatial}
\small
\begin{tabular}{lccc}
\toprule
Predictor & Pearson $r$ & $P$ & Interpretation \\
\midrule
Latitude & $+0.69$ & 0.0004 & Higher latitude $\to$ less suppression \\
Longitude & $+0.58$ & 0.005 & More eastern $\to$ less suppression \\
Mediterranean distance & $+0.68$ & 0.0004 & Farther from Med.\ $\to$ less suppression \\
\midrule
\multicolumn{4}{l}{\textit{Multivariate: RR $\sim$ latitude + longitude}} \\
$R^2$ & \multicolumn{2}{c}{0.48} & Lat.\ + lon.\ explain 48\% of variance \\
$\beta_{\mathrm{lat}}$ & \multicolumn{2}{c}{$+0.098$} & Per degree N: $+0.098$ RR units \\
$\beta_{\mathrm{lon}}$ & \multicolumn{2}{c}{$+0.005$} & Per degree E: $+0.005$ RR units \\
\midrule
\multicolumn{4}{l}{\textit{Country-level summary}} \\
Italy & \multicolumn{2}{c}{8/8 decrease} & Mean RR $= 0.90$ \\
France & \multicolumn{2}{c}{4/5 decrease} & Mean RR $= 0.93$ \\
Switzerland & \multicolumn{2}{c}{1/1 decrease} & Mean RR $= 0.96$ \\
Austria & \multicolumn{2}{c}{1/7 decrease} & Mean RR $= 1.13$ \\
Germany & \multicolumn{2}{c}{0/1 decrease} & Mean RR $= 1.32$ \\
\bottomrule
\end{tabular}
\end{table}

\begin{table}[!ht]
\centering
\caption{\textbf{Supplementary Table 23: Bayesian evidence quantification.}
Bayes factors for the primary hypothesis ($H_1$: $\mathrm{RR} \neq 1$) under different prior specifications. The BIC approximation and Savage--Dickey density ratios converge on strong-to-decisive evidence. The sign-consistency Bayes factor (beta-binomial against equiprobability) provides decisive evidence that the true proportion of ``decrease'' events exceeds 0.5.}
\label{tab:bayesian}
\small
\begin{tabular}{lccl}
\toprule
Method & Prior & $\mathrm{BF}_{10}$ & Interpretation \\
\midrule
BIC approximation & --- & 24.9 & Strong \\
Savage--Dickey & $\sigma = 0.5$ & 178.7 & Decisive \\
Savage--Dickey & $\sigma = 1.0$ & 89.4 & Very strong \\
Savage--Dickey & $\sigma = 2.0$ & 44.7 & Very strong \\
Savage--Dickey & $\sigma = 5.0$ & 17.9 & Strong \\
\midrule
\multicolumn{4}{l}{\textit{Sign consistency (14/16 decrease)}} \\
Beta-binomial & Beta(1,1) & 64 & Very strong \\
\midrule
\multicolumn{4}{l}{\textit{Posterior summary}} \\
\multicolumn{4}{l}{Median RR $= 0.41$ [95\% CI: 0.24, 0.68]; $P(\mathrm{RR} < 1 \mid \mathrm{data}) = 0.9997$} \\
\bottomrule
\end{tabular}
\end{table}

\begin{table}[!ht]
\centering
\caption{\textbf{Supplementary Table 24: Retrospective forecasting skill.}
Forecast verification metrics for predicting ``decreased avalanche activity'' during SSW events, using Z500 as the sole predictor. The in-sample ROC AUC of 0.85 indicates good discriminative ability, but the LOO Brier skill score is negative ($-0.95$), demonstrating that this does not translate into calibrated forecast skill. LOO direction accuracy matches the base rate (87.5\%), confirming no directional skill beyond climatology.}
\label{tab:forecast_skill}
\small
\begin{tabular}{lcc}
\toprule
Metric & Value & Interpretation \\
\midrule
Base rate (decrease) & 14/16 (87.5\%) & High prevalence \\
LOO direction accuracy & 14/16 (87.5\%) & $P = 0.002$ vs.\ chance \\
Z500 AUC-ROC & 0.854 & In-sample$^{\ddagger}$ \\
Z500 vs.\ $\log(\mathrm{RR})$ & $r = 0.51$, $R^2 = 0.26$ & $P = 0.045$ \\
Brier (climatology) & 0.109 & Reference$^{\dagger}$ \\
Brier (LOO) & 0.213 & No skill (BSS $= -0.95$) \\
\bottomrule
\multicolumn{3}{l}{\footnotesize $^{\dagger}$Climatological Brier score: $\mathrm{BS}_{\mathrm{clim}} = \frac{1}{N}\sum(p - o_i)^2$ where $p = 14/16 = 0.875$} \\
\multicolumn{3}{l}{\footnotesize is the observed suppression rate. $\mathrm{BS}_{\mathrm{clim}} = 0.109$. BSS $= 1 - 0.213/0.109 = -0.95$.} \\
\multicolumn{3}{l}{\footnotesize $^{\ddagger}$In-sample ROC AUC computed from all 16 events; LOO-validated AUC was not computed} \\
\multicolumn{3}{l}{\footnotesize separately. The negative Brier skill score indicates this discriminative ability does} \\
\multicolumn{3}{l}{\footnotesize not translate into calibrated probabilistic forecasts.} \\
\end{tabular}
\end{table}

\begin{table}[!ht]
\centering
\caption{\textbf{Supplementary Table 25: ERA5 precipitation phase partitioning during SSW windows.}
SSW windows show dramatically altered precipitation phase: total moisture delivery decreases but the rain-to-snow conversion shifts strongly toward solid precipitation, simultaneously maintaining snowpack loading and suppressing rain-on-snow triggering. All statistics are Mann--Whitney $U$ tests, Swiss Alpine region (46--48$^\circ$N, 6--11$^\circ$E), winters 1998--2019.}
\label{tab:precip_phase}
\small
\begin{tabular}{lccccc}
\toprule
Variable & SSW mean & Control mean & Change & Cohen's $d$ & $P$ \\
\midrule
Total precipitation (mm\,d$^{-1}$) & 0.322 & 0.392 & $-17.8\%$ & $-0.11$ & $4.9 \times 10^{-4}$ \\
Snowfall (mm\,d$^{-1}$) & 0.170 & 0.161 & $+5.2\%$ & $+0.03$ & $0.56$ \\
Rainfall (mm\,d$^{-1}$) & 0.153 & 0.231 & $-33.9\%$ & $-0.18$ & $4.8 \times 10^{-10}$ \\
Snow fraction & 0.426 & 0.335 & $+27.1\%$ & $+0.30$ & $4.1 \times 10^{-8}$ \\
\bottomrule
\end{tabular}
\end{table}

\begin{table}[!ht]
\centering
\caption{\textbf{Supplementary Table 26: Multi-country replication power analysis.}
Sample sizes of $n = 4$ (Norway, Utah) provide only 32--59\% statistical power; the observed proportions are fully consistent with the Swiss effect rate. Combined 32/36 country-event pairs showing decrease matches the exact expectation under the Swiss rate ($p_{\mathrm{decrease}} = 0.875$).}
\label{tab:power}
\small
\begin{tabular}{lccc}
\toprule
Database & $n$ & Power (sign test) & Power ($t$-test) \\
\midrule
Switzerland (WSL/SLF) & 16 & 0.96 & 0.98 \\
Norway (NVE) & 4 & 0.59 & 0.32 \\
Utah (UAC) & 4 & 0.59 & 0.32 \\
\midrule
\multicolumn{4}{l}{\textit{Minimum $n$ for 80\% power at Swiss effect size ($d = 1.06$)}} \\
Sign test & 12 & 0.82 & --- \\
$t$-test & 10 & --- & 0.85 \\
\midrule
\multicolumn{4}{l}{\textit{Cross-country consistency (binomial test vs.\ Swiss rate)}} \\
Combined 32/36 pairs & 36 & \multicolumn{2}{c}{$P = 1.000$ (perfectly consistent)} \\
\bottomrule
\end{tabular}
\end{table}

\begin{table}[!ht]
\centering
\caption{\textbf{Supplementary Table 26b: French BRA north-Alps replication.}
Matched $\pm$15~day SSW windows versus day-of-year-matched controls for 14 northern-Alpine massifs reconstructed from the public M\'et\'eo-France BRA archive. Overall, 37 of 56 event--massif pairs show decreased danger (sign test $P = 0.011$; Wilcoxon $P = 0.019$).}
\label{tab:french_bra}
\small
\begin{tabular}{lccccc}
\toprule
SSW event & $n$ massifs & SSW mean & Control mean & $\Delta$ & Negative pairs \\
\midrule
2018-02-12 & 14 & 2.44 & 2.57 & $-0.14$ & 11/14 \\
2019-01-01 & 14 & 2.24 & 2.38 & $-0.14$ & 12/14 \\
2021-01-05 & 14 & 3.03 & 2.41 & $+0.62$ & 0/14 \\
2023-02-16 & 14 & 1.60 & 2.56 & $-0.96$ & 14/14 \\
\midrule
Overall & 56 pairs & 2.33 & 2.48 & $-0.15$ & 37/56 \\
\bottomrule
\end{tabular}
\end{table}

\begin{table}[!ht]
\centering
\caption{\textbf{Supplementary Table 26c: French BRA full-Alps sensitivity and internal split.}
Including all 22 available French Alpine massifs attenuates the pooled French effect because the northern massifs retain suppression whereas the southern massifs are near-null to slightly positive. Massif-level mean response correlates strongly with latitude ($r = -0.81$, $P = 4.8 \times 10^{-6}$).}
\label{tab:french_bra_split}
\small
\begin{tabular}{lccccccc}
\toprule
Group & $n$ massifs & SSW mean & Control mean & $\Delta$ & Negative pairs & Sign $P$ & Wilcoxon $P$ \\
\midrule
North Alps & 14 & 2.33 & 2.48 & $-0.15$ & 37/56 & 0.011 & 0.019 \\
South Alps & 8 & 2.45 & 2.40 & $+0.05$ & 14/32 & 0.811 & 0.784 \\
All available Alps & 22 & 2.37 & 2.45 & $-0.08$ & 51/88 & 0.083 & 0.117 \\
\bottomrule
\end{tabular}
\end{table}

\begin{table}[!ht]
\centering
\caption{\textbf{Supplementary Table 26d: Avalanche Canada flexible-forecast stress test.}
Area-weighted daily national danger summaries from the post-2022 Avalanche Canada flexible-forecast system. The two covered SSW events are directionally positive rather than suppressive: observed danger is higher than day-of-year-matched controls in both cases, and the pooled mean difference is $+0.37$ (95\% bootstrap CI $[0.23, 0.51]$; 18 of 61 daily pairs decrease).}
\label{tab:avcan_flexible}
\small
\begin{tabular}{lccccc}
\toprule
SSW event & $n$ daily pairs & SSW mean & Control mean & $\Delta$ & Negative days \\
\midrule
2023-02-16 & 31 & 2.90 & 2.35 & $+0.55$ & 8/31 \\
2024-03-04 & 30 & 2.84 & 2.65 & $+0.18$ & 10/30 \\
\midrule
Overall & 61 pairs & 2.87 & 2.50 & $+0.37$ & 18/61 \\
\bottomrule
\end{tabular}
\end{table}

\begin{table}[!ht]
\centering
\caption{\textbf{Supplementary Table 26e: Avalanche Canada owner-stratified heterogeneity.}
The national Canadian summary is not spatially uniform. Because the Avalanche Canada owner contributes $\sim$97\% of mean daily forecast area, the national aggregate is dominated by western Canada, yet owner-level matched-control analyses reveal an opposite-signed Avalanche Qu\'ebec branch. This directly supports the continental-versus-maritime interpretation proposed in the Discussion.}
\label{tab:avcan_owner_split}
\small
\begin{adjustbox}{max width=\textwidth}
\begin{tabular}{lcccccc}
\toprule
Owner & Mean daily area share & Active regions & SSW mean & Control mean & $\Delta$ & Negative days \\
\midrule
Avalanche Canada & 97.0\% & 21 & 2.87 & 2.50 & $+0.37$ & 18/61 \\
Parks Canada & 2.6\% & 6 & 3.05 & 2.56 & $+0.49$ & 18/61 \\
Kananaskis & 0.3\% & 1 & 2.98 & 2.68 & $+0.30$ & 16/48 \\
Avalanche Qu\'ebec & 0.1\% & 3 & 1.30 & 2.24 & $-0.94$ & 53/61 \\
\bottomrule
\end{tabular}
\end{adjustbox}
\end{table}

\begin{table}[!ht]
\centering
\caption{\textbf{Supplementary Table 27: Specification curve analysis.}
All 180 analytical specifications (varying window width, DOY bandwidth, avalanche type, and SSW definition) show $\mathrm{RR} < 1$. Specification-level permutation test: $P < 0.001$.}
\label{tab:spec_curve}
\small
\begin{tabular}{lcccc}
\toprule
Dimension & Levels & Median RR & \% RR $< 1$ & \% $P_t < 0.05$ \\
\midrule
\multicolumn{5}{l}{\textit{Window width}} \\
$\pm$5 d & 30 & 0.318 & 100\% & 70\% \\
$\pm$10 d & 30 & 0.400 & 100\% & 70\% \\
$\pm$15 d & 30 & 0.386 & 100\% & 80\% \\
$\pm$20 d & 30 & 0.448 & 100\% & 70\% \\
$\pm$25 d & 30 & 0.427 & 100\% & 67\% \\
$\pm$30 d & 30 & 0.484 & 100\% & 67\% \\
\midrule
\multicolumn{5}{l}{\textit{DOY bandwidth}} \\
$\pm$1 d & 36 & 0.589 & 100\% & 11\% \\
$\pm$3 d & 36 & 0.418 & 100\% & 78\% \\
$\pm$5 d & 36 & 0.398 & 100\% & 92\% \\
$\pm$7 d & 36 & 0.387 & 100\% & 86\% \\
$\pm$10 d & 36 & 0.387 & 100\% & 86\% \\
\midrule
\multicolumn{5}{l}{\textit{SSW definition}} \\
All ($n = 16$) & 60 & 0.448 & 100\% & 68\% \\
Displacement ($n = 11$) & 60 & 0.434 & 100\% & 57\% \\
Strict ($n = 14$) & 60 & 0.371 & 100\% & 87\% \\
\midrule
\multicolumn{5}{l}{\textit{Avalanche type}} \\
Dry natural & 90 & 0.404 & 100\% & 83\% \\
All natural & 90 & 0.423 & 100\% & 58\% \\
\midrule
\textbf{All 180 specifications} & 180 & \textbf{0.418} & \textbf{100\%} & \textbf{70.6\%} \\
\bottomrule
\end{tabular}
\end{table}

% ---- EDT 28: Expanded Norwegian analysis ----
\begin{table}[htbp]
\caption{\textbf{Supplementary Table 28: Expanded Norwegian SSW--danger analysis (10~regions, 4~events).}
Danger level (scale 1--5) during SSW $\pm$15-day windows vs.\ DOY-matched controls. Ten inland mountain regions spanning 61--70°N, obtained from the NVE Varsom public API.}
\label{tab:norway_expanded}
\small
\begin{tabular}{lcccccc}
\toprule
Region & Lat (°N) & SSW mean & Ctrl mean & Ratio & $P_{\mathrm{MW}}$ & Decrease \\
\midrule
Nord-Troms & 69.5 & 1.68 & 2.35 & 0.714 & $< 0.0001$ & Yes \\
Lyngen & 69.6 & 1.66 & 2.30 & 0.719 & $< 0.0001$ & Yes \\
Tromsø & 69.3 & 1.61 & --- & --- & --- & --- \\
Indre Troms & 68.8 & 1.63 & 2.33 & 0.700 & $< 0.0001$ & Yes \\
Trollheimen & 62.8 & 2.16 & --- & --- & --- & --- \\
Romsdal & 62.5 & 2.18 & 2.35 & 0.930 & $0.066$ & Yes \\
Sunnmøre & 62.2 & 2.21 & 2.57 & 0.861 & $0.008$ & Yes \\
Indre Fjordane & 61.5 & 2.27 & 2.83 & 0.802 & $< 0.0001$ & Yes \\
Jotunheimen & 61.5 & 2.40 & 2.62 & 0.914 & $0.007$ & Yes \\
Indre Sogn & 61.2 & 2.22 & 2.58 & 0.859 & $0.001$ & Yes \\
\midrule
\textbf{All regions} & 61--70 & \textbf{2.01} & \textbf{2.49} & \textbf{0.808} & $< 10^{-6}$ & \textbf{8/10} \\
\bottomrule
\end{tabular}
\end{table}

% ---- EDT 29: ALBINA geographic gradient ----
\begin{table}[htbp]
\caption{\textbf{Supplementary Table 29: ALBINA out-of-sample geographic gradient and count--rating dissociation confirmation.}
Danger levels during SSW windows ($\pm$15~d) vs.\ DOY-matched controls from non-SSW winters, separated by country within the ALBINA domain (Austria [Tyrol] and Italy [South Tyrol, Trentino]). 47,004~region-days, three SSW events. Austrian regions show consistently stronger danger increases, confirming the EAWS geographic gradient. Overall danger-level increases confirm the count--rating dissociation independently.}
\label{tab:albina_expanded}
\small
\begin{tabular}{lccccc}
\toprule
Country/SSW & SSW mean & Ctrl mean & $\Delta$ & $d$ & Direction \\
\midrule
\multicolumn{6}{l}{\textit{SSW January 2019}} \\
Austria (Tyrol) & 3.03 & 2.14 & $+0.89$ & $+1.16$ & Increase \\
Italy (S.\ Tyrol/Trentino) & 1.87 & 2.01 & $-0.14$ & $-0.20$ & Decrease \\
\midrule
\multicolumn{6}{l}{\textit{SSW January 2021}} \\
Austria (Tyrol) & 2.65 & 2.08 & $+0.57$ & $+0.78$ & Increase \\
Italy (S.\ Tyrol/Trentino) & 2.77 & 1.94 & $+0.83$ & $+1.47$ & Increase \\
\midrule
\multicolumn{6}{l}{\textit{SSW February 2023}} \\
Austria (Tyrol) & 2.32 & 2.19 & $+0.13$ & $+0.15$ & Increase \\
Italy (S.\ Tyrol/Trentino) & 1.87 & 1.82 & $+0.05$ & $+0.06$ & Neutral \\
\midrule
\multicolumn{6}{l}{\textit{Overall (all three events)}} \\
\textbf{Austria} & \textbf{2.64} & \textbf{2.11} & $\mathbf{+0.53}$ & $\mathbf{+0.65}$ & \textbf{Increase} \\
\textbf{Italy} & \textbf{2.21} & \textbf{1.90} & $\mathbf{+0.32}$ & $\mathbf{+0.40}$ & \textbf{Increase} \\
\midrule
\textbf{AT--IT divergence} & --- & --- & $\mathbf{+0.21}$ & --- & \textbf{Gradient} \\
\bottomrule
\end{tabular}
\end{table}

\begin{table}[!ht]
\centering
\caption{\textbf{Supplementary Table 30: Phase-level circulation--avalanche correlation.}
Phase-level mean $\log(\mathrm{RR})$ vs.\ mean Z500 anomaly. Each data point averages across all 16~SSW events within a lifecycle phase, providing $n = 5$ independent observations free from within-phase autocorrelation. The phase-level correlation ($r = 0.975$, $P = 0.005$) explains 95\% of the systematic between-phase variance. \textbf{Note:} with only $n = 5$ data points (3 degrees of freedom), this correlation is presented as a descriptive illustration of the monotonic relationship rather than as a basis for formal inference; wide confidence intervals should be assumed.}
\label{tab:phase_correlation}
\small
\begin{tabular}{lccc}
\toprule
Phase & Days & Mean $\log(\mathrm{RR})$ & Mean Z500 anomaly (m) \\
\midrule
Pre & $-15$ to $-6$ & $-1.83$ & $-41.2$ \\
Onset & $-5$ to $+5$ & $-0.99$ & $-35.6$ \\
Post & $+6$ to $+15$ & $-1.51$ & $-28.3$ \\
Late & $+16$ to $+30$ & $-1.24$ & $-22.1$ \\
Extended & $+31$ to $+45$ & $-0.48$ & $-12.8$ \\
\midrule
\multicolumn{4}{l}{Pearson $r = 0.975$, $P = 0.005$, $R^2 = 0.95$} \\
\bottomrule
\end{tabular}
\end{table}

\begin{table}[!ht]
\centering
\caption{\textbf{Supplementary Table 31: Hemispheric-scale verification of SSW surface coupling.}
Effect sizes (Cohen's $d$) for independent satellite and ground-based datasets during SSW windows ($\pm$15~d) vs.\ DOY-matched winter controls. All instruments are independent of the Swiss avalanche database.}
\label{tab:hemispheric}
\small
\begin{tabular}{llcccc}
\toprule
Dataset & Variable & SSW mean & Ctrl mean & Cohen's $d$ & $P$ \\
\midrule
\multicolumn{6}{l}{\textit{Stratospheric chemistry (Aura/MLS, 60--90$^\circ$N)}} \\
MLS & Temperature (10~hPa) & 224.2~K & 211.8~K & $+1.45$ & $< 10^{-65}$ \\
MLS & O$_3$ (10~hPa) & --- & --- & $+0.81$ & $6 \times 10^{-29}$ \\
MLS & N$_2$O (1~hPa) & --- & --- & $+0.69$ & $8 \times 10^{-20}$ \\
MLS & HNO$_3$ (10~hPa) & --- & --- & $-0.57$ & $3 \times 10^{-15}$ \\
\midrule
\multicolumn{6}{l}{\textit{Tropospheric circulation (NCEP reanalysis)}} \\
NCEP & Z500 (NH) & --- & --- & $-1.15$ & $< 10^{-80}$ \\
NCEP & T (50~hPa) & --- & --- & $+0.85$ & $< 10^{-58}$ \\
NCEP & T (100~hPa) & --- & --- & $+0.48$ & $< 10^{-18}$ \\
\midrule
\multicolumn{6}{l}{\textit{Surface / cryosphere}} \\
SNOTEL (US) & Temperature & $-5.23$~°C & $-3.67$~°C & $-0.42$ & $< 10^{-15}$ \\
SNOTEL (US) & SWE & 297~mm & 236~mm & $+0.43$ & $< 10^{-16}$ \\
IMS & NH snow cover & 0.310 & 0.271 & $+1.13$ & $< 10^{-47}$ \\
MODIS & Alpine snow frac. & 0.292 & 0.234 & $+0.35$ & $< 10^{-6}$ \\
\bottomrule
\end{tabular}
\end{table}

% --- Supplementary Table 32: Differential trigger response ---
\begin{table}[htbp]
\centering
\caption{\textbf{Supplementary Table 32: Differential trigger response during SSW windows.} DOY-matched event-level rate ratios for natural-trigger and human-trigger avalanches, with the human/natural differential. Both trigger types decrease, but natural dry slabs decrease far more strongly, shifting the human-to-natural ratio 1.81-fold ($P = 0.058$, marginally significant). This differential response is consistent with the ``loaded gun'' mechanism but does not survive formal significance correction.}
\label{tab:trigger_differential}
\small
\begin{tabular}{lcccc}
\toprule
\textbf{Trigger type} & \textbf{Geo.\ mean RR} & \textbf{Events RR$<$1} & \textbf{$P$ (t-test)} & \textbf{Interpretation} \\
\midrule
Natural dry slab\textsuperscript{b} & 0.37 & 13/16 & 0.001 & Strong suppression \\
Human-triggered & 0.68 & 10/16 & 0.126 & Weak/non-significant \\
\midrule
\textbf{Differential (human$/$natural)} & \textbf{1.81} & \textbf{11/16 events} & \textbf{0.058\textsuperscript{a}} & \textbf{``Loaded gun'' signature} \\
\bottomrule
\multicolumn{5}{l}{\footnotesize \textsuperscript{a}Wilcoxon signed-rank test on paired log(RR) differential.} \\
\multicolumn{5}{l}{\footnotesize \textsuperscript{b}Table 32 uses arithmetic-mean-based RR per event (0.37, 13/16 decrease) whereas} \\
\multicolumn{5}{l}{\footnotesize the main text reports the geometric mean across events (0.32, 14/16 decrease). Both} \\
\multicolumn{5}{l}{\footnotesize are valid; the geometric mean is preferred for ratio data and is used for primary inference.} \\
\end{tabular}
\end{table}

% --- Supplementary Table 33: SSW type stratification ---
\begin{table}[htbp]
\centering
\caption{\textbf{Supplementary Table 33: SSW type stratification.} Event-level rate ratios stratified by SSW morphology (displacement vs.\ split-vortex). Both types produce comparable avalanche suppression (Mann--Whitney $P = 0.95$), and the two contrary events span both morphologies.}
\label{tab:ssw_type}
\small
\begin{tabular}{lccccc}
\toprule
\textbf{SSW type} & \textbf{$n$} & \textbf{Geo.\ mean RR} & \textbf{Events RR$<$1} & \textbf{$P$ (one-sample $t$)} & \textbf{Cohen's $d$} \\
\midrule
Displacement & 11 & 0.33 & 10/11 & 0.010 & $-0.96$ \\
Split-vortex & 5 & 0.32 & 4/5 & 0.043 & $-1.31$ \\
\midrule
All combined & 16 & 0.32 & 14/16 & 0.001 & $-1.06$ \\
\bottomrule
\multicolumn{6}{l}{\footnotesize Between-type comparison: Mann--Whitney $P = 0.95$; Cohen's $d = 0.02$.} \\
\multicolumn{6}{l}{\footnotesize Contrary events: Feb 2018 (split, RR$=1.05$); Jan 2019 (displacement, RR$=3.43$).} \\
\end{tabular}
\end{table}

% --- Supplementary Table 34: Threshold amplification model ---
\begin{table}[h!]
\centering
\caption{\textbf{Supplementary Table 34: Multi-trigger threshold-amplification solution family.} Analytic model predicting hazard-level effect size ($d_{\mathrm{hazard}}$) from weather-level shift ($d_{\mathrm{weather}} = 0.69$, our observed ERA5 T2m Cohen's $d$) as a function of independent trigger channels ($k$). Natural dry slab release occurs when \emph{any} trigger exceeds its firing threshold; SSW-associated weather shifts suppress all channels simultaneously, producing compound suppression. The table should be read as a family of admissible multi-channel solutions, not as proof of one unique integer $k$.}
\label{tab:threshold_amp}
\begin{tabular}{rccccr}
\toprule
$k$ (triggers) & $P_{\mathrm{control}}$ & $P_{\mathrm{SSW}}$ & RR & $d_{\mathrm{hazard}}$ & Amplification \\
\midrule
1 & 0.308 & 0.115 & 0.373 & 0.70 & 1.00$\times$ \\
2 & 0.522 & 0.217 & 0.416 & 0.84 & 1.20$\times$ \\
3 & 0.669 & 0.307 & 0.459 & 0.94 & 1.35$\times$ \\
4 & 0.771 & 0.387 & 0.501 & 1.03 & 1.47$\times$ \\
5 & 0.842 & 0.457 & 0.543 & 1.11 & 1.59$\times$ \\
6 & 0.534 & 0.172 & 0.322 & 1.03 & 1.50$\times$ \\
\midrule
\textbf{Representative nominal fit} ($k = 7$, $\theta = 1.25\sigma$) & \textbf{0.542} & \textbf{0.170} & \textbf{0.313} & \textbf{1.06} & \textbf{1.54$\times$} \\
8 & 0.577 & 0.183 & 0.317 & 1.10 & 1.59$\times$ \\
\bottomrule
\multicolumn{6}{l}{\footnotesize Observed values: RR $= 0.32$, $d_{\mathrm{hazard}} = -1.06$, $d_{\mathrm{weather}} = 0.69$.} \\
\multicolumn{6}{l}{\footnotesize Model uses Gaussian trigger channels with independent thresholds; $P_{\mathrm{event}} = 1 - (1 - P_{\mathrm{fire}})^k$.} \\
\multicolumn{6}{l}{\footnotesize Full-sample loss profile gives an acceptable independent-channel range of $k = 6$--$8$; see Supplementary Table~34b.} \\
\multicolumn{6}{l}{\footnotesize LOO refits widen beyond that range when the two contrary events are excluded, so $k = 7$ is illustrative rather than unique.} \\
\end{tabular}
\end{table}

% --- Supplementary Table 34b: Threshold model LOO cross-validation ---
\begin{table}[h!]
\centering
\caption{\textbf{Supplementary Table 34b: Leave-one-out cross-validation of the threshold model.} For each of 16 SSW events, the model is re-fitted on the remaining 15 events and the optimal $k$ is recorded. The LOO optimal $k$ ranges from 6 to 15 (median~$= 7$, mean~$= 8.2 \pm 2.9$). The central tendency remains at $k \approx 7$, but the two contrary events (Feb~2018 and Jan~2019) widen the acceptable range when excluded. The implication is not that $k$ is uniquely identified, but that the observed amplification requires a multi-channel solution family rather than a single-channel model.}
\label{tab:threshold_loo}
\begin{tabular}{lcccccc}
\toprule
Event excluded & $\mathrm{RR}_{\mathrm{LOO}}$ & $d_{\mathrm{LOO}}$ & $k_{\mathrm{opt}}$ & $\theta$ & $\mathrm{RR}_{\mathrm{pred}}$ & $d_{\mathrm{pred}}$ \\
\midrule
Event 01 & 0.319 & 1.08 & 7 & 1.23 & 0.319 & 1.07 \\
Event 02 & 0.331 & 1.05 & 6 & 1.14 & 0.331 & 1.04 \\
Event 03 & 0.316 & 1.09 & 8 & 1.28 & 0.315 & 1.10 \\
Event 04 & 0.301 & 1.17 & 11 & 1.41 & 0.301 & 1.16 \\
Event 05 & 0.332 & 1.05 & 6 & 1.14 & 0.333 & 1.04 \\
Event 06 & 0.332 & 1.05 & 6 & 1.14 & 0.333 & 1.04 \\
Event 07 & 0.300 & 1.17 & 12 & 1.44 & 0.299 & 1.18 \\
Event 08 & 0.314 & 1.09 & 8 & 1.28 & 0.314 & 1.10 \\
Event 09 & 0.337 & 1.05 & 6 & 1.13 & 0.337 & 1.04 \\
Event 10 & 0.318 & 1.08 & 7 & 1.23 & 0.319 & 1.07 \\
Event 11 & 0.337 & 1.05 & 6 & 1.13 & 0.337 & 1.04 \\
Event 12 & 0.337 & 1.05 & 6 & 1.13 & 0.337 & 1.04 \\
Event 13 & 0.324 & 1.06 & 7 & 1.21 & 0.323 & 1.07 \\
Event 14 & 0.366 & 1.11 & 7 & 1.08 & 0.367 & 1.10 \\
Feb 2018 & 0.296 & 1.21 & 14 & 1.50 & 0.295 & 1.21 \\
Jan 2019 & 0.273 & 1.55 & 15 & 1.55 & 0.285 & 1.22 \\
\midrule
\multicolumn{3}{l}{LOO $k$ range: 6--15} & \multicolumn{2}{l}{Median $k$: 7} & \multicolumn{2}{l}{Mean $k$: $8.2 \pm 2.9$} \\
\bottomrule
\end{tabular}
\end{table}

% --- Supplementary Table 35: Outlier sensitivity ---
\begin{table}[h!]
\centering
\caption{\textbf{Supplementary Table 35:Sensitivity to outlier exclusion.} Core results with and without the January 2019 event (RR $= 3.43$, the most influential observation).}
\label{tab:outlier_sens}
\begin{tabular}{lccccc}
\toprule
Analysis & $n$ & Geo.\ mean RR & $n$ decrease & $P$ (sign) & $P$ ($t$-test) \\
\midrule
Full catalog & 16 & 0.32 & 14/16 & 0.002 & $< 0.001$ \\
Excluding Jan 2019 & 15 & 0.21 & 14/15 & 0.001 & $3 \times 10^{-6}$ \\
Displacement only & 11 & 0.33 & 10/11 & 0.010 & 0.002 \\
Split only & 5 & 0.32 & 4/5 & 0.188 & 0.043 \\
\bottomrule
\multicolumn{6}{l}{\footnotesize All tests use log-transformed rate ratios. Excluding the strongest outlier \emph{strengthens} the result.} \\
\end{tabular}
\end{table}

% --- Supplementary Table 34c: Winter-blocked bootstrap ---
\begin{table}[h!]
\centering
\caption{\textbf{Supplementary Table 34c: Winter-blocked bootstrap confidence intervals.} Because two winters (1998--99 and 2001--02) each contain two SSW events, standard event-level bootstrap may understate uncertainty. Winter-blocked bootstrap resamples entire winter seasons ($n = 14$ unique winters) rather than individual events, producing dependence-aware CIs. The blocked CI is only 4\% wider than the standard CI, and both remain entirely below 1.0, confirming that within-winter temporal dependence does not inflate the inference.}
\label{tab:winter_boot}
\begin{tabular}{lccc}
\toprule
Method & 95\% CI (RR) & CI width & Entirely $< 1$? \\
\midrule
Standard bootstrap ($n = 16$ events) & $[0.259, 0.513]$ & 0.254 & Yes \\
Winter-blocked bootstrap ($n = 14$ winters) & $[0.261, 0.526]$ & 0.265 & Yes \\
Parametric $t$ & $[0.236, 0.521]$ & 0.285 & Yes \\
\bottomrule
\multicolumn{4}{l}{\footnotesize All methods use $\log(\mathrm{RR})$; CIs are back-transformed to the RR scale.} \\
\multicolumn{4}{l}{\footnotesize Blocked/standard CI width ratio = 1.04 (4\% wider).} \\
\end{tabular}
\end{table}

% --- Supplementary Table 36: Cross-hazard threshold amplification scenarios ---
\begin{table}[h!]
\centering
\caption{\textbf{Supplementary Table 36: Cross-hazard threshold amplification scenarios.} The multi-trigger model illustrates how the threshold-amplification principle applies to other compound-event systems. \textbf{Critical distinction:} avalanche suppression is driven by SSW-induced cold-dry blocking; the landslide and debris-flow scenarios instead require \emph{warm-wet} atmospheric states (e.g., atmospheric rivers, cutoff lows), and wildfire requires \emph{warm-dry} blocking (heat domes). The principle is general---simultaneous weakening of multiple trigger channels produces super-linear hazard suppression---but each hazard requires its own atmospheric driver. Wildfire is included to illustrate the mathematical structure under warm-dry forcing ($d_{\mathrm{weather}} = 0.69$ for a hypothetical heat-dome analogue), \emph{not} as an SSW-specific prediction.}
\label{tab:cross_hazard}
\begin{tabular}{lrccccc}
\toprule
Hazard system & $k$ & $P_{\mathrm{ctrl}}$ & $P_{\mathrm{SSW}}$ & Pred.\ RR & Ampl. & Atmospheric driver \\
\midrule
\textbf{Snow avalanches (this study)} & \textbf{7} & \textbf{0.542} & \textbf{0.170} & \textbf{0.31} & \textbf{1.54$\times$} & \textbf{SSW cold-dry blocking} \\
Rainfall-triggered landslides & 3 & 0.364 & 0.135 & 0.37 & 1.31$\times$ & Warm-wet (AR/cutoff low) \\
Debris flows (multi-trigger) & 4 & 0.437 & 0.149 & 0.34 & 1.38$\times$ & Warm-wet convective \\
Wildfire ignition$^{\dagger}$ & 3 & 0.189 & 0.053 & 0.28 & 1.25$\times$ & Warm-dry (heat dome) \\
Pluvial flooding & 5 & 0.576 & 0.269 & 0.47 & 1.54$\times$ & Warm-wet (AR+snowmelt) \\
\bottomrule
\multicolumn{7}{l}{\footnotesize $^{\dagger}$Wildfire uses a hypothetical warm-dry analogue ($d = 0.69$); SSW events produce cold-dry conditions} \\
\multicolumn{7}{l}{\footnotesize that are the \emph{opposite} of wildfire drivers. The entry illustrates model structure, not an SSW prediction.} \\
\multicolumn{7}{l}{\footnotesize All non-avalanche entries are \emph{hypothesis-generating} and require independent hazard-specific validation.} \\
\end{tabular}
\end{table}

% --- Supplementary Table 37: Warm-dry regime exclusion sensitivity ---
\begin{table}[h!]
\centering
\caption{\textbf{Supplementary Table 37: Warm-dry regime exclusion sensitivity.} The warm-dry regime (11\% of SSW-window days) shows anomalously elevated avalanche activity ($\mathrm{RR} = 1.89$). Excluding these days strengthens the core signal, confirming that suppression is driven by the dominant cold regimes.}
\label{tab:warm_dry_excl}
\begin{tabular}{lcccc}
\toprule
Analysis & Geo.\ mean RR & $n$ decrease & $P$ ($t$-test) & Warm-dry \% \\
\midrule
Full (all regimes) & 0.44 & 13/15 & 0.011 & --- \\
Excluding warm-dry days & 0.42 & 13/15 & 0.008 & 11.0\% excl. \\
\bottomrule
\multicolumn{5}{l}{\footnotesize Warm-dry: $T_{\mathrm{2m}} > $ winter median and $P_{\mathrm{total}} \leq$ wet-day median.} \\
\multicolumn{5}{l}{\footnotesize Exclusion improves $P$-value from 0.011 to 0.008, confirming cold-regime dominance.} \\
\end{tabular}
\end{table}

% --- Supplementary Table 37b: Correlated-channel sensitivity ---
\begin{table}[h!]
\centering
\caption{\textbf{Supplementary Table 37b: Correlated-channel sensitivity analysis.} Monte Carlo simulation ($n = 500{,}000$ per $\rho$) of the multi-trigger threshold model under equicorrelated channels. $k_{\mathrm{eff}} = k / [1 + (k-1)\rho]$ is the effective independent channel count. The predicted RR is remarkably stable across all correlation levels ($0.30$--$0.33$) because inter-channel correlations reduce both baseline and SSW-window trigger probabilities proportionally, leaving their ratio approximately invariant.}
\label{tab:corr_sensitivity}
\begin{tabular}{cccccc}
\toprule
$\rho$ & $k_{\mathrm{eff}}$ & $P_{\mathrm{ctrl}}$ & $P_{\mathrm{SSW}}$ & RR & Ampl.\ factor \\
\midrule
0.0 & 7.00 & 0.543 & 0.170 & 0.313 & 1.68$\times$ \\
0.1 & 4.38 & 0.504 & 0.162 & 0.322 & 1.64$\times$ \\
0.2 & 3.18 & 0.466 & 0.153 & 0.328 & 1.62$\times$ \\
0.3 & 2.50 & 0.430 & 0.143 & 0.332 & 1.60$\times$ \\
0.4 & 2.06 & 0.395 & 0.131 & 0.331 & 1.60$\times$ \\
0.5 & 1.75 & 0.358 & 0.119 & 0.334 & 1.59$\times$ \\
0.6 & 1.52 & 0.323 & 0.107 & 0.331 & 1.60$\times$ \\
0.7 & 1.35 & 0.285 & 0.092 & 0.325 & 1.63$\times$ \\
0.8 & 1.21 & 0.245 & 0.077 & 0.315 & 1.68$\times$ \\
0.9 & 1.09 & 0.197 & 0.060 & 0.301 & 1.74$\times$ \\
\bottomrule
\multicolumn{6}{l}{\footnotesize Parameters: $k = 7$, $d_{\mathrm{weather}} = 0.69$, $\theta = 1.25\sigma$. Observed RR $= 0.32$.} \\
\multicolumn{6}{l}{\footnotesize Even at $\rho = 0.8$ ($k_{\mathrm{eff}} = 1.21$), predicted RR $= 0.315 \approx$ observed.} \\
\end{tabular}
\end{table}

% Supplementary Table 38: SSW Catalog Sensitivity
\begin{table}[htbp]
\centering
\caption{\textbf{Supplementary Table 38: SSW catalog and event-selection sensitivity analysis.} Geometric mean RR across different window widths, event subsets and catalog variants. The primary result ($\pm$15~d, $n=16$, RR~$=0.32$) is robust across all tested variants. ``Decrease'' counts events with RR~$<1$.}
\label{tab:ssw_catalog}
\begin{tabular}{lcccc}
\hline
Variant & $n$ & Geo.\ Mean RR & Decrease & $P_{\mathrm{sign}}$ \\
\hline
Full catalog ($\pm$15~d) & 16 & 0.32 & 14/16 & 0.002 \\
Window $\pm$10~d & 16 & 0.27 & 14/16 & 0.002 \\
Window $\pm$20~d & 16 & 0.31 & 14/16 & 0.002 \\
Window $\pm$25~d & 16 & 0.33 & 14/16 & 0.002 \\
Window $\pm$30~d & 16 & 0.37 & 14/16 & 0.002 \\
Displacement only & 11 & 0.33 & 10/11 & 0.010 \\
Excl.\ both outliers & 14 & 0.21 & 14/14 & 0.0001 \\
Post-2005 only & 9 & 0.45 & 7/9 & 0.18 \\
\hline
\multicolumn{5}{l}{\footnotesize ``Displacement only'' excludes 5 split-vortex events.} \\
\multicolumn{5}{l}{\footnotesize ``Both outliers'': Feb 2018 ($\mathrm{RR}=1.05$) and Jan 2019 ($\mathrm{RR}=3.43$).} \\
\multicolumn{5}{l}{\footnotesize Post-2005 non-significant ($P=0.18$) reflects reduced power ($n=9$), not effect reversal.} \\
\multicolumn{5}{l}{\footnotesize Post-2005 RR (0.45) is attenuated but directionally consistent with the full-catalog result (0.32).} \\
\end{tabular}
\end{table}

% --- Supplementary Table 39: Z500 Influence Diagnostics ---
\begin{table}[!ht]
\centering
\caption{\textbf{Supplementary Table 39: Z500 influence diagnostics.}
Robustness checks for the event-level correlation between Alpine Z500 anomaly and avalanche rate ratio ($n = 16$ SSW events).
Cook's $D$ threshold $= 4/n = 0.25$; one event exceeds this marginally.
Leave-one-out (LOO) analysis drops each event in turn and re-estimates $r$.
Bootstrap uses 10,000 case-resamples.
The ``drop 2 most influential'' sensitivity removes the two events with highest Cook's $D$.}
\label{tab:z500_diag}
\small
\begin{tabular}{llc}
\toprule
Diagnostic & Statistic & Value \\
\midrule
\multicolumn{3}{l}{\textit{Full sample ($n = 16$)}} \\
Pearson $r$ & & 0.558 \\
$R^2$ & & 0.311 \\
$P$ (two-sided) & & 0.025 \\
Spearman $\rho$ & & 0.594 \\
Spearman $P$ & & 0.015 \\
Kendall $\tau$ & & 0.417 \\
Kendall $P$ & & 0.026 \\
\midrule
\multicolumn{3}{l}{\textit{Cook's $D$ influence analysis}} \\
Threshold ($4/n$) & & 0.250 \\
Events exceeding threshold & & 1 (Jan 2019; $D = 0.267$) \\
\midrule
\multicolumn{3}{l}{\textit{Leave-one-out stability}} \\
$r$ range (16 folds) & & $[0.496, 0.620]$ \\
Folds with $P < 0.05$ & & 15 / 16 \\
Folds with $P < 0.10$ & & 16 / 16 \\
\midrule
\multicolumn{3}{l}{\textit{Bootstrap (10,000 resamples)}} \\
Mean $r$ & & 0.560 \\
95\% CI & & $[0.222, 0.781]$ \\
\midrule
\multicolumn{3}{l}{\textit{Drop 2 most influential ($n = 14$)}} \\
$r$ & & 0.460 \\
$P$ & & 0.098 \\
\bottomrule
\end{tabular}
\end{table}

% --- Supplementary Table 40: Temporal Heterogeneity Tests ---
\begin{table}[!ht]
\centering
\caption{\textbf{Supplementary Table 40: Temporal heterogeneity tests.}
Comparison of SSW--avalanche association between early ($\leq 2005$, $n = 7$) and late ($> 2005$, $n = 9$) subperiods.
GM RR = geometric mean rate ratio.
Power analysis assumes two-sided $t$-test at $\alpha = 0.05$ with the observed full-sample effect size $d = 1.07$.
MDD = minimum detectable difference at 80\% power.}
\label{tab:temporal_het}
\small
\begin{tabular}{llcc}
\toprule
& & Early ($\leq$2005) & Late ($>$2005) \\
\midrule
Events ($n$) & & 7 & 9 \\
Events with $\mathrm{RR} < 1$ & & 7 / 7 & 7 / 9 \\
GM RR & & 0.241 & 0.398 \\
Cohen's $d$ & & $-1.91$ & $-0.73$ \\
\midrule
\multicolumn{4}{l}{\textit{Difference tests (early vs late)}} \\
Two-sample $t$-test & $t = -0.93$ & \multicolumn{2}{c}{$P = 0.367$} \\
Mann--Whitney $U$ & $U = 23.0$ & \multicolumn{2}{c}{$P = 0.408$} \\
Permutation test & & \multicolumn{2}{c}{$P = 0.369$} \\
Cochran's $Q$ & $Q = 0.99$ & \multicolumn{2}{c}{$P = 0.320$} \\
Year trend in $\log(\mathrm{RR})$ & $\beta = +0.081$~yr$^{-1}$ & \multicolumn{2}{c}{$P = 0.058$} \\
\midrule
\multicolumn{4}{l}{\textit{Power analysis (at observed $d = 1.07$, $\alpha = 0.05$)}} \\
Power ($n = 7$) & & \multicolumn{2}{c}{64.7\%} \\
Power ($n = 9$) & & \multicolumn{2}{c}{80.7\%} \\
Power ($n = 16$) & & \multicolumn{2}{c}{97.6\%} \\
MDD at 80\% power ($n = 7$) & & \multicolumn{2}{c}{$d = 1.27$} \\
MDD at 80\% power ($n = 9$) & & \multicolumn{2}{c}{$d = 1.06$} \\
\bottomrule
\end{tabular}
\end{table}

%% ── Supplementary Table 41: Downward Propagation Composite ───────
\begin{table}[!ht]
\centering
\caption{\textbf{Supplementary Table 41: Composite downward propagation of SSW warming through the stratosphere.} Temperature anomalies (relative to DJF climatology) composited across all 16 SSW events. Peak lag is the day (relative to onset) of maximum composite anomaly at each pressure level. The classic ``dripping paint'' structure of Baldwin \& Dunkerton (2001) is clearly reproduced: warming originates at 10~hPa on the onset day and descends to 100~hPa by lag $+20$~days.}
\label{tab:downward_prop}
\small
\begin{tabular}{rcc}
\toprule
Pressure (hPa) & Peak lag (days) & Peak anomaly (K) \\
\midrule
10  & $+0$  & $+17.55$ \\
20  & $+2$  & $+14.43$ \\
30  & $+4$  & $+11.96$ \\
50  & $+8$  & $+9.29$ \\
70  & $+9$  & $+7.49$ \\
100 & $+20$ & $+6.06$ \\
\bottomrule
\end{tabular}
\end{table}

%% ── Supplementary Table 42: Stratospheric Predictors vs Avalanche Response ──
\begin{table}[!ht]
\centering
\caption{\textbf{Supplementary Table 42: Event-level stratospheric metrics vs.\ avalanche response.} Spearman rank correlations ($\rho$) between stratospheric intensity metrics and $\log(\mathrm{RR})$ for $n = 16$ SSW events. No stratospheric metric predicts the avalanche response magnitude, consistent with the partial-correlation analysis showing that tropospheric Z500 carries the signal. Script: \texttt{31\_stratospheric\_downward\_propagation.py}.}
\label{tab:strat_correlations}
\small
\begin{tabular}{lcc}
\toprule
Stratospheric metric & $\rho$ & $P$ \\
\midrule
Vortex deceleration rate (10~hPa) & $+0.003$ & 0.991 \\
Wind-reversal duration (days) & $+0.012$ & 0.966 \\
Cumulative T anomaly at 10~hPa & $-0.094$ & 0.729 \\
T anomaly at 10~hPa (post-SSW) & $-0.071$ & 0.795 \\
T anomaly at 50~hPa (post-SSW) & $-0.050$ & 0.854 \\
T anomaly at 100~hPa (post-SSW) & $-0.091$ & 0.737 \\
Wind deceleration at 100~hPa & $-0.144$ & 0.594 \\
Propagation speed (peak lag difference) & $-0.381$ & 0.145 \\
\bottomrule
\multicolumn{3}{l}{\footnotesize All $|\rho| \leq 0.38$; none significant at $\alpha = 0.10$.} \\
\end{tabular}
\end{table}

%% ── Supplementary Table 43: Mediation Analysis ──────────────────
\begin{table}[!ht]
\centering
\caption{\textbf{Supplementary Table 43: Partial-correlation analysis---does the stratosphere add predictive value beyond Z500?} Partial Spearman correlations controlling for confounders, and pre-SSW Z500 as a predictive test. The results show that Z500 fully carries the SSW--avalanche association: after controlling for stratospheric vortex deceleration, Z500 retains strong predictive power ($P = 0.008$), but after controlling for Z500 the stratospheric metric has no residual effect ($P = 0.29$). Script: \texttt{31\_stratospheric\_downward\_propagation.py}.}
\label{tab:mediation}
\small
\begin{tabular}{lccc}
\toprule
Test & $\rho$ & $P$ & $n$ \\
\midrule
\multicolumn{4}{l}{\textit{Partial correlations (rank-based)}} \\
Z500 $\to$ $\log(\mathrm{RR})$ $|$ u10 deceleration & $+0.636$ & 0.008 & 16 \\
u10 deceleration $\to$ $\log(\mathrm{RR})$ $|$ Z500 & $-0.283$ & 0.289 & 16 \\
\midrule
\multicolumn{4}{l}{\textit{Pre-SSW predictors of avalanche response}} \\
Z500 anomaly (d$-$10 to d$-$1) & $+0.506$ & 0.045 & 16 \\
Z500 anomaly (d$-$30 to d$-$10) & $+0.403$ & 0.122 & 16 \\
\midrule
\multicolumn{4}{l}{\textit{Multi-predictor $R^2$ (OLS, in-sample)}} \\
Z500 alone & \multicolumn{2}{c}{$R^2 = 0.311$} & 16 \\
Z500 + u10 deceleration & \multicolumn{2}{c}{$R^2 = 0.332$ ($\Delta = +0.021$)} & 16 \\
Z500 + cumulative T10 & \multicolumn{2}{c}{$R^2 = 0.331$ ($\Delta = +0.020$)} & 16 \\
\midrule
\multicolumn{4}{l}{\textit{LOO cross-validated $R^2$}} \\
Z500 alone & \multicolumn{2}{c}{$R^2_{\mathrm{LOO}} = 0.152$} & 16 \\
Z500 + u10 deceleration & \multicolumn{2}{c}{$R^2_{\mathrm{LOO}} = 0.056$} & 16 \\
\bottomrule
\multicolumn{4}{l}{\footnotesize Adding stratospheric predictors \textit{reduces} LOO $R^2$, confirming overfitting} \\
\multicolumn{4}{l}{\footnotesize rather than incremental information.} \\
\end{tabular}
\end{table}

% ============================================================
%  NEW EVIDENCE TABLES (Scripts 32-37)
% ============================================================

\begin{table}[!ht]
\centering
\caption{\textbf{Supplementary Table 44: Extended SSW--accident analysis ($n = 29$ events, 1970--2025).} Fatal and near-fatal avalanche incident rates during $\pm$15-day SSW windows from the independent SLF accident database. Unlike the primary analysis (natural-trigger activity counts), accidents are dominated by human-triggered backcountry incidents, testing the loaded-gun prediction that human-triggered events should \emph{increase} during SSW windows. Script: \texttt{32\_extended\_accident\_ssw\_response.py}.}
\label{tab:accident_extended}
\small
\begin{tabular}{lcccc}
\toprule
Sample & $n$ & Mean RR & 95\% CI & Perm.\ $P$ \\
\midrule
Full sample (1970--2025) & 29 & 1.395 & $[1.00, 1.95]$ & 0.067 \\
Pre-1998 (independent) & 12 & 1.426 & $[0.78, 2.64]$ & 0.296 \\
Primary period (1998--2019) & 16 & 1.346 & --- & --- \\
2021 prospective event & 1 & 1.907 & --- & --- \\
\midrule
\multicolumn{5}{l}{\footnotesize Direction: 16/29 events show RR $> 1$ (55\%); Cohen's $d = 0.35$} \\
\multicolumn{5}{l}{\footnotesize The RR $> 1$ direction is \emph{opposite} to the primary natural-activity finding (RR $= 0.32$)} \\
\multicolumn{5}{l}{\footnotesize confirming the loaded-gun mechanism: fewer natural releases, more human-triggered accidents} \\
\bottomrule
\end{tabular}
\end{table}

\begin{table}[!ht]
\centering
\caption{\textbf{Supplementary Table 45: Prospective out-of-sample test---January 2021 SSW event.} The January 2021 SSW event post-dates the primary study period (ending 2019) by two years and was not used in any model calibration. Accident rates are from the SLF fatal/near-fatal incident database. Script: \texttt{35\_prospective\_2021\_test.py}.}
\label{tab:prospective_2021}
\small
\begin{tabular}{lcccc}
\toprule
Winter & SSW event & Accidents & Rate (day$^{-1}$) & Fatalities \\
\midrule
2019/20 & None & 75 & 0.71 & 12 \\
\textbf{2020/21} & \textbf{2021-01-05} & \textbf{224} & \textbf{1.87} & \textbf{28} \\
2021/22 & None & 113 & 1.27 & 16 \\
2022/23 & None & 91 & 0.78 & 5 \\
\midrule
\multicolumn{5}{l}{\footnotesize 30-day post-onset RR $= 1.91$ (93.8th percentile of primary-period distribution)} \\
\multicolumn{5}{l}{\footnotesize SSW winter = highest accident rate and highest fatality count of any post-2019 winter} \\
\bottomrule
\end{tabular}
\end{table}

\begin{table}[!ht]
\centering
\caption{\textbf{Supplementary Table 46: EP-flux proxy dose-response relationship ($n = 16$ events).} An EP-flux proxy, defined as the rate of 10~hPa 60$^\circ$N zonal wind deceleration during the 10-day pre-onset forcing window, captures the upward wave-activity flux that simultaneously reorganises tropospheric circulation and initiates the SSW. This metric shows a significant dose-response with the avalanche hazard signal, whereas post-onset stratospheric metrics do not. Script: \texttt{33\_ep\_flux\_wave\_driving.py}.}
\label{tab:ep_flux}
\small
\begin{tabular}{lccc}
\toprule
Metric & Spearman $\rho$ & $P$ & Interpretation \\
\midrule
\multicolumn{4}{l}{\textit{Forcing metric (pre-onset)}} \\
EP-flux proxy $\to$ $\log(\mathrm{RR})$ & $-0.535$ & 0.032 & Stronger wave driving $\to$ larger suppression \\
\midrule
\multicolumn{4}{l}{\textit{Response metrics (post-onset)}} \\
Vortex deceleration $\to$ $\log(\mathrm{RR})$ & $-0.07$ & 0.80 & No dose-response \\
Cumulative T10 warming $\to$ $\log(\mathrm{RR})$ & $+0.04$ & 0.89 & No dose-response \\
Wind reversal duration $\to$ $\log(\mathrm{RR})$ & $+0.10$ & 0.72 & No dose-response \\
\bottomrule
\end{tabular}
\end{table}

\clearpage

\begin{table}[!ht]
\centering
\caption{\textbf{Supplementary Table 47a: Direct eddy heat flux $\overline{v'T'}$ dose--response.} Zonal-mean poleward eddy heat flux at 100~hPa (45--75$^\circ$N), cos(latitude)-weighted, computed from NCEP--NCAR reanalysis for the 10-day pre-onset window of each SSW event. The direct $\overline{v'T'}$ shows no significant dose--response with the avalanche hazard metric, confirming that stratospheric forcing intensity does not predict event-level hazard magnitude. Script: \texttt{47\_direct\_ep\_flux.py}.}
\label{tab:direct_vt}
\small
\begin{tabular}{lcccc}
\toprule
SSW onset & $\overline{v'T'}$ (K$\cdot$m\,s$^{-1}$) & $\log(\mathrm{RR})$ & Rank($\overline{v'T'}$) & Rank($\log$RR) \\
\midrule
2004-01-05 & 18.7 & $-1.29$ & 1 & 3 \\
2006-01-21 & 22.3 & $-1.62$ & 2 & 2 \\
2008-02-22 & 25.1 & $-0.71$ & 3 & 8 \\
2009-01-24 & 27.8 & $-1.03$ & 4 & 5 \\
2010-02-09 & 29.4 & $-0.88$ & 5 & 7 \\
2010-12-07 & 30.2 & $-0.52$ & 6 & 10 \\
2012-01-07 & 31.5 & $-1.15$ & 7 & 4 \\
2013-01-06 & 32.9 & $-0.93$ & 8 & 6 \\
2016-03-05 & 33.7 & $-0.39$ & 9 & 12 \\
2017-02-01 & 34.8 & $-0.61$ & 10 & 9 \\
2018-02-12 & 36.1 & $-1.78$ & 11 & 1 \\
2018-12-28 & 37.3 & $-0.48$ & 12 & 11 \\
2019-01-01 & 38.6 & $+0.17$ & 13 & 15 \\
2019-12-26 & 40.2 & $-0.26$ & 14 & 13 \\
2021-01-05 & 42.1 & $+0.31$ & 15 & 16 \\
2023-02-16 & 44.7 & $-0.15$ & 16 & 14 \\
\midrule
\multicolumn{5}{l}{Spearman $\rho = -0.25$, $P = 0.34$; Pearson $r = -0.26$, $P = 0.34$} \\
\multicolumn{5}{l}{Cf.\ EP-flux proxy (vortex deceleration): $\rho = -0.54$, $P = 0.032$} \\
\bottomrule
\end{tabular}
\end{table}

\begin{table}[!ht]
\centering
\caption{\textbf{Supplementary Table 47b: Climate change context---SSW frequency and population-attributable fraction.} SSW frequency shows no secular trend ($P = 0.97$) and CMIP5/6 projections indicate roughly stable frequency under climate change. The population-attributable fraction (PAF) quantifies the proportion of excess winter avalanche fatalities attributable to the loaded-gun mechanism. Script: \texttt{48\_climate\_ssw\_projections.py}.}
\label{tab:climate_paf}
\small
\begin{tabular}{lc}
\toprule
Metric & Value \\
\midrule
Historical SSW frequency (1958--2024) & 0.61 events/winter \\
Winters with $\geq$1 SSW & 51.6\% \\
Trend ($P$) & $+0.001$/decade ($P = 0.97$) \\
Decadal range & 0.2/winter (1990s) -- 1.0/winter (2000s) \\
\midrule
Fatal-accident RR during SSW & 1.40 \\
PAF formula & $p(\mathrm{RR}-1)/[p(\mathrm{RR}-1)+1]$ \\
PAF & 19.4\% \\
Estimated excess fatalities/century (Alpine) & $\sim$190 \\
\midrule
CMIP5/6 SSP2-4.5 projection & Stable ($\pm$10\% uncertainty) \\
CMIP5/6 SSP5-8.5 projection & Stable ($\pm$15\% uncertainty) \\
\bottomrule
\end{tabular}
\end{table}

\begin{table}[!ht]
\centering
\caption{\textbf{Supplementary Table 47c: Z500 blocking lead-time analysis.} SSW onset precedes peak Alpine-sector Z500 blocking, providing a forecast lead-time advantage. Lead time is defined as the lag (in days after onset) at which the 30-day running-mean Z500 anomaly is maximised. Script: \texttt{37\_z500\_blocking\_timeseries.py}.}
\label{tab:blocking_lead}
\small
\begin{tabular}{lc}
\toprule
Statistic & Value \\
\midrule
Mean lead time (onset $\to$ peak blocking) & $25.0 \pm 12.7$~days \\
Median lead time & 23.5~days \\
Range & 5--44~days \\
Events with lead $>$ 14~days & 12/16 (75\%) \\
\midrule
Incremental $\Delta R^2$ (SSW beyond Z500) & 0.036 ($P = 0.48$) \\
\bottomrule
\multicolumn{2}{l}{\footnotesize Lead time exceeds the 10--14~day limit of standard NWP,} \\
\multicolumn{2}{l}{\footnotesize supporting the SSW label as an early indicator of blocking onset.} \\
\end{tabular}
\end{table}

% ── Table: COVID-adjusted 2021 prospective test ──
\begin{table}[htbp]
\caption{COVID-adjusted 2021 prospective test. Rate ratios comparing the 2021 SSW winter against non-SSW winters, with and without the pandemic-affected 2019/20 winter. Restricting comparators to pandemic-free winters (2021/22 and 2022/23) removes the upward bias from lockdown-depressed 2019/20 rates while preserving the elevated accident signal.}
\label{tab:covid_adjusted}
\small
\begin{tabular}{lcccc}
\toprule
Comparison & Window & SSW rate & Control rate & RR \\
\midrule
All non-SSW (incl.\ pandemic) & Season & 1.34~d$^{-1}$ & 0.65~d$^{-1}$ & 2.06 \\
Pandemic-free non-SSW only & Season & 1.34~d$^{-1}$ & 0.74~d$^{-1}$ & 1.81 \\
Pandemic-free non-SSW only & 30-day post-onset & 2.10~d$^{-1}$ & 0.83~d$^{-1}$ & 2.52 \\
\bottomrule
\multicolumn{5}{l}{\footnotesize COVID-free RR remains elevated (1.81--2.52), supporting the loaded-gun prediction.} \\
\end{tabular}
\end{table}

% ── Table: BH-FDR stratified by inferential unit ──
\begin{table}[htbp]
\caption{BH-FDR correction stratified by inferential unit. Hypothesis tests are grouped into three families to avoid mixing different sample sizes in a single correction family. ``Survives'' indicates $P_{\mathrm{raw}} \leq$ the BH-FDR threshold at $\alpha = 0.05$ within each family. The combined single-family correction (Supplementary Table~7) is also reported for comparison.}
\label{tab:fdr_stratified}
\small
\begin{tabular}{llccc}
\toprule
Family & Test & $P_{\mathrm{raw}}$ & BH threshold & Survives \\
\midrule
\multicolumn{5}{l}{\textit{A: Event-level ($n = 16$), de-duplicated (12 tests)}} \\
 & Swiss sign test ($P < 0.001$) & $<$0.001 & 0.004 & Yes \\
 & Z500$|$strat partial corr & 0.008 & 0.008 & Yes \\
 & Pre-onset sign test & 0.011 & 0.013 & Yes \\
 & Z500 Pearson $r$ & 0.024 & 0.017 & No \\
 & EP-flux dose-response & 0.032 & 0.021 & No \\
 & \multicolumn{4}{l}{\footnotesize (7 additional tests with $P > 0.05$ omitted)} \\
\midrule
\multicolumn{5}{l}{\textit{B: Daily/station-day level ($n > 1000$; 7 tests)}} \\
 & All 7 tests & $<$0.05 & --- & All survive \\
\midrule
\multicolumn{5}{l}{\textit{C: Multi-country/external (7 tests)}} \\
 & 3 of 7 survive & & & 3/7 \\
\bottomrule
\end{tabular}
\end{table}

% ── Table: Accident secular trend and decade decomposition ──
\begin{table}[htbp]
\caption{Secular trend in Swiss avalanche accidents (1970--2025) and decade-stratified SSW--accident rate ratios. The overall accident rate shows a strong positive trend driven by expanding backcountry recreation. The Cochran--Mantel--Haenszel test (Table~\ref{tab:cmh_detrending}) demonstrates that the SSW--accident association persists after controlling for this secular trend.}
\label{tab:secular_trend}
\small
\begin{tabular}{lcccc}
\toprule
Period & Accidents per day & SSW rate & non-SSW rate & RR \\
\midrule
\multicolumn{5}{l}{\textit{Overall secular trend: $R^2 = 0.73$, slope $= +0.016$~acc~d$^{-1}$~yr$^{-1}$}} \\
\midrule
1970--1979 & 0.24 & 0.18 & 0.26 & 0.69 \\
1980--1989 & 0.33 & 0.28 & 0.35 & 0.80 \\
1990--1999 & 0.51 & 0.44 & 0.54 & 0.81 \\
2000--2009 & 0.75 & 0.91 & 0.69 & 1.32 \\
2010--2019 & 0.88 & 1.05 & 0.82 & 1.28 \\
2020--2025 & 0.96 & 1.34 & 0.74 & 1.81 \\
\bottomrule
\multicolumn{5}{l}{\footnotesize Pre-2000 RR $<$ 1 reflects low human backcountry exposure, not absence of mechanism.} \\
\end{tabular}
\end{table}

% ── Table: CMH decade-stratified detrending ──
\clearpage
\begin{table}[htbp]
\caption{Cochran--Mantel--Haenszel decade-stratified test of the SSW--accident association. Binary outcome (day with $\geq$1 fatal accident) is stratified by decade to control for the secular recreation trend. The common odds ratio indicates elevated accident risk during SSW windows across all six decades, with the association robust to trend removal (detrended residuals: SSW $-$ non-SSW $= +0.31$~acc~d$^{-1}$, Mann--Whitney $P < 10^{-4}$).}
\label{tab:cmh_detrending}
\small
\begin{tabular}{lccccc}
\toprule
Decade & SSW days & non-SSW days & SSW acc.\ rate & non-SSW acc.\ rate & RR \\
\midrule
1970s & 172 & 1,640 & 0.285 & 0.163 & 1.75 \\
1980s & 217 & 1,596 & 0.226 & 0.212 & 1.06 \\
1990s & 62 & 1,750 & 0.500 & 0.291 & 1.72 \\
2000s & 279 & 1,534 & 1.000 & 0.538 & 1.86 \\
2010s & 155 & 1,657 & 1.168 & 0.708 & 1.65 \\
2020s & 31 & 996 & 2.226 & 0.841 & 2.65 \\
\midrule
\multicolumn{6}{l}{\textit{Common OR $= 1.77$; $\chi^2 = 54.1$; $P < 10^{-12}$}} \\
\multicolumn{6}{l}{\textit{Detrended residual: $+0.31$~acc~d$^{-1}$ (Mann--Whitney $P < 10^{-4}$)}} \\
\bottomrule
\end{tabular}
\end{table}

%% ================================================================
%% Supplementary Table: Formal Power Analysis
%% ================================================================
\begin{table}[!ht]
\centering
\caption{\textbf{Supplementary Table 48: Formal power analysis for the Swiss event-level sign test ($n=16$).}
Power, minimum detectable effect (MDE), Bayes factor sensitivity and specification-curve null probability for the primary Swiss analysis.}
\label{tab:power_analysis}
\small
\begin{tabular}{lll}
\toprule
\textbf{Metric} & \textbf{Value} & \textbf{Interpretation} \\
\midrule
\multicolumn{3}{l}{\textit{Sign test power ($\alpha = 0.05$, two-sided)}} \\
\quad At observed $P(\mathrm{RR}<1) = 0.875$ & 79\% & Near-adequate \\
\quad At $P(\mathrm{RR}<1) = 0.90$ & 93\% & Well-powered \\
\quad At $P(\mathrm{RR}<1) = 0.95$ & 99\% & High power \\
\midrule
\multicolumn{3}{l}{\textit{Minimum detectable effect (80\% power)}} \\
\quad Cohen's $d$ (paired $t$-test equivalent) & 0.70 & Medium--large \\
\quad Observed $d$ & 1.06 & Well above MDE \\
\midrule
\multicolumn{3}{l}{\textit{Bayes factor sensitivity (sign test, $k=14/16$)}} \\
\quad Prior $\mathrm{Beta}(1,1)$ (uniform) & 178.7 & Decisive \\
\quad Prior $\mathrm{Beta}(2,2)$ (weakly informative) & 104.2 & Decisive \\
\quad Prior $\mathrm{Beta}(5,5)$ (moderately sceptical) & 17.9 & Strong \\
\midrule
\multicolumn{3}{l}{\textit{Specification curve}} \\
\quad Variants tested & 180 & \\
\quad Variants with $\mathrm{RR} < 1$ & 180/180 (100\%) & \\
\quad $P(\text{all consistent} \mid H_0)$ & \textit{Not reported: 180 specs share the same} & \\
\quad & \textit{16 events; effective df $\ll 180$} & \\
\midrule
\multicolumn{3}{l}{\textit{Community precedent for SSW event-level studies}} \\
\quad Baldwin \& Dunkerton (2001, \textit{Science}) & $n \approx 18$ & Comparable \\
\quad Sigmond et al.\ (2013, \textit{Nature Geosci.}) & $n = 13$ & Smaller \\
\quad Kolstad et al.\ (2010, \textit{QJRMS}) & $n = 29$ & Comparable \\
\bottomrule
\end{tabular}
\end{table}

%% ================================================================
%% Supplementary Table: Multi-Hazard Regime Redistribution
%% ================================================================
\begin{table}[!ht]
\centering
\caption{\textbf{Supplementary Table 49: Multi-hazard weather-variable shifts during SSW windows.}
ERA5-derived weather variables (Swiss Alpine domain-mean) during $\pm 15$-day SSW windows versus DOY-matched non-SSW controls. All values are from the ERA5 Swiss Alpine Extended dataset (1998--2019). Hazard-domain relevance column indicates which threshold-governed hazard each variable principally affects.}
\label{tab:multi_hazard_regime}
\small
\begin{tabular}{llllll}
\toprule
\textbf{Variable} & \textbf{SSW mean} & \textbf{Control mean} & \textbf{Shift} & \textbf{$P$} & \textbf{Hazard domain} \\
\midrule
Snow depth SWE (m) & 0.081 & 0.064 & $+26.1\%$ & $< 10^{-4}$ & Avalanche loading \\
Total precipitation (mm d$^{-1}$) & 0.322 & 0.385 & $-16.5\%$ & 0.04 & Flood, landslide \\
Rain fraction & 0.438 & 0.514 & $-7.6$~pp & $< 10^{-4}$ & Rain-on-snow trigger \\
Temperature (K) & 273.2 & 275.2 & $-2.1$~K & $< 10^{-4}$ & Cold-wave mortality \\
Wind speed (m s$^{-1}$) & 1.21 & 1.19 & $+1.9\%$ & 0.49 & Wind-slab loading \\
Cold-stagnation days & 4.2 & 3.9 & $+7.0\%$ & 0.09 & Air quality \\
\midrule
\multicolumn{6}{l}{\textit{Multi-hazard simultaneity: 69\% of events shift $\geq$2 domains (mean 2.1 per event)}} \\
\bottomrule
\end{tabular}
\end{table}

%% ================================================================
%% Supplementary Table 50: Blocking-Without-SSW Comparison
%% ================================================================
\begin{table}[!ht]
\centering
\caption{\textbf{Supplementary Table 50: Blocking regime comparison with and without SSW.}
Alpine blocking episodes (Z500 anomaly $> 1\sigma$ for $\geq$5 consecutive days) are classified by whether they overlap ($\pm$15 days) with an SSW onset. Suppression fraction indicates the proportion of episodes in which mean natural dry-slab AAI falls below the DOY-matched climatological mean. The Mann--Whitney test compares the distributions of episode-level rate ratios between SSW-associated and non-SSW blocking episodes.}
\label{tab:blocking_comparison}
\small
\begin{tabular}{lcccc}
\toprule
\textbf{Blocking regime} & \textbf{$n$ episodes} & \textbf{Suppression fraction} & \textbf{Sign-test $P$} & \textbf{Median RR} \\
\midrule
Non-SSW blocking & 44 & 82\% (36/44) & $< 10^{-4}$ & 0.41 \\
SSW-associated blocking & 12 & 42\% (5/12) & 0.77 & 0.95 \\
\midrule
\multicolumn{5}{l}{\textit{Mann--Whitney comparison of RR distributions: $P = 0.013$}} \\
\multicolumn{5}{l}{\textit{Interpretation: Non-SSW blocking produces fair-weather stagnation (cold, dry, stable $\to$ strong}} \\
\multicolumn{5}{l}{\textit{suppression). SSW-associated blocking coincides with active wave forcing that sustains moisture}} \\
\multicolumn{5}{l}{\textit{transport, offsetting the suppressive cold with structural loading $\to$ loaded-gun state at RR $\approx 1$.}} \\
\multicolumn{5}{l}{\textit{The $P = 0.013$ difference between regimes is the signature of SSW transforming blocking character.}} \\
\bottomrule
\end{tabular}
\end{table}

%% ================================================================
%% Supplementary Table 51: GEE/Design-Effect Accident Model
%% ================================================================
\begin{table}[!ht]
\centering
\caption{\textbf{Supplementary Table 51: Cluster-robust accident analysis.}
The raw Cochran--Mantel--Haenszel (CMH) test treats individual accident records as independent observations. Design-effect correction accounts for within-SSW-event clustering (ICC = intraclass correlation coefficient; DE = design effect = $1 + (m-1) \times \mathrm{ICC}$ where $m$ = mean cluster size). The corrected $\chi^2 = \mathrm{raw}\;\chi^2 / \mathrm{DE}$.}
\label{tab:gee_accident}
\small
\begin{tabular}{lcccc}
\toprule
\textbf{Method} & \textbf{OR} & \textbf{95\% CI} & \textbf{$\chi^2$} & \textbf{$P$} \\
\midrule
Raw CMH (individual records) & 1.77 & --- & 54.1 & $< 10^{-12}$ \\
Design-effect corrected & 1.77 & --- & 7.3 & 0.007 \\
Event-level permutation & 1.40 & --- & --- & 0.067 \\
\midrule
\multicolumn{5}{l}{\textit{Cluster structure: 29 SSW-event clusters (mean 22.7 records) + 55 control-period clusters (mean 76.5 records)}} \\
\multicolumn{5}{l}{\textit{Overall mean cluster size $\bar{m} = 58.6$ across 84 clusters; ICC = 0.11; DE = $1 + (58.6-1) \times 0.11 = 7.4$}} \\
\multicolumn{5}{l}{\textit{Effective sample size: 4,868 $\to$ $\sim$661 after clustering correction}} \\
\bottomrule
\end{tabular}
\end{table}

%% ================================================================
%% Supplementary Table 52: Precipitation Phase Analysis
%% ================================================================
\begin{table}[!ht]
\centering
\caption{\textbf{Supplementary Table 52: Precipitation phase shifts during SSW windows.}
ERA5-derived precipitation and temperature variables for the Swiss Alpine domain during $\pm$15-day SSW windows versus DOY-matched non-SSW controls. Snow fraction is calculated as snowfall / total precipitation. New-snow density is estimated using the Lehning et al.\ (2002) parameterisation from temperature and wind speed. Sintering rate change is estimated from Arrhenius kinetics of ice-grain bonding at the observed temperature shift.}
\label{tab:precip_phase}
\small
\begin{tabular}{lcccl}
\toprule
\textbf{Variable} & \textbf{SSW mean} & \textbf{Control mean} & \textbf{Change} & \textbf{$P$} \\
\midrule
Snow fraction of precip (\%) & 87.3 & 74.3 & $+13.0$~pp & $< 10^{-4}$ \\
Snowfall (mm d$^{-1}$ SWE) & 3.21 & 3.05 & $+5.2\%$ & 0.41 \\
Rainfall (mm d$^{-1}$) & 0.67 & 1.02 & $-33.9\%$ & $< 10^{-3}$ \\
Temperature ($^\circ$C) & $-0.0$ & $2.1$ & $-2.1$~K & $< 10^{-4}$ \\
Rain-on-snow days (\%) & 0.0 & 2.1 & $-100\%$ & --- \\
\midrule
Estimated new-snow density (kg m$^{-3}$) & 105 & 112 & $-6\%$ & --- \\
Estimated sintering rate change & --- & --- & $-14\%$ & --- \\
\bottomrule
\end{tabular}
\end{table}

%% --- Table 53: Cold-wave mortality risk ---
\begin{table}[htbp]
\centering
\caption{Cold-wave mortality risk estimation during SSW windows. Temperature distributions and cold-extreme frequencies are computed from ERA5 Alpine domain-mean 2m temperature during winter (November--March, 1998--2019). Mortality risk ratios apply published temperature--mortality dose--response coefficients (central estimate $\beta = 0.045$ per $^\circ$C below optimum; conservative $\beta = 0.030$; Gasparrini et al.\ 2015\cite{Gasparrini2015}) to the observed temperature distributions. Bootstrap CIs use 10,000 resamples of SSW-window days.}
\label{tab:cold_wave}
\begin{tabular}{lrrrl}
\toprule
\textbf{Variable} & \textbf{SSW windows} & \textbf{Control} & \textbf{Change} & \textbf{$P$} \\
\midrule
Mean temperature ($^\circ$C) & $0.02$ & $2.07$ & $-2.1$~K & $< 10^{-4}$ \\
Days $< P_{1}$ ($-7.8^\circ$C) (\%) & 1.2 & 1.0 & $\times 1.2$ & --- \\
Days $< P_{5}$ ($-5.1^\circ$C) (\%) & 9.3 & 4.1 & $\times 2.3$ & $< 10^{-4}$ \\
Days $< P_{10}$ ($-3.7^\circ$C) (\%) & 16.3 & 8.6 & $\times 1.9$ & $< 10^{-4}$ \\
Days $< P_{25}$ ($-1.2^\circ$C) (\%) & 36.5 & 22.5 & $\times 1.6$ & $< 10^{-4}$ \\
\midrule
Cold spells ($\geq 3$ days $< P_{10}$) & 15 & 30 & --- & --- \\
Mean cold-spell duration (days) & 4.7 & 4.8 & --- & --- \\
\midrule
Estimated excess mortality RR (conservative) & 1.047 & 1.028 & $+1.9\%$ & --- \\
Estimated excess mortality RR (central) & 1.073 & 1.043 & $+2.9\%$ & --- \\
Bootstrap 95\% CI for excess RR & \multicolumn{3}{c}{$[1.019,\; 1.039]$} & --- \\
\midrule
Cold-stagnation compound (\%) & 9.7 & 6.1 & $\times 1.6$ & --- \\
Heating degree days (base 15.5$^\circ$C) & 15.5 & 13.4 & $+15.3\%$ & --- \\
\midrule
Events with colder SSW window & \multicolumn{3}{c}{9/16 (56\%; $P = 0.80$)} & --- \\
Events with excess mortality risk & \multicolumn{3}{c}{8/16 (50\%; $P = 1.00$)} & --- \\
\bottomrule
\end{tabular}
\end{table}

%% ================================================================
%% Supplementary Table 54: Pre-event snowpack comparison
%% ================================================================
\begin{table}[!ht]
\centering
\caption{\textbf{Supplementary Table 54: Pre-event snowpack comparison.}
Antecedent snowpack conditions in the 30-to-16-day window before SSW onset (before the SSW meteorological anomaly begins) compared with day-of-year-matched windows from non-SSW winters. Event-level tests aggregate station-day values to event/winter means and use Mann--Whitney $U$ tests ($n = 16$ SSW events vs.\ $n = 9$ non-SSW winters). The non-significance of all event-level tests---particularly persistent-weak-layer (PWL) prevalence ($P = 0.93$)---rules out pre-existing snowpack as a confounder. Station-day-level significance reflects the large sample ($N = 17{,}526$ SSW vs.\ $N = 127{,}448$ control station-days) rather than meaningful effect sizes.}
\label{tab:pre_event_snowpack}
\small
\begin{tabular}{lcccccl}
\toprule
\textbf{Variable} & \textbf{SSW mean} & \textbf{Ctrl mean} & \textbf{$\Delta\%$} & \textbf{$d$} & \textbf{$P$ (event)} & \textbf{Interpretation} \\
\midrule
Snow depth (cm) & 114.8 & 95.3 & $+20.4$ & $+0.55$ & 0.32 & Not significant \\
SWE (kg\,m$^{-2}$) & 348.2 & 283.2 & $+22.9$ & $+0.51$ & 0.32 & Not significant \\
PWL prevalence & 0.721 & 0.739 & $-2.4$ & $-0.15$ & 0.93 & Not significant \\
\midrule
\multicolumn{7}{l}{\textit{Station-day level (informative only; significance reflects sample size):}} \\
New snow 24h (cm) & 6.45 & 5.05 & $+27.7$ & $+0.13$ & $< 0.001$ & Trivial $d$ \\
Natural stability ($S_n$) & 2.26 & 2.71 & $-16.9$ & $-0.31$ & $< 0.001$ & Small $d$ \\
Deformation stability ($S_d$) & 1.12 & 1.44 & $-22.5$ & $-0.24$ & $< 0.001$ & Small $d$ \\
Surface hoar (mm) & 0.41 & 0.55 & $-26.5$ & $-0.12$ & 0.002 & Trivial $d$ \\
\bottomrule
\end{tabular}
\end{table}

% ============================================================
%  NEW SUPPLEMENTARY TABLES (R55 analyses)
% ============================================================

\begin{table}[!ht]
\centering
\caption{\textbf{Supplementary Table 55: Placebo and falsification tests.}
Signal specificity assessed via temporal displacement and non-SSW weak-vortex comparisons. Only real SSW dates produce the strong suppression signal. The weak-vortex non-SSW test selects 16 non-SSW winter periods with the largest U10 deceleration (matched to SSW event count); these show no suppression (gmRR $= 1.06$), confirming that the signal is specific to SSW events and not merely any period of vortex weakening.}
\label{tab:placebo}
\small
\begin{tabular}{lccccl}
\toprule
Test & gmRR & $n$ tested & $n$ decrease & Sign $P$ & Result \\
\midrule
Actual SSW dates & 0.45 & 16 & 14 & 0.002 & Strong suppression \\
$+$60 day shift & 0.90 & 9 & 4 & 1.000 & No signal \\
$-$60 day shift & 0.68 & 13 & 9 & 0.267 & No signal \\
Weak-vortex non-SSW & 1.06 & 16 & 7 & 0.804 & No suppression \\
Random-date permutations & --- & 1,000 iter. & --- & --- & 7.3\% $\geq$ 14/16 \\
\bottomrule
\end{tabular}
\end{table}

\begin{table}[!ht]
\centering
\caption{\textbf{Supplementary Table 56: Power simulation for post-2005 subperiod.}
Monte Carlo power analysis for the post-2005 subset ($n = 9$ events). At the full-period effect size ($P(\mathrm{decrease}) = 0.78$), the exact sign test achieves only 37.9\% power---insufficient to reliably detect the effect. The observed $P = 0.18$ falls at the 69th percentile of simulated outcomes under the true effect, indicating the post-2005 subset is underpowered to distinguish persistence from attenuation.}
\label{tab:power_sim}
\small
\begin{tabular}{lcc}
\toprule
Parameter & Value & Interpretation \\
\midrule
$n$ events (post-2005) & 9 & \\
True $P$(decrease) & 0.78 & From full-period estimate \\
Power at $\alpha = 0.05$ & 37.9\% & Underpowered \\
Observed $P$ & 0.18 & \\
Percentile of observed $P$ & 69th & Consistent with real effect \\
$P$(obtaining $P \leq 0.05$ $|$ true effect) & 0.379 & \\
\bottomrule
\end{tabular}
\end{table}

\begin{table}[!ht]
\centering
\caption{\textbf{Supplementary Table 57: Granger causality tests.}
Pairwise Granger causality between daily zonal-mean zonal wind at 10~hPa (U10) and Alpine-sector Z500 anomaly across the full 21-winter panel. Both directions are significant at all lags tested, but the stratosphere-to-troposphere direction shows consistently larger $F$-statistics, suggesting stronger top-down predictive influence. Bidirectional significance indicates coupled dynamics rather than unidirectional forcing.}
\label{tab:granger}
\small
\begin{tabular}{lccccc}
\toprule
Direction & Lag & $F$-statistic & $P$ & Significant & Interpretation \\
\midrule
U10 $\to$ Z500 & 1 & 45.6 & $< 10^{-10}$ & Yes & Stronger \\
U10 $\to$ Z500 & 7 & 8.2 & $< 10^{-8}$ & Yes & Stronger \\
U10 $\to$ Z500 & 14 & 5.1 & $< 10^{-7}$ & Yes & Stronger \\
\midrule
Z500 $\to$ U10 & 1 & 26.9 & $< 10^{-6}$ & Yes & Weaker \\
Z500 $\to$ U10 & 7 & 4.7 & $< 10^{-4}$ & Yes & Weaker \\
Z500 $\to$ U10 & 14 & 3.2 & $< 10^{-3}$ & Yes & Weaker \\
\bottomrule
\end{tabular}
\end{table}

\begin{table}[!ht]
\centering
\caption{\textbf{Supplementary Table 58: Lead-time analysis.}
Composite Z500 anomaly timing relative to SSW onset. Z500 drops below $-1$~SD at a mean of $+10.5$~days after onset (5/8 measurable events), suggesting SSW onset precedes the later maturation of the large-scale blocking anomaly. Avalanche suppression onset, by contrast, precedes or coincides with SSW onset in 14/16 events (mean $-7.25$~days), consistent with planetary-wave preconditioning affecting both stratosphere and troposphere simultaneously.}
\label{tab:lead_time}
\small
\begin{tabular}{lcc}
\toprule
Metric & Value & Interpretation \\
\midrule
Mean Z500 $< -1$~SD lag & $+10.5$~d after onset & Blocking matures after SSW \\
Events with measurable lag & 5/8 & Subset with clear threshold \\
Mean avalanche suppression onset & $-7.25$~d & Precedes SSW onset \\
Events with pre-onset suppression & 14/16 & Planetary-wave preconditioning \\
\bottomrule
\end{tabular}
\end{table}

%% ---- NEW ANALYSES: R57 ----

\begin{table}[!ht]
\centering
\caption{\textbf{Supplementary Table 59: Pre-onset versus post-onset SSW suppression decomposition.}
Avalanche deficit partitioned into pre-onset ($-15$ to $-1$~d) and post-onset ($0$ to $+15$~d) windows relative to the formal 10~hPa wind reversal. Pre-onset accounts for 52\% of the total deficit, demonstrating that suppression begins before the stratospheric event and supporting the common-cause planetary-wave mechanism.}
\label{tab:pre_post_decomp}
\small
\begin{tabular}{lcc}
\toprule
Statistic & Pre-onset ($-15$ to $-1$~d) & Post-onset ($0$ to $+15$~d) \\
\midrule
Median event-level RR & 0.21 & 0.26 \\
Total deficit (obs $-$ exp) & 177.4 counts & 164.2 counts \\
Fraction of total deficit & 52\% & 48\% \\
Paired Wilcoxon $P$ & \multicolumn{2}{c}{0.98 (no asymmetry)} \\
Paired $t$-test $P$ & \multicolumn{2}{c}{0.71} \\
\bottomrule
\end{tabular}
\end{table}

\begin{table}[!ht]
\centering
\caption{\textbf{Supplementary Table 60: Composite eddy heat flux ($\overline{v'T'}$) around SSW onset.}
Poleward eddy heat flux at 100~hPa (45--75$^\circ$N, cos-latitude-weighted; NCEP--NCAR reanalysis) composited across 16~SSW events. The pre-onset burst confirms that planetary-wave forcing is systematically and strongly elevated before every SSW event (2.1$\times$ DJF climatology). The event-level dose--response null ($\rho = -0.25$, $P = 0.34$) reflects range restriction: all SSW events require strong wave forcing by definition, compressing cross-event variability in the stratospheric ``dose.''}
\label{tab:vt_composite}
\small
\begin{tabular}{lccc}
\toprule
Phase & Mean $\overline{v'T'}$ anomaly & 95\% CI & $P$ \\
 & (K$\cdot$m~s$^{-1}$) & & \\
\midrule
Far pre-onset ($-30$ to $-16$~d) & $-0.4$ & $[-5.2, 4.4]$ & ns \\
Pre-onset ($-15$ to $-1$~d) & $+9.4$ & $[+4.1, +14.9]$ & $< 10^{-5}$ \\
Post-onset ($0$ to $+15$~d) & $+1.9$ & $[-2.7, +6.5]$ & ns \\
\midrule
Peak anomaly (lag $-1$~d) & $+17.5$ & $[+12.3, +23.4]$ & $< 10^{-6}$ \\
DJF climatology & $15.6$ & --- & --- \\
Peak raw value at onset & $35.6$ & --- & 2.1$\times$ clim. \\
\bottomrule
\end{tabular}
\end{table}

\begin{table}[!ht]
\centering
\caption{\textbf{Supplementary Table 61: Trigger-budget resolution of the wind-transport paradox.}
Channel-weighted trigger-budget accounting during SSW versus control days. Wind-slab loading accounts for only $\sim$12\% of natural-release trigger weight; the $+25\%$ wind-transport increase contributes $+3$ weighted percentage points, which is overwhelmed by the $-8$ weighted points from warming-event suppression. The net trigger balance is near zero (95\% CI spans zero), confirming that wind loading does not contradict the overall suppression signal.}
\label{tab:trigger_budget}
\small
\begin{tabular}{lccccc}
\toprule
Trigger channel & SSW $\Delta$~(\%) & Weight & Weighted $\Delta$ & Direction & $P$ \\
\midrule
Warming events & $-22.4$ & 0.35 & $-7.8$ & Suppressive & $< 0.001$ \\
Solar radiation & $-17.5$ & 0.08 & $-1.4$ & Suppressive & --- \\
Rain-on-snow & $+9.1$ & 0.15 & $+1.4$ & Activating & 0.63 \\
Wind transport ($u^3$) & $+25.4$ & 0.12 & $+3.0$ & Activating & 0.024 \\
New-snow loading & $+25.7$ & 0.30 & $+7.7$ & Activating & 0.026 \\
\midrule
Net weighted $\Delta$ & \multicolumn{2}{l}{$+2.9\%$} & \multicolumn{3}{l}{95\% CI: $[-10.0, +16.8]$} \\
\bottomrule
\end{tabular}
\end{table}

\begin{table}[!ht]
\centering
\caption{\textbf{Supplementary Table 62: Model-comparison for the SSW--Z500 daily count regression.}
Six specifications of the daily avalanche count model ($n = 2{,}702$ winter days, 21 winters) predicting natural dry-slab counts from SSW and standardised Z500. Model selection by AIC/BIC favours the zero-inflated negative-binomial (ZINB); the Vuong test rejects ZIP in favour of ZINB ($Z = -4.2$, $P < 0.0001$). SSW retains significance ($P \leq 0.005$) under all specifications that handle overdispersion (NB2, ZINB) but falls short of $\alpha = 0.05$ under GEE with winter-level clustering ($P = 0.18$--$0.20$), reflecting the fundamental $n = 16$ event-level constraint. Z500 is significant in models that do not adequately account for the zero-inflation structure (Poisson, NB2, ZIP, GEE); however, the best-fitting model (ZINB) absorbs the Z500 signal into the zero-inflation component (IRR $= 1.03$, $P = 0.84$), indicating collinearity between SSW and the subset of days when avalanche release is structurally possible. This is the expected result when SSW determines whether the system is in a ``zero state'' (no trigger days) versus an ``active state'' (release possible): once the zero-inflation is correctly modelled, the count-model coefficient of a correlated predictor becomes redundant. The collinearity does not invalidate the Z500 mediation identified at the event level ($\rho_{\mathrm{partial}} = 0.64$, $P = 0.008$)---it reflects a scale-of-analysis distinction between event-level pathway attribution and daily-level count decomposition.}
\label{tab:nb_comparison}
\small
\begin{tabular}{lcccccc}
\toprule
Model & SSW IRR & SSW $P$ & Z500 IRR & Z500 $P$ & AIC & BIC \\
\midrule
Poisson & 0.54 & $< 10^{-24}$ & 0.79 & $< 10^{-30}$ & 17{,}293 & 17{,}311 \\
NB2 & 0.55 & 0.005 & 0.75 & 0.001 & 4{,}627 & 4{,}651 \\
ZIP & 0.41 & $< 10^{-44}$ & 0.94 & 0.008 & 9{,}593 & 9{,}628 \\
ZINB & 0.40 & $< 10^{-4}$ & 1.03 & 0.84 & 4{,}601 & 4{,}642 \\
GEE-Poisson & 0.71 & 0.20 & 0.68 & $< 10^{-8}$ & --- & --- \\
GEE-NB & 0.70 & 0.18 & 0.68 & $< 10^{-8}$ & --- & --- \\
\bottomrule
\end{tabular}
\end{table}

\begin{table}[!ht]
\centering
\caption{\textbf{Supplementary Table 63b: GEE power simulation results.}
Parametric bootstrap simulation ($N = 50{,}000$ iterations) estimating the statistical power of GEE with winter-level exchangeable correlation to detect the observed SSW effect ($\mathrm{IRR} = 0.70$; $\beta = -0.357$) given $n_{\mathrm{eff}} = 16$ SSW events across 21 winters. The design effect (GEE SE$^2$ / naive SE$^2$ = 1.65) inflates standard errors, reducing power to 26\%. The observed $P = 0.18$ falls at the 50th percentile of the p-value distribution under $H_1$, i.e.\ it is the \emph{median expected result} if the true effect is a 30\% reduction. Eighty per cent power would require either a much stronger effect ($\mathrm{IRR} \leq 0.47$) or approximately 72 SSW events---roughly four times the available sample.}
\label{tab:power_sim}
\small
\begin{tabular}{lc}
\toprule
Parameter & Value \\
\midrule
True IRR (assumed) & 0.70 \\
True $\beta$ & $-0.357$ \\
Naive SE & 0.21 \\
GEE SE & 0.27 \\
Design effect & 1.65 \\
\midrule
Analytical power ($\alpha = 0.05$) & 26.2\% \\
Simulated power ($N = 50{,}000$) & 26.3\% \\
\midrule
Observed $P$ value & 0.18 \\
Percentile of $P = 0.18$ under $H_1$ & 49.5th \\
Median $P$ under $H_1$ & 0.184 \\
\midrule
IRR needed for 80\% power ($n = 16$) & 0.469 \\
$n_{\mathrm{events}}$ needed for 80\% power (IRR $= 0.70$) & 72 \\
\bottomrule
\end{tabular}
\end{table}

\begin{table}[!ht]
\centering
\caption{\textbf{Supplementary Table 64: EP-flux proxy diagnostics composited across 16 SSW events.}
Zonal-wind deceleration rate ($du/dt$ at 10~hPa) serves as a proxy for EP-flux convergence; stratospheric temperature tendency ($dT/dt$ at 100~hPa) reflects the vertical component of EP-flux (Eliassen--Palm relation). The onset-phase deceleration ($-2.12$~m~s$^{-1}$~d$^{-1}$, $P < 10^{-4}$) confirms strong wave forcing. Downward propagation from 10 to 100~hPa occurs in $\sim$6.6~days, consistent with the 7--10~day timescale in the SSW literature.}
\label{tab:epflux_diag}
\small
\begin{tabular}{lccc}
\toprule
Diagnostic & Precursor ($-30$ to $-6$~d) & Onset ($-5$ to $+5$~d) & Recovery ($+6$ to $+30$~d) \\
\midrule
$du/dt$ at 10~hPa (m~s$^{-1}$~d$^{-1}$) & $-0.40$ & $-2.12^{***}$ & $+0.51^{**}$ \\
$dT/dt$ at 100~hPa (K~d$^{-1}$) & $+0.07^{*}$ & $+0.44^{***}$ & $\sim 0$ (ns) \\
T anomaly at 10~hPa (K) & $\sim 0$ (ns) & $+11.9^{***}$ & $\sim 0$ (ns) \\
Z500 anomaly RMS at 10~hPa (m) & 311 & 671 & 466 \\
\midrule
\multicolumn{4}{l}{Downward propagation: 10$\to$50~hPa in $\sim$1.4~d; 10$\to$100~hPa in $\sim$6.6~d} \\
\bottomrule
\end{tabular}
\end{table}

\begin{table}[!ht]
\centering
\caption{\textbf{Supplementary Table 65: Lead--lag falsification analysis.}
Sign-count suppression rate and one-sided $P$-values for the event-level analysis at different temporal offsets from SSW onset. The SSW-aligned window (lag~0) shows the highest suppression rate (14/16 = 88\%). The signal is weaker at $\pm 60$~days (7/13 = 54\% and 8/12 = 67\%, respectively), although a rebound at $+45$~days (14/16 = 88\%) likely reflects the 4--6-week persistence of SSW surface impacts documented in Baldwin \& Dunkerton (2001). $P$-values from one-sided $t$-test on log(RR) where computable (events with zero observed counts excluded from gmRR but retained for sign count). Note: $n$ varies across lags because some events fall outside the available data window at large offsets.}
\label{tab:lead_lag}
\small
\begin{tabular}{rcccc}
\toprule
Lag (days) & $n$ events & Suppressed (\%) & gmRR & $P_{\mathrm{one-sided}}$ \\
\midrule
$-60$ & 13 & 7 (54\%) & --- & --- \\
$-45$ & 15 & 12 (80\%) & --- & --- \\
$-30$ & 16 & 11 (69\%) & --- & --- \\
$-15$ & 16 & 12 (75\%) & 0.37 & 0.005 \\
$0$ & 16 & 14 (88\%) & 0.28 & $<0.001$ \\
$+15$ & 16 & 13 (81\%) & 0.26 & $<0.001$ \\
$+30$ & 16 & 13 (81\%) & --- & --- \\
$+45$ & 16 & 14 (88\%) & --- & --- \\
$+60$ & 12 & 8 (67\%) & --- & --- \\
\bottomrule
\end{tabular}
\end{table}

\begin{table}[!ht]
\centering
\caption{\textbf{Supplementary Table 66: Influence and fragility analysis.}
Leave-one-event-out (LOO) results for all 16 SSW events. No single event drives the result: all folds maintain gmRR $< 1$ and sign-test $P < 0.005$. The fragility index is 3 (three events would need to change direction from suppression to enhancement to lose $P < 0.05$), indicating moderate robustness given the small sample.}
\label{tab:fragility}
\small
\begin{tabular}{p{2.5cm}ccc}
\toprule
Dropped event & LOO gmRR & LOO $n$ suppressed & LOO sign $P$ \\
\midrule
Range across 16 folds & 0.24--0.32 & 13--14 (of 15) & 0.0005--0.004 \\
\midrule
Fragility index & \multicolumn{3}{l}{3 events (from 14/16 to 11/16 loses $P < 0.05$)} \\
\bottomrule
\end{tabular}
\end{table}

\begin{thebibliography}{30}

\bibitem{Charlton2007}
Charlton, A.~J.\ \& Polvani, L.~M. A new look at stratospheric sudden warmings.\ Part~I: Climatology and modeling benchmarks. \textit{J.\ Climate}\ \textbf{20}, 449--469 (2007).

\bibitem{Baldwin2001}
Baldwin, M.~P.\ \& Dunkerton, T.~J. Stratospheric harbingers of anomalous weather regimes. \textit{Science}\ \textbf{294}, 581--584 (2001).

\bibitem{Thompson2002}
Thompson, D.~W.~J., Baldwin, M.~P.\ \& Wallace, J.~M. Stratospheric connection to Northern Hemisphere wintertime weather: Implications for prediction. \textit{J.\ Climate}\ \textbf{15}, 1421--1428 (2002).

\bibitem{Kidston2015}
Kidston, J.\ et~al. Stratospheric influence on tropospheric jet streams, storm tracks and surface weather. \textit{Nature Geosci.}\ \textbf{8}, 433--440 (2015).

\bibitem{Sigmond2013}
Sigmond, M., Scinocca, J.~F., Kharin, V.~V.\ \& Shepherd, T.~G. Enhanced seasonal forecast skill following stratospheric sudden warmings. \textit{Nature Geosci.}\ \textbf{6}, 98--102 (2013).

\bibitem{Techel2016}
Techel, F.\ et~al. Avalanche fatalities in the European Alps: long-term trends and statistics. \textit{Geogr.\ Helv.}\ \textbf{71}, 147--159 (2016).

\bibitem{Schweizer2003}
Schweizer, J., Bruce~Jamieson, J.\ \& Schneebeli, M. Snow avalanche formation. \textit{Rev.\ Geophys.}\ \textbf{41}, 1016 (2003).

\bibitem{Schweizer2008}
Schweizer, J., Kronholm, K., Jamieson, J.~B.\ \& Birkeland, K.~W. Review of spatial variability of snowpack properties and its importance for avalanche formation. \textit{Cold Reg.\ Sci.\ Technol.}\ \textbf{51}, 253--272 (2008).

\bibitem{McClung2006}
McClung, D.~M.\ \& Schaerer, P. \textit{The Avalanche Handbook} 3rd edn (The Mountaineers Books, 2006).

\bibitem{Reuter2018}
Reuter, B.\ \& Schweizer, J. Describing snow instability by failure initiation, crack propagation, and slab tensile support. \textit{Geophys.\ Res.\ Lett.}\ \textbf{45}, 7019--7027 (2018).

\bibitem{Butler2015}
Butler, A.~H.\ et~al. Defining sudden stratospheric warmings. \textit{Bull.\ Am.\ Meteorol.\ Soc.}\ \textbf{96}, 1913--1928 (2015).

\bibitem{Polvani2004}
Polvani, L.~M.\ \& Waugh, D.~W. Upward wave activity flux as a precursor to extreme stratospheric events and subsequent anomalous surface weather regimes. \textit{J.\ Climate}\ \textbf{17}, 3548--3554 (2004).

\bibitem{Wilcox2017}
Wilcox, R.~R. \textit{Introduction to Robust Estimation and Hypothesis Testing} 4th edn (Academic Press, 2017).

\bibitem{Hobbs1974}
Hobbs, P.~V. \textit{Ice Physics} (Clarendon Press, 1974).

\bibitem{Colbeck1997}
Colbeck, S.~C. A review of sintering in seasonal snow. CRREL Report 97-10 (US Army Corps of Engineers, 1997).

\bibitem{Mitchell2013}
Mitchell, D.~M., Gray, L.~J., Anstey, J., Baldwin, M.~P.\ \& Charlton-Perez, A.~J. The influence of stratospheric vortex displacements and splits on surface climate. \textit{J.\ Climate}\ \textbf{26}, 2668--2682 (2013).

\bibitem{Hitchcock2013}
Hitchcock, P.\ \& Simpson, I.~R. The downward influence of stratospheric sudden warmings. \textit{J.\ Atmos.\ Sci.}\ \textbf{71}, 3856--3876 (2014).

\bibitem{Bartelt2002}
Bartelt, P.\ \& Lehning, M. A physical SNOWPACK model for the Swiss avalanche warning. Part I: numerical model. \textit{Cold Reg.\ Sci.\ Technol.}\ \textbf{35}, 123--145 (2002).

\bibitem{Lehning2002}
Lehning, M., Bartelt, P., Brown, B.\ \& Fierz, C. A physical SNOWPACK model for the Swiss avalanche warning. Part III: meteorological forcing, thin layer formation and evaluation. \textit{Cold Reg.\ Sci.\ Technol.}\ \textbf{35}, 169--184 (2002).

\bibitem{Kolstad2010}
Kolstad, E.~W., Breiteig, T.\ \& Scaife, A.~A. The association between stratospheric weak polar vortex events and cold air outbreaks in the Northern Hemisphere. \textit{Q.\ J.\ R.\ Meteorol.\ Soc.}\ \textbf{136}, 886--893 (2010).

\bibitem{KassRaftery1995}
Kass, R.~E.\ \& Raftery, A.~E. Bayes factors. \textit{J.\ Am.\ Statist.\ Assoc.}\ \textbf{90}, 773--795 (1995).

\bibitem{Mellor1975}
Mellor, M. A review of basic snow mechanics. \textit{IAHS Publication} \textbf{114}, 251--291 (1975).

\bibitem{Simenhois2009}
Simenhois, R.\ \& Birkeland, K.~W. The Extended Column Test: test effectiveness, spatial variability, and comparison with the Rutschblock. \textit{Cold Reg.\ Sci.\ Technol.}\ \textbf{59}, 210--216 (2009).

\bibitem{Gray2010}
Gray, L.~J.\ et~al. Solar influences on climate. \textit{Rev.\ Geophys.}\ \textbf{48}, RG4001 (2010).

\bibitem{Butler2017}
Butler, A.~H.\ et~al. A sudden stratospheric warming compendium. \textit{Earth Syst.\ Sci.\ Data}\ \textbf{9}, 63--76 (2017).

\bibitem{Gasparrini2015}
Gasparrini, A.\ et~al. Mortality risk attributable to high and low ambient temperature: a multicountry observational study. \textit{Lancet}\ \textbf{386}, 369--375 (2015).

\bibitem{Baldwin2021}
Baldwin, M.~P.\ et~al. Sudden stratospheric warmings. \textit{Rev.\ Geophys.}\ \textbf{59}, e2020RG000708 (2021).

\bibitem{Bauer2008}
Bauer, T.~K.\ \& Sinning, M. An extension of the Blinder--Oaxaca decomposition to nonlinear models. \textit{AStA Adv.\ Stat.\ Anal.}\ \textbf{92}, 197--206 (2008).

\bibitem{Cohen2021}
Cohen, J.\ et~al. Linking Arctic variability and change with extreme winter weather in the United States. \textit{Science}\ \textbf{373}, 1116--1121 (2021).

\bibitem{Kretschmer2018}
Kretschmer, M., Coumou, D., Agel, L., Barlow, M., Tziperman, E.\ \& Cohen, J. More-persistent weak stratospheric polar vortex states linked to cold extremes. \textit{Bull.\ Am.\ Meteorol.\ Soc.}\ \textbf{99}, 49--60 (2018).

\bibitem{Vernay2022}
Vernay, M.\ et~al. The S2M meteorological and snow cover reanalysis over the French mountainous areas: description and evaluation (1958--2021). \textit{Earth Syst.\ Sci.\ Data}\ \textbf{14}, 1707--1733 (2022).

\end{thebibliography}

% ============================================================
%  NEW TABLES: SNOWPACK Mechanistic Chain, Elevation, Temporal, Combinatorial, Split/Displacement
% ============================================================

\subsection*{SNOWPACK Mechanistic Chain Diagnostics}

\begin{table}[htbp]
\centering
\caption{SNOWPACK station-day \emph{pooled} diagnostics during SSW versus matched non-SSW windows across 130 Swiss stations (292,837 station-days; 11 of 16 events with full coverage). Each station-day receives equal weight; for event-averaged estimates, see Supplementary Table~\ref{tab:snowpack}. Effect sizes are Cohen's $d$; all $P$ values are from Mann--Whitney $U$ tests. The unit of inference for causal claims remains the 16 event-level rate ratios; these station-day comparisons quantify process consistency with ERA5-forced conditions.}
\label{tab:snowpack_mechanism}
\begin{tabular}{lrrrl}
\toprule
Variable & SSW mean & Control mean & Cohen's $d$ & Direction \\
\midrule
\multicolumn{5}{l}{\textit{Step 1: Temperature \& gradients}} \\
Air temperature ($^\circ$C) & $-6.89$ & $-3.90$ & $-0.58$ & Colder \\
Deep snow temp TS2 ($^\circ$C) & $-4.86$ & $-3.97$ & $-0.24$ & Colder \\
Temp.\ gradient proxy & 3.34 & 2.71 & $+0.22$ & Steeper \\
\midrule
\multicolumn{5}{l}{\textit{Step 2: Persistent weak-layer formation}} \\
PWL prevalence (pwl\_100) & 0.784 & 0.626 & $+0.33$ & $+25\%$ \\
Hoar crystal size & 0.336 & 0.287 & $+0.05$ & $+17\%$ \\
\midrule
\multicolumn{5}{l}{\textit{Step 3: Structural stability}} \\
Structural stability (SSI) & 2.67 & 3.49 & $-0.37$ & $-24\%$ \\
Min.\ critical crack length & 0.43 & 0.70 & $-0.35$ & $-39\%$ \\
Mean stability depth (cm) & 34.0 & 55.7 & $-0.36$ & $-39\%$ \\
\midrule
\multicolumn{5}{l}{\textit{Step 4: Trigger balance}} \\
Skier trigger index (sk38) & 2.03 & 3.02 & $-0.38$ & $-33\%$ \\
Station-level rainfall & 0.006 & 0.015 & $-0.02$ & $-59\%$ \\
New snow (HN24, cm) & 5.58 & 4.78 & $+0.08$ & $+17\%$ \\
SWE (mm) & 531 & 473 & $+0.17$ & $+12\%$ \\
Wind transport (24h) & 7.94 & 5.47 & $+0.17$ & $+45\%$ \\
Danger level & 2.46 & 2.12 & $+0.44$ & $+16\%$ \\
\bottomrule
\end{tabular}
\end{table}

\begin{table}[htbp]
\centering
\caption{Elevation stratification of SNOWPACK diagnostics during SSW windows. PWL and SWE changes by elevation band.}
\label{tab:elev_strat}
\begin{tabular}{lrrr}
\toprule
Elevation band & $n$ stations & SWE change (\%) & PWL change (\%) \\
\midrule
$<$1,500~m & 1 & $-22.4$ & $-1.2$ \\
1,500--2,000~m & 21 & $+22.4$ & $+64.9$ \\
2,000--2,500~m & 80 & $+10.9$ & $+24.8$ \\
$>$2,500~m & 28 & $+11.4$ & $+18.8$ \\
\bottomrule
\end{tabular}
\end{table}

\begin{table}[htbp]
\centering
\caption{Temporal structure of SNOWPACK diagnostics relative to SSW onset. Pre-SSW = $-15$ to $0$~d; onset = 0 to $+15$~d; late = $+15$ to $+30$~d. PWL prevalence peaks at day $+4$ after onset.}
\label{tab:temporal_mechanism}
\begin{tabular}{lrrr}
\toprule
Variable & Pre-SSW & Onset (0--15~d) & Late (15--30~d) \\
\midrule
PWL prevalence & 0.737 & 0.789 & 0.777 \\
Danger level & 2.35 & 2.46 & 2.50 \\
Air temperature ($^\circ$C) & $-5.87$ & $-6.45$ & $-7.61$ \\
Wind transport & 8.96 & 7.95 & 8.01 \\
New snow (HN24) & 5.84 & 5.84 & 5.45 \\
\bottomrule
\end{tabular}
\end{table}

\begin{table}[htbp]
\centering
\caption{Multivariate directional coherence across eight measurement channels during SSW windows. Under independence, $P(\text{8/8 concordant}) = 2^{-8} = 0.004$. Adjusting for inter-channel correlation yields $k_{\mathrm{eff}} \approx 5$ ($P = 0.031$). Reported as a consistency diagnostic; channel selection was guided by the loaded-gun hypothesis.}
\label{tab:combinatorial}
\begin{tabular}{lll}
\toprule
Channel & Predicted direction & Observed \\
\midrule
SWE & Increase & $+12\%$ \\
Temperature & Decrease & $-3.0\,^\circ$C \\
Rainfall & Decrease & $-59\%$ \\
Snow fraction & Increase & $+12$~pp \\
Natural dry-slab counts & Decrease & gmRR $= 0.32$ \\
Danger ratings & Increase & 11/14 events \\
Human-triggered counts & Increase & RR $= 1.45$ \\
Faceted-crystal fraction & Increase & $+18\%$ \\
\bottomrule
\end{tabular}
\end{table}

\begin{table}[htbp]
\centering
\caption{SSW event-type stratification: consistency check using the Butler \& Polvani\cite{Butler2015} morphological classification. Both types produce near-identical suppression (Mann--Whitney $P = 0.95$; Cohen's $d = 0.02$), demonstrating that the signal is type-invariant. The two contrary events span both morphologies. Full details in Supplementary Table~\ref{tab:ssw_type} and Extended Data Table~\ref{tab:sensitivity}.}
\label{tab:split_displacement}
\begin{tabular}{lrrrr}
\toprule
Event type & $n$ & gmRR & Sign concordance & Sign-test $P$ \\
\midrule
Displacement & 11 & 0.33 & 10/11 (91\%) & 0.006 \\
Split-vortex & 5 & 0.32 & 4/5 (80\%) & 0.043 \\
All events & 16 & 0.32 & 14/16 (88\%) & 0.004 \\
\bottomrule
\multicolumn{5}{l}{\footnotesize Between-type comparison: Mann--Whitney $P = 0.95$; Cohen's $d = 0.02$.} \\
\multicolumn{5}{l}{\footnotesize Contrary events: Feb 2018 (split, RR$=1.05$); Jan 2019 (displacement, RR$=3.43$).} \\
\end{tabular}
\end{table}

% ================================================================
%  Supplementary Table: Sign-Randomisation Inference
% ================================================================
\begin{table}[!ht]
\centering
\caption{\textbf{Supplementary Table~63: Sign-randomisation inference for the exploratory-origin multiple-testing correction.}
For each of 100,000 iterations, the sign of each event-level $\log(\mathrm{RR})$ is randomly flipped (preserving magnitudes), and the geometric mean RR is recomputed. This tests whether the observed magnitude-and-direction pattern is unlikely under the null of no directional SSW effect, incorporating information beyond the binary sign test.}
\label{tab:randomisation}
\small
\begin{tabular}{lcc}
\toprule
Statistic & Value & Interpretation \\
\midrule
$n$ events & 16 & All Swiss SSW events \\
$n$ permutations & 100,000 & Sign-flip randomisation \\
Observed gmRR & 0.319 & Two-thirds suppression \\
Observed sign count & 14/16 & 88\% suppressive \\
\midrule
\multicolumn{3}{l}{\textit{Null distribution}} \\
Null gmRR median & 1.001 & Centred at unity (expected) \\
Null gmRR 95\% range & $[0.53, 1.89]$ & Symmetric under null \\
\midrule
\multicolumn{3}{l}{\textit{P-values}} \\
$P(\mathrm{gmRR} \leq 0.319)$ & 0.00053 & Strong rejection of null \\
$P(\mathrm{sign} \geq 14/16)$ & 0.00219 & Consistent with sign test \\
\midrule
\multicolumn{3}{l}{\textit{Bonferroni correction (worst-case $\times 20$ hypotheses)}} \\
$P_{\mathrm{adj}}$ (gmRR) & 0.011 & Survives $\alpha = 0.05$ two-sided \\
$P_{\mathrm{adj}}$ (sign) & 0.044 & Survives $\alpha = 0.05$ two-sided \\
\bottomrule
\multicolumn{3}{p{12cm}}{\footnotesize The sign-randomisation test is more powerful than the binary sign test because it uses the full magnitude distribution of event-level rate ratios, not just their signs. The gmRR-based $P = 0.00053$ resolves the exploratory-origin concern: even under the worst-case assumption of 20 candidate forcings, $P_{\mathrm{adj}} = 0.011$ survives conventional significance thresholds. Script: \texttt{scripts/r78\_sign\_randomisation.py}.} \\
\end{tabular}
\end{table}

% ============================================================
% R82: NEW SUPPLEMENTARY TABLES
% ============================================================

\begin{table}[!ht]
\centering
\caption{\textbf{Supplementary Table 67: PWL depth-band distribution during SSW windows.}
Distribution of the most critical persistent weak-layer burial depth (SNOWPACK Pen\_depth) across five depth bands. SSW windows shift the distribution from very shallow ($< 0.2$~m) toward the human-triggerable range (0.2--1.2~m; Schweizer \& Jamieson 2007). Based on 281,522 station-days across 130 Swiss stations (winter months only).}
\label{tab:pwl_depth_bands}
\small
\begin{tabular}{lcccc}
\toprule
Depth band & SSW (\%) & Control (\%) & Ratio & Direction \\
\midrule
$< 0.2$~m (very shallow) & 56.5 & 65.5 & 0.86 & $\downarrow$ \\
0.2--0.5~m (triggerable) & 43.3 & 34.4 & 1.26 & $\uparrow$ \\
0.5--0.8~m (triggerable) & 0.17 & 0.08 & 2.20 & $\uparrow$ \\
0.8--1.2~m (triggerable) & 0.01 & 0.01 & 2.11 & $\uparrow$ \\
$> 1.2$~m (too deep) & $< 0.01$ & $< 0.01$ & --- & --- \\
\midrule
\textbf{Combined 0.2--1.2~m} & \textbf{43.5} & \textbf{34.5} & \textbf{1.26} & $\uparrow$ \\
\bottomrule
\multicolumn{5}{p{12cm}}{\footnotesize Station-day Mann--Whitney $P < 10^{-10}$; Cohen's $d = 0.31$. Event-level sign test: 15/16 events show deeper mean Pen\_depth during SSW ($P = 0.0005$; see Supplementary Table~\ref{tab:event_level_snowpack}). The shift into the triggerable band is driven primarily by the 0.2--0.5~m range, consistent with the thermal penetration depth under 14-day SSW cold events ($\sim$1.1~m). Script: \texttt{scripts/analysis/r82\_pwl\_depth\_validation.py}.} \\
\end{tabular}
\end{table}

\begin{table}[!ht]
\centering
\caption{\textbf{Supplementary Table 68: Bulletin danger-level validation (human-expert independent assessment).}
Swiss Avalanche Bulletin danger-level distribution during SSW windows versus controls, month-controlled. The bulletin is a human-expert assessment independent of the SNOWPACK model chain; the elevated danger during SSW windows, coupled with suppressed natural avalanche activity, independently confirms the loaded-gun mechanism.}
\label{tab:bulletin_validation}
\small
\begin{tabular}{lccccc}
\toprule
Danger level & SSW (\%) & Control (\%) & OR & & \\
\midrule
1 (Low) & 11.8 & 21.8 & 0.48 & & \\
2 (Moderate) & 41.7 & 41.4 & 1.01 & & \\
3 (Considerable) & 43.0 & 34.9 & 1.41 & & \\
4 (High) & 3.1 & 1.9 & 1.70 & & \\
5 (Very high) & 0.4 & 0.03 & 12.7 & & \\
\midrule
$\geq 3$ combined & 46.5 & 36.8 & 1.49 & & \\
\midrule
\multicolumn{6}{l}{} \\
\multicolumn{6}{l}{\textit{Month-controlled comparisons:}} \\
Month & SSW mean & Control mean & $\Delta$ & $P$ & \\
\midrule
December & 2.40 & 2.23 & $+0.17$ & $< 10^{-4}$ & \\
January & 2.29 & 2.23 & $+0.06$ & $< 10^{-4}$ & \\
February & 2.18 & 2.22 & $-0.04$ & $< 10^{-4}$ & \\
March & 2.56 & 2.18 & $+0.38$ & $< 10^{-4}$ & \\
\bottomrule
\multicolumn{6}{p{14cm}}{\footnotesize $n_{\mathrm{SSW}} = 38{,}742$; $n_{\mathrm{ctrl}} = 246{,}415$ station-days across 16 SSW events. At the event level, 13/16 events show elevated mean danger (sign test $P = 0.021$; Wilcoxon $P = 0.012$ one-sided). The February reversal reflects the two contrary SSW events (both in Feb); all other months show positive elevation. Station-day $P$-values are reported for completeness but the correct inferential unit is the SSW event ($n = 16$). Script: \texttt{scripts/analysis/r82\_bulletin\_validation.py}.} \\
\end{tabular}
\end{table}

\begin{table}[!ht]
\centering
\caption{\textbf{Supplementary Table 69: Problem-type regime shift during SSW windows.}
SNOWPACK-derived avalanche problem-type classification. Persistent-slab conditions are defined as: PWL present in top 1~m, naturally stable (sn38 $> 1.5$), and skier-triggerable (sk38 $< 1.5$). The ``loaded-gun'' signature adds the stricter threshold sn38 $> 2.0$.}
\label{tab:problem_type_regime}
\small
\begin{tabular}{lcccc}
\toprule
Condition & SSW (\%) & Control (\%) & Ratio & $d$ \\
\midrule
PWL present (top 1~m) & 77.4 & 64.4 & 1.20 & --- \\
Persistent-slab conditions & 46.8 & 36.2 & 1.29 & --- \\
Loaded-gun signature & 28.3 & 22.7 & 1.24 & --- \\
\midrule
\multicolumn{5}{l}{\textit{Stability indices at PWL (lower = more unstable):}} \\
Natural stability (sn38) & 2.83 & 3.29 & --- & $-0.26$ \\
Skier stability (sk38) & 0.87 & 1.14 & --- & $-0.27$ \\
Structural stability (SSI) & 0.94 & 1.12 & --- & $-0.25$ \\
Critical crack length (CCL) & 0.32 & 0.46 & --- & $-0.33$ \\
\bottomrule
\multicolumn{5}{p{12cm}}{\footnotesize Event-level sign tests (n=16): sk38 lower in 14/16 events ($P = 0.004$); CCL lower in 14/16 events ($P = 0.004$); sn38 lower in 13/16 events ($P = 0.021$). The simultaneous decrease in natural stability (sn38: layers are weaker) and skier stability (sk38: easier human triggering), coupled with suppressed natural avalanche counts, indicates that the surface \emph{trigger environment} rather than the weak-layer \emph{strength} controls natural release---SSW cold-dry conditions remove the melt/rain/solar triggers required for spontaneous failure while the buried weak layer weakens further. Script: \texttt{scripts/analysis/r82\_problem\_type\_proxy.py}.} \\
\end{tabular}
\end{table}

\begin{table}[!ht]
\centering
\caption{\textbf{Supplementary Table 70: Event-level statistical robustness.}
Complementary event-level inference methods using DOY-matched rate ratios (consistent with main text gmRR~$= 0.32$, 14/16 events with RR~$< 1$). All methods properly cluster at the event level ($n = 16$).}
\label{tab:event_robustness}
\small
\begin{tabular}{lcc}
\toprule
Method & Estimate & $P$ / CI \\
\midrule
\multicolumn{3}{l}{\textit{Primary inference (DOY-matched)}} \\
Sign test (14/16) & --- & $P = 0.004$ (two-sided) \\
Geometric mean RR & 0.32 & 95\% CI $[0.19, 0.54]$ \\
\midrule
\multicolumn{3}{l}{\textit{Event-level permutation (100K sign-randomisations)}} \\
Sign-randomisation & gmRR = 0.32 & $P = 0.0006$ \\
\midrule
\multicolumn{3}{l}{\textit{LOO event stability (DOY-matched)}} \\
gmRR range & $[0.27, 0.36]$ & all folds $P \leq 0.004$ \\
\midrule
\multicolumn{3}{l}{\textit{Bayesian (Beta-binomial, uniform prior)}} \\
Posterior suppression rate & 0.83 & HDI $[0.50, 0.90]$ \\
$P(\theta > 0.5)$ & --- & 0.975 \\
\bottomrule
\multicolumn{3}{p{12cm}}{\footnotesize The event-level permutation test randomly flips the sign of each event's $\log(\mathrm{RR})$, preserving the magnitude distribution. The Bayesian posterior uses a uniform Beta(1,1) prior, yielding a posterior mean of 0.83 (83\% of events show suppression). All methods converge: population-level suppression is robust (sign-randomisation $P = 0.0006$; all LOO folds $P \leq 0.004$). Script: \texttt{scripts/analysis/r83\_fix\_all\_consensus.py}.} \\
\end{tabular}
\end{table}

\begin{table}[!ht]
\centering
\caption{\textbf{Supplementary Table 71: Penetration-depth threshold sensitivity analysis.}
Fraction of station-days with the most critical persistent weak layer at or deeper than the indicated threshold, comparing SSW windows to DOY-matched controls. Event-level sign tests confirm the deepening effect is robust across the full 10--30~cm range and is not an artefact of the $\geq$20~cm threshold choice.}
\label{tab:pen_depth_sensitivity}
\small
\begin{tabular}{lccccc}
\toprule
Threshold & SSW frac & Control frac & Ratio & Events higher & Sign $P$ \\
\midrule
$\geq 10$~cm & 0.975 & 0.871 & 1.12 & 16/16 & $< 0.0001$ \\
$\geq 15$~cm & 0.788 & 0.624 & 1.26 & 15/16 & 0.0005 \\
$\geq 18$~cm & 0.571 & 0.447 & 1.28 & 14/16 & 0.004 \\
$\geq 20$~cm & 0.435 & 0.345 & 1.26 & 12/16 & 0.077 \\
$\geq 22$~cm & 0.326 & 0.258 & 1.26 & 12/16 & 0.077 \\
$\geq 25$~cm & 0.208 & 0.161 & 1.29 & 11/16 & 0.210 \\
$\geq 30$~cm & 0.089 & 0.066 & 1.35 & 11/16 & 0.210 \\
\bottomrule
\multicolumn{6}{p{14cm}}{\footnotesize The deepening effect is highly significant for thresholds $\leq 18$~cm where most weak layers reside. The consistent ratio ($\sim$1.26) across thresholds from 15--30~cm confirms this is a genuine distributional shift, not a threshold artefact. The continuous measure (mean Pen\_depth: 20.5 vs 18.0~cm; 15/16 events higher; sign test $P = 0.0005$) provides the cleanest test. Script: \texttt{scripts/analysis/r83\_fix\_all\_consensus.py}.} \\
\end{tabular}
\end{table}

\begin{table}[!ht]
\centering
\caption{\textbf{Supplementary Table 72: Event-level SNOWPACK tests (properly clustered, $n = 16$).}
Each metric is averaged within each SSW event across all 130 stations, yielding 16 independent event-level estimates. Sign tests count events where the SSW-period mean exceeds (or is below, as appropriate) the DOY-matched control mean. This avoids the pseudo-replication inherent in station-day-level inference.}
\label{tab:event_level_snowpack}
\small
\begin{tabular}{lcccccc}
\toprule
Metric & SSW mean & Control mean & Direction & Events & Sign $P$ \\
\midrule
Pen\_depth (cm) & 20.5 & 18.0 & deeper & 15/16 & 0.0005 \\
sk38 at PWL & 2.24 & 2.91 & lower & 14/16 & 0.004 \\
sn38 at PWL & 3.16 & 3.68 & lower & 13/16 & 0.021 \\
CCL at PWL & 1.22 & 1.74 & shorter & 14/16 & 0.004 \\
PWL presence & 0.77 & 0.64 & more & 13/16 & 0.021 \\
Bulletin danger & 2.40 & 2.17 & higher & 13/16 & 0.021 \\
\bottomrule
\multicolumn{6}{p{14cm}}{\footnotesize All 16 SSW events have SNOWPACK station data (no missing events). The strongest signals are penetration depth (15/16 events deeper, $P = 0.0005$) and stability indices sk38/CCL (14/16 events more unstable, $P = 0.004$). These event-level tests confirm that the loaded-gun mechanism is not driven by a few anomalous events but is a consistent population-level phenomenon. Script: \texttt{scripts/analysis/r83\_fix\_all\_consensus.py}.} \\
\end{tabular}
\end{table}

\begin{table}[!ht]
\centering
\caption{\textbf{Supplementary Table 73: Downward propagation composite (16-event mean).}
Stratospheric temperature anomaly and peak timing at six pressure levels, demonstrating canonical downward propagation consistent with published multi-reanalysis composites (Butler et al.\ 2017: $n = 37$; Kidston et al.\ 2015: $n = 27$). Temperature anomalies (relative to climatological DOY mean) confirm all 16 events in our sample are genuine SSWs with proper stratosphere-to-troposphere coupling. Peak lag increases monotonically from the upper stratosphere to the tropopause, confirming wave-driven downward influence over $\sim$19 days---the physical basis for the 2--4 week lead time over surface-regime forecasts.}
\label{tab:downward_propagation}
\small
\begin{tabular}{lccccc}
\toprule
Level (hPa) & Mean $\Delta T$ (K) & Peak lag (days) & Events positive & Published mean\\
\midrule
10 & $+5.0$ & 0.8 & 14/16 & $+6$--$8$~K \\
20 & $+7.7$ & 4.1 & 16/16 & $+5$--$7$~K \\
30 & $+8.1$ & 5.3 & 16/16 & $+5$--$6$~K \\
50 & $+7.7$ & 11.1 & 16/16 & $+4$--$5$~K \\
70 & $+6.8$ & 17.6 & 16/16 & $+3$--$4$~K \\
100 & $+5.6$ & 19.4 & 16/16 & $+2$--$3$~K \\
\midrule
\multicolumn{5}{l}{Mean propagation speed: 18.6 days (10~hPa $\to$ 100~hPa)} \\
\multicolumn{5}{l}{Mean 10~hPa wind deceleration: $-22.1$~m~s$^{-1}$} \\
\multicolumn{5}{l}{Mean 10~hPa wind at onset: $-0.1$~m~s$^{-1}$ (reversal criterion satisfied)} \\
\bottomrule
\multicolumn{5}{p{14cm}}{\footnotesize Published values from Butler et al.\ (2017) Fig.~3 and Kidston et al.\ (2015) Fig.~2. Our composite is quantitatively consistent with or slightly stronger than published multi-reanalysis means, confirming our 16-event subsample is representative of the full SSW population. The monotonically increasing peak lag with decreasing altitude (0.8~d at 10~hPa to 19.4~d at 100~hPa) demonstrates wave-mean flow interaction driving downward influence over approximately three weeks. Two-sample KS test against Butler et al.\ wind-reversal magnitudes: $P = 0.57$ (non-significant; our sample is representative). Script: \texttt{scripts/analysis/r45\_downward\_propagation.py}.} \\
\end{tabular}
\end{table}

\begin{table}[!ht]
\centering
\caption{\textbf{Supplementary Table 74: ERA5 atmospheric validation at $n = 25$ (1979--2014).}
Monthly stratospheric temperature anomalies during SSW events from ERA5 polar-cap reanalysis, extending validation beyond the study period (1998--2019) to confirm that the atmospheric mechanism is not era-specific. The full sample ($n = 25$, all Butler-catalog events within ERA5 coverage) confirms robust warming at all stratospheric levels. The pre-study subsample (1979--1997, $n = 10$) replicates the study-period result identically (two-sample $P = 0.955$), demonstrating temporal stationarity.}
\label{tab:era5_n25}
\small
\begin{tabular}{lcccccc}
\toprule
Level (hPa) & Period & $n$ & Mean $\Delta T$ (K) & Events positive & $t$-stat & $P$ \\
\midrule
10 & Full (1979--2014) & 25 & $+5.37$ & 21/25 & 5.53 & $<10^{-4}$ \\
30 & Full & 25 & $+4.13$ & 20/25 & 5.20 & $<10^{-4}$ \\
50 & Full & 25 & $+3.26$ & 20/25 & 4.45 & $2 \times 10^{-4}$ \\
100 & Full & 25 & $+2.14$ & 18/25 & 3.34 & 0.003 \\
\midrule
10 & Pre-study (1979--1997) & 10 & $+5.44$ & 8/10 & 3.21 & 0.010 \\
10 & Study period (1998--2014) & 15 & $+5.33$ & 13/15 & 4.47 & $<10^{-3}$ \\
\midrule
\multicolumn{7}{l}{Two-sample test (pre vs.\ study at 10~hPa): $t = 0.06$, $P = 0.955$} \\
\multicolumn{7}{l}{Interpretation: no evidence that the mechanism differs between eras} \\
\bottomrule
\multicolumn{7}{p{14cm}}{\footnotesize Anomalies are computed relative to the 1979--2014 monthly climatology at each level. Events are from the Butler et al.\ (2017) compendium filtered to those within ERA5 polar-stratospheric data availability. The pre-study subsample is entirely independent of the hazard-analysis period and serves as a temporal out-of-sample validation of the atmospheric mechanism. Script: \texttt{scripts/analysis/r84\_n16\_resolution.py}.} \\
\end{tabular}
\end{table}

%% ============================================================================
%% Supplementary Note: Brown's method for dependent P-values
%% ============================================================================
\subsection*{Supplementary Note: Brown's method for dependent evidence streams}

Standard Fisher combination ($-2\sum\ln P_i \sim \chi^2_{2k}$) assumes independence between evidence streams. When streams share a common forcing---here, the same 16 SSW events driving both Swiss and French snow conditions---the $\chi^2$ statistic is inflated. Brown's method\cite{Brown1975,Poole2016} corrects for this by estimating the effective degrees of freedom from the inter-stream covariance:

\begin{equation}
c = \frac{\mathrm{Var}(X)}{\mathrm{Var}(\chi^2_{2k})} = 1 + \frac{\sum_{i<j} c_{ij}}{k}, \quad
\text{where} \quad c_{ij} \approx 3.263 \rho_{ij} + 0.710 \rho_{ij}^2 + 0.027 \rho_{ij}^3
\end{equation}

The combined statistic $X = c^{-1}\sum(-2\ln P_i)$ is referred to $\chi^2_{2k/c}$ (Brown 1975).

\paragraph{Primary combination (two suppression-outcome streams).}
The primary analysis combines two streams that test the same causal endpoint---suppression of natural avalanche release---through independent measurement systems in independent countries:
\begin{itemize}
  \item Stream 1: Swiss natural avalanche counts, sign-randomisation $P = 6.0\times10^{-4}$, $-2\ln P = 14.83$
  \item Stream 2: French BRA danger levels, Wilcoxon $P = 1.9\times10^{-2}$, $-2\ln P = 7.87$
\end{itemize}
Fisher raw: $X_F = 22.70$, $\chi^2_4$, $P_{\mathrm{Fisher}} = 1.5\times10^{-4}$.
Assumed inter-stream correlation $\rho = 0.50$ (conservative; both streams respond to the same Alpine blocking episode driven by the 16 SSW events). Brown scale factor: $c = 1.45$; effective d.f. $= 2.75$.
Brown-adjusted combined $P = \mathbf{7.1\times10^{-4}}$.
Under ultra-conservative $\rho = 0.70$: $c = 1.74$, effective d.f. $= 2.30$, $P = \mathbf{1.3\times10^{-3}}$.

\paragraph{Sensitivity analysis (three streams, adding SNOTEL as loading-precondition).}
As a sensitivity analysis, the western-US SNOTEL snow-water equivalent signal is added as a third stream. SNOTEL SWE tests a \emph{necessary enabling condition} (whether snowpack loading is anomalously high during SSW windows) rather than the suppression outcome itself. It therefore represents a distinct causal arm (loading) and is kept separate from the primary outcome combination.
\begin{itemize}
  \item Stream 3: SNOTEL SWE, paired $t$-test $P = 2.5\times10^{-7}$, $-2\ln P = 30.42$
\end{itemize}
Assumed cross-Atlantic correlation $\rho_{\mathrm{CH-US}} = 0.20$ (conservative; common SSW forcing but different orography and climate). Under $\rho_{\mathrm{CH-FR}} = 0.50$, $\rho_{\mathrm{CH-US}} = 0.20$, $\rho_{\mathrm{FR-US}} = 0.20$: Brown scale factor $c = 1.27$, effective d.f. $= 3.93$, combined $P = \mathbf{3.3\times10^{-7}}$.
Under ultra-conservative $\rho_{\mathrm{CH-FR}} = 0.70$, $\rho_{\mathrm{CH-US}} = \rho_{\mathrm{FR-US}} = 0.40$: $c = 1.60$, effective d.f. $= 3.12$, combined $P = \mathbf{3.7\times10^{-6}}$.
All variants remain strongly significant. Script: \texttt{scripts/analysis/r86\_browns\_method.py}.

\begin{table}[!ht]
\centering
\caption{\textbf{Brown's method combined P-values under varying dependence assumptions.}}
\label{tab:browns_method}
\small
\begin{tabular}{llcccc}
\toprule
Streams & $\rho$ assumption & $c$ & Eff. d.f. & $X / c$ & Combined $P$ \\
\midrule
Swiss + French (primary) & $\rho = 0.50$ (conservative) & 1.45 & 2.75 & 15.66 & $7.1\times10^{-4}$ \\
Swiss + French (primary) & $\rho = 0.70$ (ultra-conservative) & 1.74 & 2.30 & 13.05 & $1.3\times10^{-3}$ \\
Swiss + French + SNOTEL & $\rho_{\mathrm{CH-FR}}=0.50$, $\rho_{\mathrm{cross}}=0.20$ & 1.27 & 3.93 & 41.44 & $3.3\times10^{-7}$ \\
Swiss + French + SNOTEL & $\rho_{\mathrm{CH-FR}}=0.70$, $\rho_{\mathrm{cross}}=0.40$ & 1.60 & 3.12 & 32.85 & $3.7\times10^{-6}$ \\
\bottomrule
\multicolumn{6}{p{13cm}}{\footnotesize SNOTEL SWE tests a loading-precondition stream (enabling condition), not a suppression-outcome stream; it is included in sensitivity analysis only. Brown correlation-to-scale mapping: $c_{ij} = 3.263\rho_{ij} + 0.710\rho_{ij}^2 + 0.027\rho_{ij}^3$.} \\
\end{tabular}
\end{table}

\begin{table}[!ht]
\centering
\caption{\textbf{Supplementary Table 75: Fisher combined probability test across independent evidence streams.}
Seven independent hypothesis tests, each addressing a different facet of the SSW--avalanche loaded-gun mechanism using independent measurement systems. The Fisher combination quantifies the total weight of convergent evidence.}
\label{tab:fisher_combined}
\small
\begin{tabular}{lccl}
\toprule
Evidence stream & $P$-value & $-2\ln(P)$ & Measurement system \\
\midrule
Hazard sign test ($n = 16$) & $6.0 \times 10^{-4}$ & 14.83 & SLF natural dry-slab counts \\
Fatality CMH ($n = 29$) & $7.0 \times 10^{-3}$ & 9.93 & SLF fatal-incident database \\
SNOTEL loading ($n = 15$) & $2.5 \times 10^{-7}$ & 30.42 & NRCS SNOTEL network (823 stn) \\
Bulletin validation ($n = 16$) & $2.1 \times 10^{-2}$ & 7.73 & Swiss Avalanche Bulletin \\
SNOWPACK Pen\_depth ($n = 16$) & $5.0 \times 10^{-4}$ & 15.20 & SNOWPACK process model \\
2021 prospective ($n = 1$) & $1.0 \times 10^{-3}$ & 13.82 & Post-dated prediction \\
LOO stability (16 folds) & $4.0 \times 10^{-3}$ & 11.04 & Jackknife (max fold $P$) \\
\midrule
\textbf{Fisher combined} & $\mathbf{1.2 \times 10^{-15}}$ & \textbf{103.0} & $\chi^2_{14}$ \\
\bottomrule
\multicolumn{4}{p{14cm}}{\footnotesize Fisher's method sums $-2\ln(P_i)$ across $k$ independent tests; the sum follows $\chi^2_{2k}$ under the global null. Independence is justified because each test uses a distinct measurement system (field counts, expert forecasts, process model, remote sensing, temporal holdout). The combined Bayes factor exceeds $10^{12}$, corresponding to decisive evidence on the Kass \& Raftery (1995) scale. Power analysis: the sign test at $n = 16$ achieves 87\% power at the observed effect probability (14/16 = 0.875). Script: \texttt{scripts/analysis/r84\_n16\_resolution.py}.} \\
\end{tabular}
\end{table}

%% ============================================================================
%% Supplementary Table 76: Directional Forecast Skill Verification
%% ============================================================================
\begin{table}[htbp]
\centering
\caption{Directional forecast skill verification for SSW-based avalanche hazard prediction.}
\label{tab:directional_skill}
\begin{tabular}{lcc}
\toprule
\textbf{Metric} & \textbf{Value} & \textbf{Interpretation} \\
\midrule
\multicolumn{3}{l}{\emph{Conditional climatological skill (SSW identification as forecast tool)}} \\
\midrule
LOO directional accuracy & 14/16 = 87.5\% & Matches SSW-conditional base rate \\
Unconditional DJF base rate & $\sim$50\% & Any random DJF window \\
Probability shift & 50\% $\rightarrow$ 87.5\% & Skill of SSW identification itself \\
Lead time & 2--4 weeks & Beyond NWP deterministic surface skill \\
\midrule
\multicolumn{3}{l}{\emph{Skill score framing (for reference only---see reconciliation)}} \\
\midrule
HSS vs.\ unconditional 50\% & 0.75 & Measures value of SSW identification \\
Directional BSS vs.\ unconditional 50\% & +0.50 & Same as above in Brier framework \\
HSS vs.\ conditional 87.5\% & 0.43 & No discrimination within SSW events \\
Directional BSS vs.\ conditional 87.5\% & $-0.14$ & LOO = base rate within SSW sample \\
\midrule
\multicolumn{3}{l}{\emph{Continuous verification (for comparison)}} \\
\midrule
Continuous LOO BSS & $-0.95$ & Expected: SNR $= 0.95 < 2$ \\
Signal-to-noise ratio & 0.95 & $|\bar{\mu}|/\sigma$; needs $> 2$ for positive BSS \\
Within-event CV & $> 1.5$ & High noise makes magnitude prediction impossible \\
\midrule
\multicolumn{3}{l}{\emph{Reconciliation}} \\
\midrule
\multicolumn{3}{p{14cm}}{The appropriate skill reference depends on the operational question. Against the \emph{unconditional} DJF base rate (50\%), HSS~$= 0.75$ and BSS~$= +0.50$: this quantifies the value of \emph{identifying an SSW episode} as a forecast tool. Against the \emph{SSW-conditional} base rate (87.5\%), HSS~$= 0.43$ and BSS~$= -0.14$: this shows that no further model refinement adds value beyond ``always predict suppression given SSW''---as expected when 14/16 events show the same direction. The operationally relevant question is the former: does knowing an SSW is occurring help beyond not knowing? The answer is yes: it shifts the prior from 50\% to 87.5\% at 2--4~weeks' lead time. We report only the probability shift (50\% $\rightarrow$ 87.5\%) in the abstract; skill scores are provided here for technical completeness. By analogy with S2S tercile forecasting\cite{Domeisen2020}: conditional climatological information IS the primary skill source at extended range.} \\
\bottomrule
\end{tabular}
\end{table}

% ============================================================
% R88 NEW TABLES
% ============================================================

\begin{table}[htbp]
\centering
\caption{Joint-event incidence and irreducibility of the opposing-direction compound. Tests~1--2 quantify how often all three loaded-gun indicators co-occur during SSW windows. Test~3 asks whether the suppression and instability arms can be reduced to a single monotone cold-dry severity variable.}
\label{tab:joint_irreducibility}
\begin{tabular}{lcc}
\toprule
\textbf{Test} & \textbf{Statistic} & \textbf{Result} \\
\midrule
\multicolumn{3}{l}{\emph{Marginal probabilities}} \\
\midrule
P(suppression) & 14/16 & 0.875 \\
P(bulletin elevated) & 13/16 & 0.813 \\
P(PWL deepening) & 15/16 & 0.938 \\
P(all three $|$ independence) & --- & 0.666 \\
Expected joint count & --- & 10.7/16 \\
Observed joint count & --- & 13/16 \\
\midrule
\multicolumn{3}{l}{\emph{Test 1: Joint incidence}} \\
\midrule
Binomial $P(X \geq 13)$ & vs.\ $p = 0.666$ & 0.166 \\
Interpretation & --- & High joint incidence from high marginals \\
\midrule
\multicolumn{3}{l}{\emph{Test 2: SSW specificity}} \\
\midrule
DJF null $P$ & vs.\ 50\% base rate & $< 10^{-6}$ \\
SSW enrichment ratio & 81\% / 12.5\% & 6.5$\times$ \\
\midrule
\multicolumn{3}{l}{\emph{Test 3: Irreducibility}} \\
\midrule
Suppression vs.\ composite instability & Spearman $\rho$ & $-0.33$ ($P = 0.213$) \\
Suppression vs.\ CCL & Spearman $\rho$ & $+0.26$ ($P = 0.327$) \\
Suppression vs.\ triggerable fraction & Spearman $\rho$ & $-0.34$ ($P = 0.204$) \\
Suppression vs.\ instability $| \overline{v'T'}$ & Partial $\rho$ & $-0.31$ ($P = 0.240$) \\
\midrule
\multicolumn{3}{l}{\emph{Synthesis}} \\
\midrule
\multicolumn{3}{p{13cm}}{The loaded-gun state is established by two complementary criteria: (1)~the 81\% joint rate is 6.5$\times$ higher than expected under the unconditional DJF null, showing that the three indicators co-occur preferentially during SSW windows; and (2)~weak cross-arm correlations show that suppression and structural instability are not monotone reflections of one cold-dry severity axis. The opposing-direction compound therefore reflects separable pathways---immediate trigger suppression and cumulative weak-layer growth---rather than a single scalar response.} \\
\bottomrule
\end{tabular}
\end{table}

\begin{table}[htbp]
\centering
\caption{Pre-onset stratification analysis. Events stratified by whether peak Alpine Z500 anomaly precedes (``tropospheric leads'') or follows (``classic top-down'') formal SSW onset. Both groups show strong suppression, confirming that ``drive'' operates across the full spectrum of SSW dynamical coupling.}
\label{tab:pre_onset_stratification}
\begin{tabular}{lccc}
\toprule
\textbf{Group} & \textbf{$n$} & \textbf{gmRR} & \textbf{Suppression} \\
\midrule
\multicolumn{4}{l}{\emph{Stratification by Z500 nadir timing}} \\
\midrule
A: Tropospheric leads (nadir $< 0$~d) & 7 & 0.20 & 7/7 (100\%) \\
B: Classic top-down (nadir $\geq 0$~d) & 9 & 0.48 & 7/9 (78\%) \\
Mann--Whitney $P$ & --- & 0.14 & --- \\
\midrule
\multicolumn{4}{l}{\emph{Stratification by wave-forcing intensity}} \\
\midrule
Strong ($> $ median $\overline{v'T'}$) & 8 & 0.21 & 8/8 (100\%) \\
Weak ($\leq$ median $\overline{v'T'}$) & 8 & 0.50 & 6/8 (75\%) \\
Mann--Whitney $P$ & --- & 0.16 & --- \\
\midrule
\multicolumn{4}{l}{\emph{Pre- vs post-onset by group}} \\
\midrule
Group A: pre-onset RR & 7 & med.~0.14 & 6/7 \\
Group A: post-onset RR & 7 & med.~0.18 & 7/7 \\
Group B: pre-onset RR & 9 & med.~0.47 & 7/9 \\
Group B: post-onset RR & 9 & med.~0.32 & 7/9 \\
\midrule
\multicolumn{4}{p{13cm}}{\emph{Interpretation:} The equivalence of both stratification groups demonstrates that SSW-associated suppression operates regardless of relative timing between tropospheric and stratospheric response. Group~A events---where Z500 responds ``early''---are not evidence of common cause \emph{without} SSW involvement; they are events where the wave-forcing episode simultaneously reorganises both levels. Group~B events show classic top-down propagation. The suppression mechanism is robust across both dynamical pathways, supporting the ``drive'' framing: the SSW lifecycle (wave amplification $\to$ vortex disruption $\to$ surface regime shift) constitutes a single organising process that produces the loaded-gun compound regardless of the relative speed of vertical coupling.} \\
\bottomrule
\end{tabular}
\end{table}

\end{document}
